\documentclass[11pt,reqno]{amsart}

\usepackage[T1]{fontenc}
\usepackage[utf8]{inputenc}
\usepackage{lmodern}
\usepackage[margin=1in]{geometry}
\usepackage{amsmath,amssymb,mathtools,mathrsfs}
\usepackage{booktabs,longtable,array,tabularx}
\usepackage{enumitem}
\usepackage{microtype}
\usepackage{xurl}
\usepackage[
  hidelinks,
  pdftitle={Hydra Maps, Numens, and (p,q)-Adic Analysis: A Critical Expository Edition of the Mathematical Work of Maxwell C. Siegel},
  pdfauthor={Maxwell C. Siegel (source works)},
  pdfsubject={Canonical critical-expository edition of the mathematical corpus of Maxwell C. Siegel}
]{hyperref}
\usepackage[nameinlink,noabbrev]{cleveref}

\newtheorem{theorem}{Theorem}[section]
\newtheorem{lemma}[theorem]{Lemma}
\newtheorem{proposition}[theorem]{Proposition}
\newtheorem{corollary}[theorem]{Corollary}
\newtheorem{conjecture}[theorem]{Conjecture}
\theoremstyle{definition}
\newtheorem{definition}[theorem]{Definition}
\newtheorem{construction}[theorem]{Construction}
\newtheorem{sourceissue}[theorem]{Source issue}
\newtheorem{status}[theorem]{Status record}
\theoremstyle{remark}
\newtheorem{remark}[theorem]{Remark}
\newtheorem{editorial}[theorem]{Editorial note}

\newcommand{\N}{\mathbb N}
\newcommand{\Z}{\mathbb Z}
\newcommand{\Q}{\mathbb Q}
\newcommand{\R}{\mathbb R}
\newcommand{\C}{\mathbb C}
\newcommand{\F}{\mathbb F}
\newcommand{\OK}{\mathcal O_K}
\newcommand{\End}{\operatorname{End}}
\newcommand{\Aff}{\operatorname{Aff}}
\newcommand{\Tr}{\operatorname{Tr}}
\newcommand{\supp}{\operatorname{supp}}
\newcommand{\Span}{\operatorname{span}}
\newcommand{\Dig}{\operatorname{Dig}_p}
\newcommand{\sh}{\vartheta_p}
\newcommand{\Num}{X_H}
\newcommand{\mult}{M_H}
\newcommand{\ind}[1]{\mathbf 1_{#1}}
\newcommand{\vpp}{\vphantom{\Large A}}

\title[Hydra maps, numens, and $(p,q)$-adic analysis]
{Hydra Maps, Numens, and $(p,q)$-Adic Analysis:\\
A Critical Expository Edition of the Mathematical Work of Maxwell C. Siegel}
\author{Mathematical works by Maxwell C. Siegel}
\thanks{The critical apparatus, typed reconstructions, counterexamples, and
edition-supplied proofs are components of this edition and are not attributed
to Siegel as source text.}
\date{ArXiv corpus through 12 February 2026; source audit and edition dated
24 August 2026}

\begin{document}

\begin{abstract}
This critical expository edition reconstructs the mathematical programme of
Maxwell C. Siegel from his author-controlled and arXiv corpus.  Its organizing
object is a residue-class piecewise-affine arithmetic dynamical system, called
a Hydra map, together with the digit-recursive function called its numen.
The defining recursion is derived from affine word composition, its extension
from non-negative integers to $p$-adic integers is separated from the
place-dependent convergence hypotheses, and the Correspondence Principle is
stated with all of the hypotheses used in the latest source formulation.
The analytic developments are then placed around that core: singularity
conservation and dreamcatchers; Dirichlet series, Mellin--Perron inversion and
the contour formulation of weak Collatz; $(p,q)$-adic Fourier analysis and
Wiener--Tauberian cyclicity; rising-continuous functions and variable
topologies of convergence; and the later algebra of products of F-series.

This is an edition of Siegel's corpus, not a new joint research article.
Siegel retains authorship of the source mathematics and of his terms
\emph{Hydra map}, \emph{numen}, \emph{dreamcatcher}, and
\emph{rising-continuity}.  Source forms are preserved; any running notation
is placed beside them by an explicit crosswalk.  Published, current-preprint,
superseded, and withdrawn witnesses remain distinct, and source defects are
recorded before a corrected statement is proved.
\end{abstract}

\maketitle
\tableofcontents

\section{Editorial method, attribution, and mathematical status}

\subsection{What this edition is}

The primary sources are the mathematical works written by Maxwell Charles
Siegel and distributed through arXiv or through his author-controlled site.

The purpose of the present text is narrower and more exact.  It fixes a
canonical witness for each work, states how older versions relate to it,
places any running notation beside the typed source notation through an
explicit crosswalk, and arranges the results by mathematical dependence rather
than by upload date.  The result is a modern recasting of the corpus: neither
a facsimile anthology nor a survey that replaces formulas by descriptions.

\subsection{Four statement statuses}

The edition uses the following four dispositions.  A source-issue or status
environment states its disposition locally; Chapter~10 records the source
identity and disposition of each theorem-like item admitted to its result
table.  A theorem-style heading by itself neither assigns authorship nor
certifies that a source proof has been completed.

\begin{enumerate}[label=\textup{(\roman*)}]
\item A \emph{reconstructed theorem} is a statement in Siegel's corpus whose
      source hypotheses and conclusion are preserved and placed beside any
      running notation through an explicit typed crosswalk, and whose proof is
      reproduced or reconstructed at the level needed by the subsequent
      chain.
\item A \emph{source theorem} is stated as a theorem of the named source.  If
      this edition gives only its proof mechanism, that limitation is said
      explicitly; the label does not mean that this edition supplies an
      independent proof of every auxiliary assertion in the source.
\item A \emph{source issue} records a typing error, a changed convention, an
      omitted hypothesis, a withdrawn paper, or a point at which the source
      itself supplies only a sketch.  The issue is not repaired by pretending
      that the altered statement was what the source always said.
\item An \emph{editorial proposition} is a formal consequence proved in this
      edition in order to join two source formulations.  It is marked as
      editorial and is not attributed to Siegel as an original theorem.
\end{enumerate}

Extraction is not verification.  Conversely, a typographical defect does not
by itself refute the intended theorem.  The edition therefore distinguishes
the exact source string, the mathematically typed reconstruction, and the
status of the proof.

\subsection{Notation policy}

The source convention $x\overset{a}{\equiv}y$ is rendered as
$x\equiv y\pmod a$.  The source symbol $[x]_a$ is retained for the standard
representative when it is part of a digit formula.  Iverson brackets are
written $\ind{P}$.  Composition is read from right to left.  Thus, for a word
$\mathbf j=(j_0,\ldots,j_{m-1})$, the ordered branch composition used below
is
\[
 H_{\mathbf j}=H_{j_0}\circ H_{j_1}\circ\cdots\circ H_{j_{m-1}}.
\]
The source sometimes indexes a word from $1$ and sometimes from $0$; the
edition uses $0$-based digit indices and states the convention whenever an
order-sensitive identity is used.

The name \emph{$(p,q)$-adic analysis} is retained as Siegel's name for
analysis of functions whose domain carries a $p$-adic topology and whose
values lie in a field with a distinct $q$-adic (or, more generally,
place-dependent) topology.  This is not the standard adjective ``adic''
applied to a single topological ring; the domain and codomain places must
always be displayed.

The name \emph{rising-continuity} is retained for the source condition
\[
 f(z)=\lim_{n\to\infty} f([z]_{p^n}),
\]
where the limit is taken in the specified codomain.  In standard language it
is continuity along the canonical truncation sequence, not automatically
topological continuity with respect to every sequence converging to $z$.
The distinction is used in \cref{sec:rising-fseries}.

\subsection{Version precedence}

The latest non-withdrawn arXiv version controls a work unless an
author-controlled later witness explicitly corrects it.  A later technical
manual controls a definition when it declares that earlier definitions are
inconsistent and fixes a forward convention.  Under this rule the 2026
Hydra/Numen manual controls the basic definitions of Hydra maps, numens, and
the general Correspondence Principle; the 2022 dissertation remains the
largest source for the analytic development.  A withdrawn record remains in
the historical corpus but cannot control a theorem contradicted by its own
withdrawal notice.

\begin{sourceissue}[Canonical does not mean infallible]
The word \emph{canonical} refers to source identity, version precedence, and
the stable presentation adopted by this edition.  It does not convert an
author claim into a proved theorem.  In particular, the 2026 paper labels its
Correspondence Principle proof a sketch; this edition preserves that fact and
separates the fully checked affine-word calculation from the remaining
dynamical implications.
\end{sourceissue}

\subsection{The exactness standard}

Each comparison is attached to an explicit object.  For example, the chain
used for a characteristic distribution is the following typed sequence, not
the sentence ``Fourier analysis detects periodic points.''  Starting from a
centered admissible pre-Hydra
\[
 \mathcal H=(K,\Lambda,(H_j)_{j\in\mathcal A_p})
 \in\{(K,\Lambda,\mathbf H):H_j\in\Aff_{\F}(K)\},
\]
affine word multiplication gives
\[
 (M_H,X_H)\in
 \operatorname{Map}(\Sigma_p^*,\operatorname{GL}_{\F}(K))
 \times\operatorname{Map}(\Sigma_p^*,K).
\]
Centeredness proves that the second component descends along
$\Dig:\Sigma_p^*\twoheadrightarrow\N_0$ to
$X_H:\N_0\to K$.  At a place $\ell$ let
\[
 D_{H,\ell}=\{z\in\Z_p^{\mathrm{dig}}:
             (X_H([z]_{p^N}))_{N\geq0}\text{ converges in }K_\ell\}.
\]
The truncation limit is a map
\[
 X_{H,\ell}:D_{H,\ell}\longrightarrow K_\ell.
\]
These measurability assertions are automatic.  Each truncation map
$S_N(z)=X_H([z]_{p^N})$ is locally constant.  Writing
$d_\ell(x,y)=|x-y|_\ell$, completeness of $K_\ell$ gives
\[
 D_{H,\ell}
 =\bigcap_{r\geq1}\bigcup_{N_0\geq0}
   \bigcap_{m,n\geq N_0}
   \{z:d_\ell(S_m(z),S_n(z))<1/r\},
\]
so $D_{H,\ell}$ is Borel.  On that Borel subspace,
$X_{H,\ell}=\lim_NS_N$ is Borel as the pointwise limit of metric-space-valued
Borel maps.  Thus only conullness remains a hypothesis.  If
$m_p(D_{H,\ell})=1$, then the pushforward is the Borel probability
\[
 \nu_{H,\ell}=(X_{H,\ell})_*
   (m_p|_{D_{H,\ell}})\in\operatorname{Prob}(K_\ell).
\]
Its Fourier--Stieltjes transform is the bounded continuous function
\[
 \widehat\nu_{H,\ell}:\widehat{K_\ell}\longrightarrow\C,
 \qquad
 \widehat\nu_{H,\ell}(\gamma)
 =\int_{K_\ell}\gamma(-x)\,d\nu_{H,\ell}(x).
\]
If a cyclic-translate space is subsequently considered, translations are
indexed by the character group, not by $K_\ell$: for
$\eta\in\widehat{K_\ell}$,
$(\tau_\eta f)(\gamma)=f(\gamma+\eta)$, and the asserted closure must name
its ambient function space and norm.  No self-duality of $K_\ell$ is being
silently chosen here.  Likewise, the numen is not merely ``a function
encoding the dynamics'': it is the translation term of an ordered affine
branch composition and satisfies a digit recursion with a specified initial
condition.  This level of explicitness is the governing editorial standard
throughout the edition.

\section{The canonical Siegel corpus}
\label{sec:corpus}

\subsection{Logical works and version witnesses}

The primary research corpus occupies ten arXiv lineages as of 23 August
2026.  A lineage is a bibliographic container, not necessarily a homogeneous
sequence of minor revisions.  In particular, 2007.15936v1,
2111.07883v1, and 2208.11082v1 differ substantially from the later versions
filed under the same identifiers and are retained in the apparatus as
distinct substantive witnesses.  A withdrawal stub is a bibliographic
event, not a replacement source for the withdrawn mathematics.  The table
records the controlling witness used by this edition.

\begingroup\small
\begin{longtable}{@{}p{0.13\textwidth}p{0.44\textwidth}p{0.13\textwidth}p{0.22\textwidth}@{}}
\toprule
arXiv & Work & Witness & Editorial status\\
\midrule
\endfirsthead
\toprule
arXiv & Work & Witness & Editorial status\\
\midrule
\endhead
1909.09733 & \emph{Conservation of Singularities in Functional Equations Associated to Collatz-Type Dynamical Systems; or, Dreamcatchers for Hydra Maps} & v3 & Current historical source for dreamcatchers and singularity conservation.\\
2002.12762 & \emph{Fourier Analysis of the Parity-Vector Parameterization of the Generalized Collatz $px+1$ Maps} & v1 & Withdrawn in v2 as obsolete; retained as a superseded stage, not controlling.\\
2007.15936 & \emph{The Collatz Conjecture \& Non-Archimedean Spectral Theory: Part I} & v3 & Published core for the shortened $qx+1$ numen and Correspondence Principle.\\
2111.07882 & \emph{Functional Equations Associated to Collatz-Type Maps on Integer Rings of Algebraic Number Fields} & v1 & Withdrawn in v2 for an error in the section on quotient rings; historical only.\\
2111.07883 & \emph{Non-Archimedean Spectral Theory, Part I.5: How to Write the (Weak) Collatz Conjecture as a Contour Integral} & v2 & Current source for the Dirichlet/Mellin/contour formulation.\\
2208.11082 & \emph{Non-Archimedean Spectral Theory, Part II: $(p,q)$-Adic Fourier Analysis and Wiener's Tauberian Theorem} & v2 & Current and published analytic source.\\
2307.00436 & \emph{Infinite Series Whose Topology of Convergence Varies From Point to Point} & v1 & Published source for frames, rising-continuity, and point-dependent convergence.\\
2412.02902 & \emph{$(p,q)$-Adic Analysis and the Collatz Conjecture} & v1 & ArXiv witness of the 2022 dissertation; principal synthetic source.\\
2507.13358 & \emph{Algebras of $p$-Adic Distributions Induced by Pointwise Products of F-Series} & v1 & Current extension to products, frames, distributions, and breakdown varieties.\\
2601.17030 & \emph{The Hydra Map and Numen Formalisms for Collatz-Type Problems} & v3 & Latest terminology and foundational manual; v3 corrects Proposition 7.6.\\
\bottomrule
\end{longtable}
\endgroup

The author site also supplies mathematically substantive version witnesses.
Its numbering must not be confused with the numbering of the papers:

\begingroup\small
\begin{longtable}{@{}p{0.14\textwidth}p{0.43\textwidth}p{0.35\textwidth}@{}}
\toprule
Witness & Author-site title & Canonical relation\\
\midrule
\endfirsthead
\toprule
Witness & Author-site title & Canonical relation\\
\midrule
\endhead
Web installment I & \emph{The Numen of a Hydra Map} & Expository/version
witness to paper Part I, 2007.15936v3; it does not control a conflicting
formula.\\
Web installment II & \emph{The Correspondence Principle} & Continuation of
Web I and likewise a witness to paper Part I.  It is \emph{not} paper Part
II.  Its Dirichlet--Perron coda is a precursor to 2111.07883v2.\\
Web installment III & \emph{$(p,q)$-adic Fourier Analysis \& Tauberian
Spectral Theory} & Evolving exposition to be collated against paper Part II,
2208.11082v2.  A later webpage modification date alone gives no theorem
precedence.\\
Web installment IV & \emph{A Fourier Analysis of the Numen} & Site-only
application/exposition, collated chiefly against the one-dimensional
dissertation material.\\
\bottomrule
\end{longtable}
\endgroup

Web installments I and II were recovered as current WordPress snapshots
during the audit; the public API does not expose their revision histories.
The site also contains a March 2025 talk note on arithmetic properties of
F-series and companion expository material.  Those witnesses are part of the
author-controlled corpus but do not supersede a later arXiv theorem merely
because their prose is newer.

The authority order used when witnesses conflict is
\[
\begin{gathered}
 \text{2601 definitions}\longrightarrow
 \text{2307 F-series/frame definitions}\\
 \longrightarrow\text{latest specialized paper}\longrightarrow
 \text{dissertation}\\
 \longrightarrow\text{web exposition}\longrightarrow
 \text{historical or withdrawn witness}.
\end{gathered}
\]
This order is claim-specific: it does not erase unique material in an older
witness, and it does not turn a later unproved claim into a theorem.

\subsection{Mathematical dependency order}

The upload chronology obscures the mathematical order.  The canonical
reading order used in this edition is:
\[
\begin{gathered}
\text{residue-class affine dynamics}\longrightarrow
\text{word composition}\longrightarrow\text{numen recursion}\\
\longrightarrow\text{$p$-adic extension}\longrightarrow
\text{Correspondence Principle},\\[3pt]
\text{set-series/permutation operators}\longrightarrow
\text{dreamcatchers}\longrightarrow
\text{singularity conservation},\\[3pt]
\text{Haar integration and Fourier series}\longrightarrow
\text{Wiener--Tauberian cyclicity}\longrightarrow
\text{spectral reformulation},\\[3pt]
\text{rising-continuity and frames}\longrightarrow
\text{F-series products}\longrightarrow
\text{distribution algebras and breakdown loci}.
\end{gathered}
\]
The contour/Mellin formulation branches from the one-dimensional numen before
returning to the spectral programme through its Dirichlet series.

\subsection{Supersession and withdrawal rules}

The 2026 manual explicitly warns that earlier definitions of Hydra maps may
be inconsistent with its definitions.  Accordingly, this edition does not
splice the 2019, dissertation, and 2026 definitions into an artificial single
quotation.  It gives the 2026 object first and then records the older
specializations.

The withdrawn 2111.07882 paper is not used to certify the number-ring
generalization.  The later 2601.17030 manual supplies a new formulation on
$\mathcal O_K$ and is treated independently.  Similarly, 2002.12762v1 is a
historical predecessor to the numen formulation, not a coequal canonical
statement of the current theory.

\begin{status}[Coverage of the edition]
All ten arXiv lineages, the materially distinct earlier versions, all four
numbered Web installments, and the three public dissertation witnesses are
represented in the corpus register or apparatus.  The mathematical
exposition gives full definitions and the principal theorem chains of the
non-withdrawn research strands.  The two withdrawn works are described only
to the point needed to identify what was superseded or invalidated.
Author-site posts are used as explicit version or expository witnesses, not
as silent replacements for formal proofs.  General teaching notes lie
outside this research edition; the specialized $p$-adic background notes are
apparatus rather than independent core works.
\end{status}

\section{Arithmetic dynamics before the encoding}
\label{sec:arithmetic-dynamics}

The numen formalism does not begin with a transform.  It begins with an
ordinary discrete dynamical system and with the elementary distinction
between a bounded recurrent trajectory and a trajectory which escapes every
bounded subset.  This section fixes those notions before any coding is
introduced.

\subsection{Forward orbits and eventual repetition}

Let $Y$ be a set and let $T:Y\to Y$.  Write $T^{\circ 0}=\operatorname{id}_Y$
and $T^{\circ(n+1)}=T\circ T^{\circ n}$.

\begin{definition}
For $y\in Y$, the \emph{forward orbit} is
\[
  \mathcal O_T^+(y)=\bigl(T^{\circ n}(y)\bigr)_{n\geq 0}.
\]
The point $y$ is \emph{periodic} if $T^{\circ m}(y)=y$ for some $m\geq1$,
and \emph{preperiodic} if $T^{\circ a}(y)$ is periodic for some $a\geq0$.
It is \emph{strictly preperiodic} if it is preperiodic but not periodic.
The least positive period of a periodic point is its \emph{exact period}.
\end{definition}

\begin{lemma}[eventual repetition]
\label{lem:eventual-repetition}
For $y\in Y$, the following are equivalent.
\begin{enumerate}[label=\textup{(\roman*)}]
\item $y$ is preperiodic;
\item $T^{\circ m}(y)=T^{\circ n}(y)$ for some $0\leq m<n$;
\item the set $\{T^{\circ n}(y):n\geq0\}$ is finite.
\end{enumerate}
\end{lemma}

\begin{proof}
If $T^{\circ m}(y)=T^{\circ n}(y)$ with $m<n$, then
$T^{\circ m}(y)$ is periodic of period dividing $n-m$.  This proves
(ii)$\Rightarrow$(i).  A preperiodic orbit consists of a finite initial tail
followed by a finite cycle, proving (i)$\Rightarrow$(iii).  Finally, an
infinite sequence taking values in a finite set repeats a value, so
(iii)$\Rightarrow$(ii).
\end{proof}

The relation used in Siegel's latest manual is
\[
 x\sim_T y
 \quad\Longleftrightarrow\quad
 T^{\circ m}(x)=T^{\circ n}(y)
 \quad\text{for some }m,n\geq0.
\]
It is an equivalence relation: symmetry and reflexivity are immediate, and
transitivity follows by iterating the two common-forward-image equalities
until their middle iterates agree.  Its equivalence classes are the weak
connected components of the functional graph of $T$.

\subsection{The compact-bounded dichotomy}

Suppose now that $Y$ is contained in a normed finite-dimensional real vector
space $V$.  In the arithmetic applications the unconditional locally finite
model is the image of the ring of integers of a number field under its full
Minkowski embedding.  The image under one archimedean embedding need not be
locally finite when the field has degree greater than one.  A subset
$Y\subset V$ is \emph{locally finite} when
$Y\cap B$ is finite for every bounded $B\subset V$.

\begin{proposition}[the full Minkowski image is a lattice]
\label{prop:minkowski-lattice}
Let $K$ be a number field of degree $d=r_1+2r_2$.  Choose its real
embeddings $\sigma_1,\ldots,\sigma_{r_1}$ and one embedding
$\tau_j$ from each pair of complex-conjugate embeddings.  Then
\[
 \iota_K:K\longrightarrow\R^{r_1}\times\C^{r_2}\cong\R^d,
 \qquad
 x\longmapsto
 (\sigma_1(x),\ldots,\sigma_{r_1}(x),
  \tau_1(x),\ldots,\tau_{r_2}(x))
\]
maps $\OK$ onto a full lattice.  In particular, $\iota_K(\OK)$ is locally
finite.
\end{proposition}

\begin{proof}
Choose an integral basis $\omega_1,\ldots,\omega_d$ of $\OK$.  The matrix
whose columns are the images of these basis elements under all $d$
embeddings has nonzero determinant: its squared determinant is the nonzero
field discriminant of the basis.  Replacing each conjugate pair of complex
coordinates by its real and imaginary parts is an invertible real-linear
change of coordinates.  Hence the real-linear map
\[
 A:\R^d\longrightarrow\R^{r_1}\times\C^{r_2},
 \qquad e_j\longmapsto\iota_K(\omega_j),
\]
is an isomorphism, and
$\iota_K(\OK)=A(\Z^d)$.  If $B$ is bounded, then $A^{-1}(B)$ is bounded and
therefore meets $\Z^d$ in finitely many points.  Thus $B$ meets
$\iota_K(\OK)$ in finitely many points.
\end{proof}

\begin{proposition}[preperiodic or escaping]
\label{prop:preperiodic-or-escaping}
Let $Y\subset V$ be locally finite and let $T:Y\to Y$.  For every $y\in Y$
exactly one of the following occurs:
\begin{enumerate}[label=\textup{(\roman*)}]
\item $y$ is preperiodic;
\item $\|T^{\circ n}(y)\|\to\infty$ as $n\to\infty$.
\end{enumerate}
\end{proposition}

\begin{proof}
The alternatives are disjoint because a preperiodic orbit is finite.  If
$y$ is not preperiodic, \cref{lem:eventual-repetition} says that all members
of its forward orbit are distinct.  Fix $R>0$.  The set
$Y\cap\{v:\|v\|\leq R\}$ is finite, so only finitely many of those distinct
orbit points can lie in it.  Hence there is $N_R$ such that
$\|T^{\circ n}(y)\|>R$ for all $n\geq N_R$.  Since $R$ was arbitrary, the
norms tend to infinity.
\end{proof}

\begin{sourceissue}[the discreteness hypothesis in the 2026 statement]
The source states \cref{prop:preperiodic-or-escaping} for a uniformly
discrete subset of an arbitrary valued field and infers that every bounded
ball meets that subset in finitely many points.  Uniform separation alone
does not imply this in a non-proper metric space.  The proof is valid under
the locally finite hypothesis above; in particular it is valid for the full
Minkowski lattice.  It is not automatically valid after projection to one
archimedean embedding, where bounded sets may contain infinitely many
algebraic integers.  This edition does not silently retain either the more
general but insufficient source hypothesis or that single-embedding
specialization without an added discreteness assumption.
\end{sourceissue}

\subsection{The Collatz and shortened
\texorpdfstring{$qx+1$}{qx+1} maps}

The classical Collatz map on $\Z$ is
\[
 C(n)=
 \begin{cases}
 n/2,&n\equiv0\pmod2,\\
 3n+1,&n\equiv1\pmod2.
 \end{cases}
\]
Its odd branch is followed by an even value.  Compressing that forced even
step gives the shortened map
\[
 T_3(n)=
 \begin{cases}
 n/2,&n\equiv0\pmod2,\\
 (3n+1)/2,&n\equiv1\pmod2.
 \end{cases}
\]
More generally, for an odd integer $q$, put
\begin{equation}
\label{eq:Tq}
 T_q(n)=
 \begin{cases}
 n/2,&n\equiv0\pmod2,\\
 (qn+1)/2,&n\equiv1\pmod2.
 \end{cases}
\end{equation}
The restriction to odd $q$ is structural: it makes the second expression an
integer exactly on the odd residue class.  The parity itinerary of $n$ is
\[
 \varepsilon_k(n)=[T_q^{\circ k}(n)]_2\in\{0,1\},\qquad k\geq0,
\]
where $[a]_2$ denotes the standard residue representative.  For a finite
orbit segment of length $m$ we package the itinerary as
\[
 \boldsymbol\varepsilon_m(n)
   =(\varepsilon_0(n),\ldots,\varepsilon_{m-1}(n)).
\]

There are two orders in the corpus which must not be conflated.  The entries
of an itinerary are chronological: the branch indexed by $\varepsilon_0$ is
applied first.  The numen word $\mathbf j=(j_0,\ldots,j_{m-1})$ is used in
the ordered affine composition
\[
 H_{\mathbf j}=H_{j_0}\circ\cdots\circ H_{j_{m-1}},
\]
so the rightmost branch acts first.  Consequently the numen word attached to
a chronological orbit segment is its reversal unless a statement explicitly
adopts the opposite composition convention.  Every order-sensitive formula
below is derived rather than inferred from prose.

\subsection{What the coding is required to preserve}

The goal is not merely to replace an integer orbit by a digit string.  A
usable coding must simultaneously retain:
\begin{enumerate}[label=\textup{(\alph*)}]
\item the residue class which makes each branch the genuine branch of the
      piecewise map;
\item the linear part of a branch word;
\item the translation part of that branch word;
\item convergence at a specified completion when an infinite digit string
      is substituted; and
\item enough branch correctness to turn a fixed point of a formal branch
      composition back into a periodic point of the original map.
\end{enumerate}
Hydra maps provide (a), the pair $(M_H,X_H)$ provides (b)--(c), the extension
lemma provides (d), and integrality supplies the principal mechanism for (e).

\section{Hydra maps and their affine word calculus}
\label{sec:hydra-maps}

The 2026 manual is the controlling source for the word \emph{Hydra}.  Earlier
papers use related residue-class maps, but not always the same formal data.
We first type the current object, then isolate an internal inconsistency in
the source definition and state the convention used in the rest of the
edition.

\subsection{Global-field data}

Let $\F$ be either $\Q$ or a rational function field $\F_q(T)$, and let $K$
be a finite extension of $\F$.  Write $\OK$ for the corresponding ring of
integers and $d=[K:\F]$.  If $\ell$ is a place of $K$, let $K_\ell$ be the
completion and $|\cdot|_\ell$ a fixed representative absolute value.  For
$A\in\End_\F(K)$ put
\begin{equation}
\label{eq:place-operator-seminorm}
 \|A\|_\ell=\sup_{x\in K,\,x\ne0}\frac{|Ax|_\ell}{|x|_\ell}
 \in[0,\infty].
\end{equation}
This quantity is finite exactly when $A$ is continuous for the $\ell$-adic
topology, in which case $A$ extends uniquely to a bounded operator on
$K_\ell$.  All placewise operator statements below include this finiteness
hypothesis.  A general $\F$-linear endomorphism need not be continuous at a
chosen single place: for example, a field automorphism can interchange two
places over the same place of $\F$.  Continuity is automatic when the map is
already typed on the selected local factor, as for scalar multiplication.

Let $\Lambda\subsetneq\OK$ be a nonzero ideal.  Its quotient is finite; set
\[
 p=|\OK/\Lambda|\geq2.
\]
Choose an enumeration of the cosets, with the zero coset indexed by $0$:
\[
 \OK/\Lambda=\{\Lambda_0,\ldots,\Lambda_{p-1}\},\qquad
 \Lambda_0=\Lambda.
\]
This enumeration is additional data.  It supplies a bijection from residue
classes to the digit alphabet $\mathcal A_p=\{0,\ldots,p-1\}$; there is no
canonical digit order for a general ideal quotient.

\begin{definition}[affine $\F$-automorphism]
The affine group used by Siegel is the group of affine $\F$-linear
automorphisms of the affine space underlying $K$:
\[
 \Aff_\F(K)=\operatorname{GL}_\F(K)\ltimes K
            =\End_\F(K)^\times\ltimes K.
\]
An element $(r,c)$ acts by $x\mapsto rx+c$, where
$r\in\End_\F(K)^\times$ and $c\in K$.  Its multiplication law is
\begin{equation}
\label{eq:affine-product}
 (r,c)(s,d)=(r\circ s,\;r(d)+c),
\end{equation}
and its inverse is $(r,c)^{-1}=(r^{-1},-r^{-1}c)$.
\end{definition}

The use of $\End_\F(K)$ rather than scalar multiplication by elements of $K$
is essential.  When $[K:\F]>1$, an $\F$-linear map need not be $K$-linear.
The controlling source calls these maps ``affine lattice morphisms.''  That
name is not used as a type assertion here: an arbitrary element of
$\Aff_\F(K)$ need not preserve $\OK$, $\Lambda$, or any lattice in $K$.
Lattice preservation is always stated separately when it is needed.

\begin{definition}[pre-Hydra]
A \emph{pre-Hydra} is a tuple
\[
 \mathcal H=(K,\Lambda,(H_j)_{j\in\mathcal A_p}),
 \qquad H_j=(r_j,c_j)\in\Aff_\F(K).
\]
The $H_j$ are its branches.  The associated piecewise expression, whenever
it is integer-valued on the selected residue classes, is
\begin{equation}
\label{eq:piecewise-hydra}
 H(x)=\sum_{j=0}^{p-1}\ind{x\in\Lambda_j}\bigl(r_jx+c_j\bigr),
 \qquad x\in\OK.
\end{equation}
\end{definition}

\begin{definition}[admissible and exact-integral]
A pre-Hydra is \emph{admissible} when
\begin{equation}
\label{eq:branch-admissibility}
 H_j(\Lambda_j)\subseteq\OK
 \qquad(0\leq j<p).
\end{equation}
It is \emph{exact-integral} when, more strongly,
\begin{equation}
\label{eq:exact-integrality}
 H_j(x)\in\OK\quad\Longleftrightarrow\quad x\in\Lambda_j
 \qquad(x\in\OK,\ 0\leq j<p).
\end{equation}
This is the domain used in Definition 4.5 of the controlling source.  No
field-wide predicate over $x\in K$ is substituted for it.  An admissible
pre-Hydra defines a map $H:\OK\to\OK$ by
\cref{eq:piecewise-hydra}; this is the Hydra map used below.
\end{definition}

\begin{sourceissue}[the two integrality clauses in the manual]
In Definition 4.3 of arXiv:2601.17030v3, the displayed Hydra is required to
satisfy $r_jz+c_j\in\OK$ for every $z\in\OK$ and every $j$.  Definition 4.5
then calls a Hydra integral when $H_j(z)\in\OK$ \emph{if and only if}
$z\in\Lambda_j$, with $z\in\OK$.  For $p\geq2$ these conditions cannot both
hold: after fixing $j$, choose $z\in\OK\setminus\Lambda_j$.  The first clause
puts $H_j(z)$ in $\OK$, while the second forbids it.  The first clause also
fails for the manual's own principal example $T_3$, whose half-integral
branch formulas are integer-valued only on their assigned parity classes.
The minimal condition needed to define $H:\OK\to\OK$ is
\cref{eq:branch-admissibility}; the iff condition
\cref{eq:exact-integrality} is the stronger branch-correctness property used
in the Correspondence Principle.  These two source clauses and their common
integral domain are preserved rather than replaced by a field-wide
predicate.
\end{sourceissue}

\begin{definition}[proper and centered]
A Hydra is
\begin{enumerate}[label=\textup{(\roman*)}]
\item \emph{proper} if $I-r_0\in\End_\F(K)$ is invertible;
\item \emph{centered} if $H_0(0)=0$, equivalently $c_0=0$.
\end{enumerate}
\end{definition}

Properness gives a unique fixed point for the affine branch labelled $0$,
namely $(I-r_0)^{-1}c_0$, but neither makes that point zero nor makes a
translation by it preserve the branch labelled $0$.  The residue-class
permutation must also be tracked.

\begin{proposition}[translation conjugacy and the zero-class branch]
\label{prop:translation-conjugacy-cosets}
Let $a\in\OK$, let $\tau_a(x)=x+a$, and define
$H^{(a)}=\tau_a^{-1}\circ H\circ\tau_a:\OK\to\OK$.  Translation induces the
permutation $\pi_a$ of $\mathcal A_p$ determined uniquely by
\begin{equation}
\label{eq:translation-coset-permutation}
 \Lambda_{\pi_a(j)}=\Lambda_j+a.
\end{equation}
For $x\in\Lambda_j$, the branch of $H^{(a)}$ on $\Lambda_j$ is
\begin{equation}
\label{eq:translated-Hydra-branch}
 H^{(a)}_j(x)
 =\tau_a^{-1}H_{\pi_a(j)}\tau_a(x)
 =r_{\pi_a(j)}x+
   \bigl(r_{\pi_a(j)}a+c_{\pi_a(j)}-a\bigr).
\end{equation}
Admissibility is preserved by this conjugacy, and exact-integrality is
preserved when $H$ is exact-integral.  In the label-dependent terminology of
the definition, $H^{(a)}$ is proper if and only if
$I-r_{\pi_a(0)}$ is invertible.  The branch of $H^{(a)}$ on the zero class
$\Lambda$ is
centered if and only if
\begin{equation}
\label{eq:translation-centering-condition}
 H_{\pi_a(0)}(a)=a.
\end{equation}
In particular, if $H$ is proper and
$a_0=(I-r_0)^{-1}c_0\in\OK$, then conjugation by $\tau_{a_0}$ is guaranteed
to center the zero-class branch when $a_0\in\Lambda$.  If
$a_0\notin\Lambda$, the equation $H_0(a_0)=a_0$ alone gives no such
conclusion; the exact additional condition is
$H_{\pi_{a_0}(0)}(a_0)=a_0$.
\end{proposition}

\begin{proof}
Addition by $a$ is a bijection of $\OK$ carrying each coset
$\Lambda_j$ onto the coset $\Lambda_j+a$.  Since the displayed enumeration
contains every coset exactly once, \cref{eq:translation-coset-permutation}
defines a permutation, with inverse $\pi_{-a}$.  If $x\in\Lambda_j$, then
$x+a\in\Lambda_{\pi_a(j)}$, so the piecewise definition of $H$ selects
$H_{\pi_a(j)}$.  Applying $\tau_a^{-1}$ and expanding the affine map proves
\cref{eq:translated-Hydra-branch}.  If $x\in\Lambda_j$, admissibility of $H$
puts $H_{\pi_a(j)}(x+a)$ in $\OK$, so subtracting $a\in\OK$ proves
admissibility of $H^{(a)}$.  If $H$ is exact-integral, then, for $x\in\OK$,
\[
 H^{(a)}_j(x)\in\OK
 \Longleftrightarrow H_{\pi_a(j)}(x+a)\in\OK
 \Longleftrightarrow x+a\in\Lambda_{\pi_a(j)}
 \Longleftrightarrow x\in\Lambda_j,
\]
which proves exact-integrality.  At $j=0$, centeredness means precisely
$H^{(a)}_0(0)=0$; by \cref{eq:translated-Hydra-branch}, this is equivalent to
$r_{\pi_a(0)}a+c_{\pi_a(0)}-a=0$, which is
\cref{eq:translation-centering-condition}.  The linear part of that same
zero-class branch is $r_{\pi_a(0)}$, proving the properness criterion.  Finally,
$\pi_{a_0}(0)=0$ exactly when $\Lambda+a_0=\Lambda$, equivalently
$a_0\in\Lambda$; in that case the fixed-point equation for $H_0$ supplies
the required equality.
\end{proof}

\subsection{Exact scalar prior-art boundary}

The scalar integer specialization is not a new class of maps under a new
name.  Matthews's generalized Collatz maps already have exactly the same
fixed-modulus residue-class presentation, while Kohl's residue-class-wise
affine maps give the ambient mapping monoid
\cite[Sec.~2]{MatthewsGeneralized}\cite[Defs.~1.1 and~1.4]{KohlRCWA}.
The extra global-field data, non-scalar endomorphisms, centeredness, and
properness in Siegel's formalism remain additional structure.  In scalar
coordinates, exact-integrality is a predicate of the displayed presentation,
whereas Matthews attaches \emph{relatively prime type} to the map at its least
modulus.  The following comparison proves the exact relation without
identifying those two levels prematurely.

\begin{proposition}[fixed-modulus scalar Hydra--Matthews equivalence]
\label{prop:scalar-hydra-matthews}
Fix $D\in\Z_{\geq2}$, put $K=\F=\Q$, $\OK=\Z$, $\Lambda=D\Z$, and use the standard
classes $\Lambda_i=i+D\Z$ for $0\leq i<D$.  Let
\[
 H_i(x)=\alpha_i x+\beta_i,
 \qquad \alpha_i\in\Q^\times,\quad \beta_i\in\Q.
\]
Let $H:\Z\to\Q$ be the associated piecewise function
\[
 H(x)=H_i(x)\quad\text{for }x\in i+D\Z.
\]
Then the following are equivalent for each $i$.
\begin{enumerate}[label=\textup{(\roman*)}]
\item $H_i(i+D\Z)\subseteq\Z$.
\item There are unique integers
      $m_i\ne0$ and $\rho_i$ such that
      \begin{equation}
      \label{eq:matthews-scalar-branch}
       \rho_i\equiv i m_i\pmod D,
       \qquad
       H_i(x)=\frac{m_i x-\rho_i}{D}
       \quad(x\in\Q).
      \end{equation}
\end{enumerate}
Consequently an admissible scalar Hydra with this fixed standard modulus is
exactly a Matthews generalized Collatz presentation with the same displayed
modulus, and its associated map has type $H:\Z\to\Z$.  Write
$d_{\mathrm{Mat}}(H)$ for Matthews's least affine residue modulus, and put
$g_i=\gcd(D,m_i)$.  The full integral locus of the $i$th formal
branch is
\[
 \bigl\{x\in\Z:H_i(x)\in\Z\bigr\}
 =i+\frac{D}{g_i}\Z.
\]
Therefore
\begin{equation}
\label{eq:scalar-exact-integral-relatively-prime}
 \begin{aligned}
 &\text{the scalar Hydra is exact-integral}\\
 &\quad\Longleftrightarrow\quad
   \gcd(D,m_i)=1\ \text{for every }i\\
 &\quad\Longleftrightarrow\quad
   \gcd(m_0\cdots m_{D-1},D)=1.
 \end{aligned}
\end{equation}
The last condition is the relatively-prime numerical condition for this
displayed $D$-presentation.  Matthews assigns the name
\emph{relatively prime type} using the least modulus
$d_{\mathrm{Mat}}(H)$ \cite[Sec.~2]{MatthewsGeneralized}; consequently it is
the same map-level condition only when $D=d_{\mathrm{Mat}}(H)$.

The qualifier is necessary.  The shortened $3x+1$ map has least-modulus
presentation $D=2$ with multipliers $(1,3)$, hence is of Matthews relatively
prime type.  Refining the same map to $D=4$ gives multipliers
$(2,6,2,6)$ and remainder parameters $(0,-2,0,-2)$ in
\cref{eq:matthews-scalar-branch}.  Each formal branch is then integral on a
whole parity class, not on only its assigned class modulo $4$; the refined
$D=4$ presentation is not exact-integral.
\end{proposition}

\begin{proof}
Assume (i).  Since both $i$ and $i+D$ belong to the assigned class,
\[
 m_i:=D\alpha_i=H_i(i+D)-H_i(i)\in\Z.
\]
The branch is an affine automorphism of $\Q$, so $m_i\ne0$.  Put
\[
 \rho_i:=m_i i-DH_i(i)=-D\beta_i.
\]
This is an integer and satisfies
$\rho_i\equiv m_i i\pmod D$; expanding the definition gives
\cref{eq:matthews-scalar-branch}.  The coefficients of an affine function
make $m_i=D\alpha_i$ and $\rho_i=-D\beta_i$ unique.

Conversely, if (ii) holds and $x=i+Dt$, then
\[
 H_i(x)=\frac{m_i i-\rho_i}{D}+m_it\in\Z,
\]
  so (i) holds.  Applying the equivalence independently to all $D$ classes
identifies the two fixed-modulus presentations.

For $x\in\Z$, the congruence in
\cref{eq:matthews-scalar-branch} gives
\begin{align*}
 H_i(x)\in\Z
 &\Longleftrightarrow D\mid m_i x-\rho_i\\
 &\Longleftrightarrow D\mid m_i(x-i)\\
 &\Longleftrightarrow \frac{D}{g_i}\mid x-i.
\end{align*}
This proves the displayed integral locus.  It equals the assigned class
$i+D\Z$ exactly when $g_i=1$.  Requiring this for all $i$ is equivalent to
$\gcd(m_0\cdots m_{D-1},D)=1$, which proves
\cref{eq:scalar-exact-integral-relatively-prime} and the source-term
identification.
\end{proof}

\begin{proposition}[typed inclusion in the rcwa mapping monoid]
\label{prop:scalar-hydra-rcwa}
Inherit $D\in\Z_{\geq2}$ and the branch data from
\cref{prop:scalar-hydra-matthews}, assume its equivalent conditions hold for
every $0\leq i<D$, and define the piecewise map
\[
 H:\Z\longrightarrow\Z,
 \qquad H(x)=H_i(x)\quad\text{when }x\equiv i\pmod D.
\]
Then $H$ is a residue-class-wise affine (rcwa) mapping in Kohl's sense.  If
$d_{\mathrm{Mat}}(H)$ is Matthews's least affine residue modulus and
$\operatorname{Mod}(H)$ is Kohl's least rcwa modulus, then
\begin{equation}
\label{eq:rcwa-modulus-divides-hydra-modulus}
 d_{\mathrm{Mat}}(H)=\operatorname{Mod}(H)\mid D.
\end{equation}
If the displayed scalar Hydra is exact-integral, then
\begin{equation}
\label{eq:exact-integral-minimal-modulus}
 d_{\mathrm{Mat}}(H)=\operatorname{Mod}(H)=D.
\end{equation}
More precisely,
\begin{equation}
\label{eq:scalar-exact-integral-minimal-relative-type}
 \text{the displayed $D$-presentation is exact-integral}
 \quad\Longleftrightarrow\quad
 \left\{
 \begin{array}{l}
  D=d_{\mathrm{Mat}}(H),\\
  H\text{ is of Matthews relatively prime type}.
 \end{array}
 \right.
\end{equation}
Finally,
\[
 H\in\operatorname{RCWA}(\Z)
 \quad\Longleftrightarrow\quad
 H:\Z\to\Z\text{ is bijective},
\]
because $\operatorname{RCWA}(\Z)$ denotes Kohl's group of rcwa
\emph{permutations}, not the full rcwa mapping monoid.
\end{proposition}

\begin{proof}
On the class $i+D\Z$, formula
\cref{eq:matthews-scalar-branch} has Kohl's integer-coefficient form
$(a_i x+b_i)/c_i$ with $(a_i,b_i,c_i)=(m_i,-\rho_i,D)$; hence $D$ is an
rcwa presentation modulus.

Matthews and Kohl minimize over the same positive integers in this scalar
setting.  Matthews asks that the restriction to each residue class be affine;
Kohl writes the same rational-affine restriction as $(ax+b)/c$ with integer
coefficients, which is possible after clearing denominators.  Thus
$d_{\mathrm{Mat}}(H)=\operatorname{Mod}(H)$.

For completeness, let $m=\operatorname{Mod}(H)$ and
$g=\gcd(m,D)$.  Fix a class modulo $g$.  Every class modulo $m$ and every
class modulo $D$ lying above it have a nonempty infinite intersection by the
Chinese remainder theorem.  The affine formulas from the two presentations
agree with $H$ on that intersection, hence agree as rational affine functions.
It follows that all formulas above the fixed class modulo $g$ are identical;
therefore $g$ is also an rcwa presentation modulus.  Minimality of $m$ and
$g\leq m$ force $g=m$, proving $m\mid D$.  The last equivalence is precisely
Kohl's definition of the group $\operatorname{RCWA}(\Z)$.

Now assume exact-integrality.  If $m<D$, then $m\mid D$ and the two distinct
$D$-classes $D\Z$ and $m+D\Z$ lie in the single class $m\Z$.  Let $A$ be the
rational-affine formula of a modulus-$m$ presentation on $m\Z$.  It agrees
with $H_0$ on the infinite set $D\Z$ and with $H_m$ on the infinite set
$m+D\Z$; equality on an infinite subset forces equality of rational-affine
functions, so $H_0=A=H_m$.  Exact-integrality says that the integral locus of
$H_0$ is $D\Z$ and that of $H_m$ is $m+D\Z$.  One function cannot have these
two distinct integral loci, a contradiction.  Hence $m=D$, proving
\cref{eq:exact-integral-minimal-modulus}.  Together with
\cref{eq:scalar-exact-integral-relatively-prime}, this proves the forward
implication in
\cref{eq:scalar-exact-integral-minimal-relative-type}: the displayed modulus
is now the least modulus, and its multipliers satisfy Matthews's defining
coprimality condition.  Conversely, if $D=d_{\mathrm{Mat}}(H)$ and $H$ is of
Matthews relatively prime type, its unique least-$D$ multipliers satisfy
$\gcd(m_0\cdots m_{D-1},D)=1$; the earlier equivalence then proves that the
displayed presentation is exact-integral.
\end{proof}

Thus neither ``Hydra'' nor ``rcwa group element'' may be used as an
unqualified synonym for the other.  The shortened Collatz map itself is an
rcwa mapping which is surjective but not injective
\cite[Rem.~1.2]{KohlRCWA}; it therefore lies outside
$\operatorname{RCWA}(\Z)$.  Conversely, rcwa membership alone supplies no
centeredness, properness, exact-integrality, branch-correctness, or
global-field structure.  The scalar equivalence between exact-integrality and
Matthews's relatively-prime condition also requires that the displayed
modulus be the least one, as proved above; it is not a consequence of rcwa
membership alone.

\subsection{The shortened \texorpdfstring{$qx+1$}{qx+1} Hydra}

For odd $q$, \cref{eq:Tq} is a $2$-Hydra over $K=\F=\Q$ with
$\Lambda=2\Z$ and branches
\begin{equation}
\label{eq:Tq-branches}
 H_0(x)=\frac{x}{2},
 \qquad
 H_1(x)=\frac{qx+1}{2}.
\end{equation}
It is centered and proper because $c_0=0$ and $1-r_0=1/2$.  On integral
inputs it has the exact branch-correctness property
\[
 H_0(x)\in\Z\Longleftrightarrow x\in2\Z,
 \qquad
 H_1(x)\in\Z\Longleftrightarrow x\in1+2\Z
 \quad(x\in\Z).
\]
The unshortened Collatz map is not exact-integral in this sense because its
odd branch $3x+1$ is integer-valued on both parity classes.

\subsection{Words and order}

Let $\Sigma_p^*$ be the free monoid of finite words in $\mathcal A_p$,
including the empty word $\varnothing$.  Concatenation is denoted by
$\mathbf i\wedge\mathbf j$.  For
$\mathbf j=(j_0,\ldots,j_{m-1})$, set
\begin{equation}
\label{eq:word-composition}
 H_{\mathbf j}=H_{j_0}\circ H_{j_1}\circ\cdots\circ H_{j_{m-1}},
 \qquad H_\varnothing=\operatorname{id}_K.
\end{equation}
Thus
\begin{equation}
\label{eq:concat-composition}
 H_{\mathbf i\wedge\mathbf j}=H_{\mathbf i}\circ H_{\mathbf j}.
\end{equation}
There are unique maps
\[
 M_H:\Sigma_p^*\to\End_\F(K)^\times,
 \qquad X_H:\Sigma_p^*\to K
\]
such that
\begin{equation}
\label{eq:affine-word-decomposition}
 H_{\mathbf j}(x)=M_H(\mathbf j)x+X_H(\mathbf j).
\end{equation}
The second component is Siegel's \emph{numen}; equivalently,
$X_H(\mathbf j)=H_{\mathbf j}(0)$.

\begin{proposition}[affine concatenation identities]
\label{prop:affine-concatenation}
For all $\mathbf i,\mathbf j\in\Sigma_p^*$,
\begin{align}
 M_H(\mathbf i\wedge\mathbf j)
   &=M_H(\mathbf i)M_H(\mathbf j),\label{eq:M-concat}\\
 X_H(\mathbf i\wedge\mathbf j)
   &=M_H(\mathbf i)X_H(\mathbf j)+X_H(\mathbf i)
     =H_{\mathbf i}\bigl(X_H(\mathbf j)\bigr).
     \label{eq:X-cocycle}
\end{align}
\end{proposition}

\begin{proof}
Substitute \cref{eq:affine-word-decomposition} twice into
\cref{eq:concat-composition}:
\[
 H_{\mathbf i}(H_{\mathbf j}(x))
 =M_H(\mathbf i)M_H(\mathbf j)x
  +M_H(\mathbf i)X_H(\mathbf j)+X_H(\mathbf i).
\]
Uniqueness of the linear and translation parts gives both identities.
\end{proof}

\begin{corollary}[closed word formulas]
\label{cor:closed-word-formulas}
For $\mathbf j=(j_0,\ldots,j_{m-1})$,
\begin{align}
 M_H(\mathbf j)
   &=r_{j_0}r_{j_1}\cdots r_{j_{m-1}},\label{eq:M-word}\\
 X_H(\mathbf j)
   &=\sum_{a=0}^{m-1}
       \bigl(r_{j_0}r_{j_1}\cdots r_{j_{a-1}}\bigr)c_{j_a},
       \label{eq:X-word}
\end{align}
where an empty product is $I$ and products mean composition in the displayed
order.
\end{corollary}

\begin{proof}
Induct on $m$.  Appending the last letter and applying
\cref{eq:X-cocycle} gives the new summand after the existing left prefix;
\cref{eq:M-concat} gives the product formula.
\end{proof}

When the $r_j$ do not commute, the product order in
\cref{eq:M-word,eq:X-word} is indispensable.  Expressions such as
$(1-M_H(\mathbf j))^{-1}$ mean the inverse of an endomorphism of $K$, not
scalar division.

\subsection{Repeated words and affine fixed points}

For a word $\mathbf j$, write $\mathbf j^{\wedge n}$ for its $n$-fold
concatenation, and abbreviate $M=M_H(\mathbf j)$ and $X=X_H(\mathbf j)$.

\begin{proposition}[geometric sum in the endomorphism algebra]
\label{prop:repeated-word}
For $n\geq1$,
\begin{align}
 M_H(\mathbf j^{\wedge n})&=M^n,\\
 X_H(\mathbf j^{\wedge n})&=(I+M+\cdots+M^{n-1})X.\label{eq:operator-geometric}
\end{align}
If $I-M$ is invertible, then
\begin{equation}
\label{eq:operator-geometric-inverse}
 X_H(\mathbf j^{\wedge n})=(I-M^n)(I-M)^{-1}X.
\end{equation}
The affine map $H_{\mathbf j}$ has a unique fixed point precisely when
$I-M$ is invertible, in which case
\begin{equation}
\label{eq:word-fixed-point}
 \operatorname{Fix}(H_{\mathbf j})=\{(I-M)^{-1}X\}.
\end{equation}
\end{proposition}

\begin{proof}
The first two identities follow by induction from
\cref{prop:affine-concatenation}.  Because $M$ commutes with every polynomial
in $M$,
\[
 (I-M)(I+M+\cdots+M^{n-1})=I-M^n,
\]
which gives \cref{eq:operator-geometric-inverse}.  Finally,
$H_{\mathbf j}(z)=z$ is exactly $(I-M)z=X$.
\end{proof}

\begin{remark}
Properness requires $I-r_0$ to be invertible.  It does not imply that
$I-M_H(\mathbf j)$ is invertible for every nonempty word.  The latter is a
separate hypothesis whenever \cref{eq:word-fixed-point} is used.
\end{remark}

\subsection{Correct words versus formal words}

Every word defines an affine map on $K$, but it need not describe a genuine
orbit segment of the piecewise map $H$.  A chronological branch word
$\boldsymbol\varepsilon=(\varepsilon_0,\ldots,\varepsilon_{m-1})$ is
\emph{correct at $z\in\OK$} if
\[
 H^{\circ k}(z)\in\Lambda_{\varepsilon_k}
 \qquad(0\leq k<m).
\]
Then
\begin{equation}
\label{eq:chronological-composition}
 H^{\circ m}(z)
 =H_{\varepsilon_{m-1}}\circ\cdots\circ H_{\varepsilon_0}(z)
 =H_{\operatorname{rev}(\boldsymbol\varepsilon)}(z).
\end{equation}
Exact-integrality is designed to recover correctness from integrality of
successive branch values.  The passage from a formal fixed point of
$H_{\mathbf j}$ to a periodic point of $H$ therefore requires more than the
affine equation alone; the needed arithmetic step is isolated in
\cref{sec:correspondence}.

\section{The numen: finite digits, infinite digits, and place-dependent limits}
\label{sec:numen}

The numen is the translation cocycle of the affine word representation.  Its
extension to infinite digits is a second construction, and the topology in
which that extension converges is part of the data.  Keeping these two facts
separate removes several ambiguities present when the word ``continuation''
is used without a codomain place.

\subsection{Finite base-\texorpdfstring{$p$}{p} words}

For $\mathbf j=(j_0,\ldots,j_{m-1})\in\Sigma_p^*$ define
\begin{equation}
\label{eq:digsum}
 \Dig(\mathbf j)=\sum_{a=0}^{m-1}j_ap^a.
\end{equation}
The least significant digit occurs first.  The map $\Dig:\Sigma_p^*\to
\N_0$ is surjective and is not injective: appending zeroes on the right does
not change its value.  These are the only identifications.

\begin{lemma}[finite-word quotient]
\label{lem:word-quotient}
If $\Dig(\mathbf i)=\Dig(\mathbf j)$, then one of $\mathbf i,\mathbf j$ is
obtained from the other by appending finitely many zeroes, after deleting
zeroes already appended to both.  Consequently
$\Sigma_p^*/\!\sim\;\cong\N_0$, where
$\mathbf i\sim\mathbf j$ means equality of digit sums.
\end{lemma}

\begin{proof}
Every nonzero integer has a unique finite base-$p$ expansion with nonzero
last digit.  The zero integer is represented by the empty word and by words
consisting only of zeroes.  The claim follows.
\end{proof}

\begin{proposition}[descent of the word numen]
\label{prop:numen-descent}
If $H$ is centered, then $X_H(\mathbf j)$ depends only on
$\Dig(\mathbf j)$.  Hence there is a well-defined function
\[
 X_H:\N_0\longrightarrow K,
 \qquad X_H(\Dig(\mathbf j))=X_H(\mathbf j).
\]
\end{proposition}

\begin{proof}
Centeredness gives $H_0(0)=0$.  Therefore
\[
 X_H(\mathbf j\wedge(0))
 =H_{\mathbf j}(H_0(0))=H_{\mathbf j}(0)=X_H(\mathbf j).
\]
Iterate this identity and use \cref{lem:word-quotient}.
\end{proof}

The multiplier requires a separate convention: unlike the centered numen, it
does not descend through the finite-word quotient.

\begin{definition}[canonical integer multiplier]
\label{def:canonical-integer-multiplier}
For $n>0$, let
\[
 \operatorname{dig}_p(n)=(d_0(n),\ldots,d_{L-1}(n))
\]
be the unique shortest base-$p$ word representing $n$, so that
$d_{L-1}(n)\ne0$, and put $\operatorname{dig}_p(0)=\varnothing$.  Define
\begin{equation}
\label{eq:canonical-integer-multiplier}
 M_H(n):=M_H\bigl(\operatorname{dig}_p(n)\bigr),
 \qquad M_H(0)=I.
\end{equation}
Thus $M_H(n)$ always refers to the canonical word, not to an arbitrary word
with digit sum $n$.
\end{definition}

Indeed, the concatenation identity gives
\[
 M_H\bigl(\mathbf j\wedge(0)^{\wedge k}\bigr)
 =M_H(\mathbf j)r_0^k.
\]
Consequently equivalent representatives obtained by appending zeroes need
not have the same multiplier.  The word multiplier descends through all such
representatives if and only if $r_0=I$.

For a noncentered proper Hydra, the natural digit-recursive function still
exists, but it is not the translation term $H_{\mathbf j}(0)$ alone.  Its
terminal value must be the fixed point of the zero branch.

\begin{theorem}[digit-recursive characterization]
\label{thm:digit-recursive-characterization}
Let $(r_j,c_j)_{0\leq j<p}$ be affine branches on $K$.  Functions
$f:\N_0\to K$ satisfying
\begin{equation}
\label{eq:numen-FE-N}
 f(pn+j)=r_jf(n)+c_j
 \qquad(n\in\N_0,\ 0\leq j<p)
\end{equation}
are in bijection with solutions $a\in K$ of
\begin{equation}
\label{eq:zero-fixed-equation}
 (I-r_0)a=c_0.
\end{equation}
The bijection sends $f$ to $f(0)=a$.  If $H$ is proper, the solution is
unique and $a=(I-r_0)^{-1}c_0$.  If $H$ is centered and proper, $a=0$ and
$f=X_H$.
\end{theorem}

\begin{proof}
Putting $n=j=0$ in \cref{eq:numen-FE-N} gives
\cref{eq:zero-fixed-equation}.  Conversely, choose a solution $a$.  Set
$f(0)=a$.  Every positive $n$ has a unique decomposition $n=pm+j$ with
$0\leq j<p$ and $m<n$, so \cref{eq:numen-FE-N} recursively defines one and
only one value $f(n)$.  The recursion automatically satisfies every
displayed equation because the decomposition is unique.

If $\mathbf d=\operatorname{dig}_p(n)=(d_0,\ldots,d_{L-1})$, repeated
substitution gives
\begin{equation}
\label{eq:finite-general-solution}
 f(n)=H_{d_0}\circ\cdots\circ H_{d_{L-1}}(a)
     =M_H(\mathbf d)a+X_H(\mathbf d).
\end{equation}
For centered $H$, appending zeroes leaves the value unchanged and this is
exactly the descended word numen of \cref{prop:numen-descent}.
\end{proof}

Combining \cref{eq:X-word} with the canonical digit word gives the explicit
finite formula
\begin{equation}
\label{eq:numen-finite-digit-formula}
 X_H(n)=\sum_{a=0}^{L-1}
 \bigl(r_{d_0}\cdots r_{d_{a-1}}\bigr)c_{d_a},
 \qquad \operatorname{dig}_p(n)=(d_0,\ldots,d_{L-1}),
\end{equation}
for a centered Hydra, with the empty sum understood when $n=0$.  This
formula, not an analogy with an orbit, is the primitive F-series of the later
papers.

\subsection{The base-\texorpdfstring{$p$}{p} digit compactum}

For any integer $p\geq2$, put
\begin{equation}
\label{eq:composite-p-adics}
 \Z_p^{\mathrm{dig}}=\varprojlim_N\Z/p^N\Z.
\end{equation}
Every element has a unique digit expansion
$\mathfrak z=\sum_{a\geq0}d_a(\mathfrak z)p^a$ with
$d_a(\mathfrak z)\in\mathcal A_p$.  If $p$ is prime this is the usual ring
$\Z_p$.  For composite $p$, \cref{eq:composite-p-adics} is still the compact
base-$p$ digit ring, but it is not a discrete valuation ring and should not
be treated as a field's ring of integers.

\begin{proposition}[digits, the inverse limit, and Haar probability]
\label{prop:digit-compactum-haar}
Give $\mathcal A_p^{\N_0}$ the product topology, with each copy of
$\mathcal A_p$ discrete.  The map
\[
 \Phi_p:\mathcal A_p^{\N_0}\longrightarrow\Z_p^{\mathrm{dig}},
 \qquad
 (d_a)_{a\geq0}\longmapsto
 \left(\sum_{a=0}^{N-1}d_ap^a\bmod p^N\right)_{N\geq1}
\]
is a homeomorphism.  Under $\Phi_p$, normalized Haar probability $m_p$ on
$\Z_p^{\mathrm{dig}}$ is the product of the uniform probability measures on
$\mathcal A_p$.  Equivalently, for every $N\geq1$ and
$e_0,\ldots,e_{N-1}\in\mathcal A_p$,
\begin{equation}
\label{eq:digit-cylinder-mass}
 m_p\{\mathfrak z:d_0(\mathfrak z)=e_0,\ldots,
                   d_{N-1}(\mathfrak z)=e_{N-1}\}=p^{-N}.
\end{equation}
Thus the coordinate digits are independent and each is uniform on
$\mathcal A_p$.
\end{proposition}

\begin{proof}
The displayed residues are compatible under reduction, so $\Phi_p$ is
well-defined.  Conversely, take a compatible sequence
$(x_N\bmod p^N)_N$.  Its reduction modulo $p$ determines a unique
$d_0\in\mathcal A_p$.  After $d_0,\ldots,d_{N-1}$ have been determined,
compatibility makes
\[
 x_{N+1}-\sum_{a=0}^{N-1}d_ap^a
\]
divisible by $p^N$, and its quotient modulo $p$ determines the unique next
digit $d_N$.  This constructs the inverse of $\Phi_p$.  Both maps are
continuous because each output coordinate or digit depends on only finitely
many input coordinates; alternatively, the continuous bijection is a
homeomorphism because the product is compact and the inverse limit is
Hausdorff.

Fixing the first $N$ digits specifies one coset of the open subgroup
$p^N\Z_p^{\mathrm{dig}}$.  That subgroup has index $p^N$, so translation
invariance and total mass one give every such cylinder mass $p^{-N}$.
These are exactly the finite-dimensional cylinder probabilities of the
uniform product measure, proving the measure identity and independence.
\end{proof}

Write
\begin{align}
 [\mathfrak z]_{p^N}&=\sum_{a=0}^{N-1}d_a(\mathfrak z)p^a,\\
 \sh(\mathfrak z)&=\frac{\mathfrak z-d_0(\mathfrak z)}p
                  =\sum_{a\geq0}d_{a+1}(\mathfrak z)p^a.
\end{align}
The first is the canonical truncation and the second is the left digit shift.
The digit and measure assertions used below are therefore consequences of
\cref{prop:digit-compactum-haar}, including when $p$ is composite.

\subsection{The primitive F-series}

Assume throughout this subsection that $H$ is centered.  Fix a place $\ell$
of $K$ and view all $r_j$ as continuous operators on the completion
$K_\ell$.  Let
\begin{equation}
\label{eq:partial-F-series}
 S_N(\mathfrak z)=X_H([\mathfrak z]_{p^N})
 =\sum_{a=0}^{N-1}
   \left(\prod_{k=0}^{a-1}r_{d_k(\mathfrak z)}\right)
   c_{d_a(\mathfrak z)}.
\end{equation}
Here and below the empty product is $I$.  When $S_N(\mathfrak z)$ converges
in $K_\ell$, define
\begin{equation}
\label{eq:rising-extension}
 X_{H,\ell}(\mathfrak z)=\lim_{N\to\infty}S_N(\mathfrak z).
\end{equation}
This is a \emph{rising extension}: its values are prescribed by limits along
the canonical truncation sequence.  It is not a topology-free interpolation.

\begin{proposition}[pointwise convergence criterion]
\label{prop:numen-pointwise-convergence}
Let $C=\max_j|c_j|_\ell$.  Suppose that for a given $\mathfrak z$ there are
$\eta>0$ and $A\in\R$ such that
\begin{equation}
\label{eq:negative-lyapunov}
 \sum_{k=0}^{a-1}\log\|r_{d_k(\mathfrak z)}\|_\ell
 \leq A-\eta a
\end{equation}
for all sufficiently large $a$.  Then the series in
\cref{eq:partial-F-series} converges absolutely in an archimedean completion
and converges in a nonarchimedean completion.  The same estimate makes the
remainder term
$\bigl(\prod_{k<N}r_{d_k}\bigr)a_0$ tend to zero for every fixed
$a_0\in K_\ell$.
\end{proposition}

\begin{proof}
The norm of the $a$th summand is at most
$C\exp(A-\eta a)$.  This is summable in the archimedean case and tends to
zero in the nonarchimedean case; a series in a complete nonarchimedean field
converges exactly when its terms tend to zero.  The identical product bound
applies to the remainder.
\end{proof}

\begin{sourceissue}[the pointwise digit-density criterion]
For a fixed digit sequence, write
\[
 \nu_{j,N}=\frac1N\#\{0\leq k<N:d_k(\mathfrak z)=j\},
 \qquad \alpha_j=\log\|r_j\|_\ell.
\]
The upper Lyapunov exponent is controlled by
\begin{equation}
\label{eq:density-lyapunov-upper-bound}
 \limsup_{N\to\infty}\sum_j\alpha_j\nu_{j,N}
 \leq
 \sum_{\alpha_j<0}\alpha_j\liminf_{N\to\infty}\nu_{j,N}
 +\sum_{\alpha_j>0}\alpha_j\limsup_{N\to\infty}\nu_{j,N}.
\end{equation}
For a negative coefficient, an upper bound on $\alpha_j\nu_{j,N}$ requires a
lower bound on the frequency, hence a liminf; multiplication by
$\alpha_j<0$ reverses the inequality.  For a positive coefficient, it
requires an upper bound on the frequency, hence a limsup.  Terms with
$\alpha_j=0$ are immaterial.  The pointwise criterion in the 2026 source
reverses these choices, using a limsup for contractive branches and a liminf
for expansive branches.  Its displayed quantity is therefore a lower, not
an upper, control and its negativity is not a sufficient convergence
condition.  We avoid that error here by assuming the direct eventual bound
\cref{eq:negative-lyapunov}; negativity of the right side of
\cref{eq:density-lyapunov-upper-bound} is one valid way to obtain that bound.
\end{sourceissue}

\begin{theorem}[almost-everywhere extension]
\label{thm:numen-ae-extension}
Assume $H$ is centered and proper and
\begin{equation}
\label{eq:geometric-mean-contraction}
 \prod_{j=0}^{p-1}\|r_j\|_\ell<1.
\end{equation}
Then \cref{eq:rising-extension} converges for Haar-almost every
$\mathfrak z\in\Z_p^{\mathrm{dig}}$ and defines a measurable
$K_\ell$-valued function on its conull domain.  On that domain it obeys
\begin{equation}
\label{eq:numen-FE-adic}
 X_{H,\ell}(p\mathfrak z+j)
 =r_jX_{H,\ell}(\mathfrak z)+c_j.
\end{equation}
\end{theorem}

\begin{proof}
For Haar-almost every digit sequence, the strong law of large numbers gives
\[
 \frac1N\sum_{k=0}^{N-1}\log\|r_{d_k(\mathfrak z)}\|_\ell
 \longrightarrow \frac1p\sum_{j=0}^{p-1}\log\|r_j\|_\ell<0.
\]
Thus \cref{eq:negative-lyapunov} holds eventually, and
\cref{prop:numen-pointwise-convergence} applies.  Each $S_N$ is locally
constant and therefore measurable; its pointwise limit is measurable.
Finally, the digits of $p\mathfrak z+j$ are $j,d_0(\mathfrak z),\ldots$.
Separate the first summand in \cref{eq:partial-F-series} and pass to the
limit to obtain \cref{eq:numen-FE-adic}.
\end{proof}

\begin{theorem}[uniform contraction]
\label{thm:numen-uniform-contraction}
Suppose
\[
 \rho=\max_{0\leq j<p}\|r_j\|_\ell<1.
\]
Then the series in \cref{eq:partial-F-series} converges uniformly on
$\Z_p^{\mathrm{dig}}$.  Its sum is continuous and satisfies
\cref{eq:numen-FE-adic}.  More precisely, in the archimedean case,
\begin{equation}
\label{eq:uniform-tail-bound}
 \sup_{\mathfrak z}|X_{H,\ell}(\mathfrak z)-S_N(\mathfrak z)|_\ell
 \leq \frac{C\rho^N}{1-\rho},
\end{equation}
while in the nonarchimedean case the right side may be replaced by
$C\rho^N$.
\end{theorem}

\begin{proof}
The $a$th summand has norm at most $C\rho^a$, uniformly in
$\mathfrak z$.  The ordinary and ultrametric tail estimates give
\cref{eq:uniform-tail-bound}.  Uniform limits of the locally constant
functions $S_N$ are continuous.  The functional equation follows exactly as
in \cref{thm:numen-ae-extension}.
\end{proof}

\begin{sourceissue}[centeredness and uniqueness in the extension lemma]
The 2026 extension lemma assumes properness but not centeredness, while its
displayed F-series is the translation term $H_{\mathbf j}(0)$ and therefore
uses terminal value $0$.  For a noncentered proper Hydra, the unique solution
on $\N_0$ instead has terminal value
$a=(I-r_0)^{-1}c_0$ and finite values
$M_H(\mathbf d)a+X_H(\mathbf d)$ as in
\cref{eq:finite-general-solution}.  The extra product term disappears under
a contraction hypothesis, but it cannot be omitted formally.  Likewise, a
functional equation on all infinite digit strings need not determine an
arbitrary function without a rising-limit, regularity, or boundedness
condition.  The uniqueness asserted here is uniqueness of the prescribed
rising extension, and under uniform contraction uniqueness in the natural
continuous bounded class.
\end{sourceissue}

\subsection{The shortened \texorpdfstring{$qx+1$}{qx+1} numen}

For the Hydra $T_q$, write $\chi_q=X_{T_q}$.  From
\cref{eq:Tq-branches,eq:numen-FE-N},
\begin{align}
 \chi_q(2n)&=\frac12\chi_q(n),\label{eq:chi-even}\\
 \chi_q(2n+1)&=\frac{q\chi_q(n)+1}{2},\label{eq:chi-odd}
\end{align}
with $\chi_q(0)=0$.  If
$\operatorname{dig}_2(n)=(d_0,\ldots,d_{L-1})$ and
$s_a=d_0+\cdots+d_{a-1}$, then
\begin{align}
 M_q(n)&=\frac{q^{d_0+\cdots+d_{L-1}}}{2^L},
 \label{eq:Mq-explicit}\\
 \chi_q(n)&=\sum_{\substack{0\leq a<L\\d_a=1}}
               \frac{q^{s_a}}{2^{a+1}}.
 \label{eq:chi-explicit}
\end{align}
Here $L$ is the length of the canonical shortest binary word; in particular,
$M_q(0)=1$.  Appending zeroes would multiply the word multiplier by extra
factors of $1/2$, so a padded representative must not be used in
\cref{eq:Mq-explicit}.
For an odd prime $q$, the same series converges in $\Q_q$ for every
$\mathfrak z\in\Z_2$: it is finite when the binary expansion has finitely
many ones, and otherwise its nonzero terms have $q$-adic norms
$q^{-s_a}\to0$.  Thus
\begin{equation}
\label{eq:chi-q-adic-series}
 \chi_q(\mathfrak z)=
 \sum_{\substack{a\geq0\\d_a(\mathfrak z)=1}}
 \frac{q^{s_a(\mathfrak z)}}{2^{a+1}}
 \in\Z_q.
\end{equation}
It is rising-continuous everywhere.  It is $2$-adic-to-$q$-adic continuous
at every binary sequence with infinitely many ones, but not at the
nonnegative integers: if $n\in\N_0$ has $k=\#_1(n)$ ones, a new $1$ placed
arbitrarily far beyond its last digit contributes a term of the fixed
nonzero $q$-adic norm $q^{-k}$, independent of that remote position (it is a
unit only when $n=0$).  This pointwise distinction is the prototype for the
later theory of variable topologies of convergence.

\section{Periodic codes and the Correspondence Principle}
\label{sec:correspondence}

The algebraic heart of the Correspondence Principle is a fixed-point
identity for a repeated word.  Its dynamical heart is branch correctness.
Those are independent steps: the first takes place in an affine group, while
the second identifies a formal composition with an iterate of a piecewise
map.

\subsection{Purely periodic digit strings}

Let $\mathbf j=(j_0,\ldots,j_{L-1})$ be a word of length $L\geq1$, and put
\[
 n=\Dig(\mathbf j)=\sum_{a=0}^{L-1}j_ap^a.
\]
Its infinite repetition represents the base-$p$ digit-ring element
\begin{equation}
\label{eq:Bp}
 B_p(\mathbf j)=B_{p,L}(n)
 =n(1+p^L+p^{2L}+\cdots)=\frac{n}{1-p^L}.
\end{equation}
The length $L$ is part of this expression.  When $\mathbf j$ is the shortest
word for $n>0$, write $\lambda_p(n)=L$ and
$B_p(n)=n/(1-p^{\lambda_p(n)})$.

\begin{theorem}[repeated-code fixed-point identity]
\label{thm:repeated-code-fixed-point}
Let $H$ be centered, let $\ell$ be a place of $K$, and set
$M=M_H(\mathbf j)$ and $X=X_H(\mathbf j)$.  Suppose
$\|M\|_\ell<1$.  Then $I-M$ is invertible on $K_\ell$, the rising numen
converges at $B_p(\mathbf j)$, and
\begin{equation}
\label{eq:repeated-code-value}
 X_{H,\ell}(B_p(\mathbf j))=(I-M)^{-1}X.
\end{equation}
This value is the unique fixed point of $H_{\mathbf j}$ in $K_\ell$.
\end{theorem}

\begin{proof}
The $mL$-digit truncation of $B_p(\mathbf j)$ is represented by
$\mathbf j^{\wedge m}$.  By \cref{prop:repeated-word},
\[
 X_H(\mathbf j^{\wedge m})
 =(I+M+\cdots+M^{m-1})X.
\]
The operator series converges in the complete endomorphism algebra because
$\|M\|_\ell<1$, and its sum is $(I-M)^{-1}$.  Truncations ending inside the
next block have the same limit because the remaining finite prefix is
multiplied by $M^m$.  This proves the rising limit.  The fixed-point assertion
is \cref{eq:word-fixed-point}.
\end{proof}

\begin{remark}
In a scalar one-dimensional setting \cref{eq:repeated-code-value} is often
written $X/(1-M)$.  For a general Hydra, $M$ is an endomorphism and the typed
formula is $(I-M)^{-1}X$.
\end{remark}

\begin{proposition}[rational points are the eventually periodic digit strings]
\label{prop:rational-eventually-periodic}
Under the diagonal embedding into the completions at the primes dividing
$p$,
\[
 \Q\cap\Z_p^{\mathrm{dig}}
 =\Z_{(p)}:=\left\{\frac ab\in\Q:\gcd(b,p)=1
                   \text{ after reducing }a/b\right\}.
\]
An element of $\Z_p^{\mathrm{dig}}$ lies in this subring if and only if its
base-$p$ digit sequence is eventually periodic.
\end{proposition}

\begin{proof}
Write $p=\prod_{\substack{\ell\mid p\\ \ell\ {\mathrm{prime}}}}\ell^{e_\ell}$.
The compatible Chinese-remainder
isomorphisms give
\[
 \Z_p^{\mathrm{dig}}
 =\varprojlim_N\Z/p^N\Z
 \xrightarrow{\sim}
 \prod_{\ell\mid p}\varprojlim_N\Z/\ell^{e_\ell N}\Z
 \xrightarrow{\sim}\prod_{\ell\mid p}\Z_\ell.
\]
The subsequence $(e_\ell N)_{N\geq1}$ is cofinal, which proves the second
isomorphism.  Under the resulting diagonal embedding of $\Q$ into
$\prod_{\ell\mid p}\Q_\ell$, a reduced rational belongs to the displayed
product of integer rings exactly when no prime $\ell\mid p$ divides its
denominator.  Hence
$\Q\cap\Z_p^{\mathrm{dig}}=\Z_{(p)}$ with the meaning stated above.

If a digit string has a prefix of length $R$ with value $u$ and thereafter
repeats a block of length $L$ with value
$v=\sum_{h=0}^{L-1}e_hp^h$, then its value in the inverse limit is
\[
 u+p^R v(1+p^L+p^{2L}+\cdots)
 =u+\frac{p^Rv}{1-p^L}.
\]
Because $1-p^L\equiv1\pmod p$, its denominator is coprime to $p$.

Conversely, let $x=a/b$ with $a,b\in\Z$ and $\gcd(b,p)=1$.  For every
$N\geq0$ there is a unique compatible residue
$A_N\in\{0,\ldots,p^N-1\}$ satisfying
$bA_N\equiv a\pmod{p^N}$; these residues define the digit expansion of $x$.
Here $A_0=0$.  Put
\[
 u_N=\frac{a-bA_N}{p^N}\in\Z.
\]
If $d_N$ is the next digit, then
\begin{equation}
\label{eq:rational-digit-state}
 u_{N+1}=\frac{u_N-bd_N}{p},
 \qquad
 d_N\in\mathcal A_p,\qquad bd_N\equiv u_N\pmod p.
\end{equation}
The congruence determines $d_N$ uniquely from $u_N$ because $b$ is a unit
modulo $p$.  Iterating
$|u_{N+1}|\leq(|u_N|+|b|(p-1))/p$ gives
\[
 |u_N|\leq p^{-N}|u_0|+|b|(1-p^{-N})
 \leq |u_0|+|b|.
\]
Thus the integer sequence $(u_N)$ takes values in a finite set, and the
deterministic transition
\cref{eq:rational-digit-state} must eventually repeat.  Its associated
digits are therefore eventually periodic.
\end{proof}

Let $H$ be centered, fix a place $\ell$ at which all branches are bounded,
and suppose $\|M_H(\mathbf j)\|_\ell<1$.  If $\mathfrak z$ has a finite
prefix $\mathbf i$ followed by the infinite repetition of $\mathbf j$, then
\begin{equation}
\label{eq:eventually-periodic-numen}
 X_{H,\ell}(\mathfrak z)
 =H_{\mathbf i}\bigl((I-M_H(\mathbf j))^{-1}X_H(\mathbf j)\bigr).
\end{equation}
This is an affine image of a formal periodic fixed point.  It becomes a
statement about the actual Hydra only after correctness is proved.

\subsection{Classical \texorpdfstring{$2$}{2}-adic conjugacy and the exact
periodic overlap}

The parity-vector conjugacy predates the numen.  Its digit order is
chronological, whereas Siegel's displayed affine word acts from right to
left.  The difference is not cosmetic: it forces a word reversal in the
precise locus where the two constructions have a common rational value.

Let $q$ be any odd integer and extend $T_q$ to $\Z_2$ using the parity
branches in \cref{eq:Tq-branches}.  Let
\[
 S=\vartheta_2:\Z_2\longrightarrow\Z_2,
 \qquad S(z)=\frac{z-[z]_2}{2},
\]
be the binary digit shift.
For every prime $\ell$, write
$\iota_\ell:\Q\hookrightarrow\Q_\ell$ for the canonical field embedding.

\begin{proposition}[the Bernstein--Lagarias conjugacy, with maps typed]
\label{prop:BL-conjugacy-typed}
Define
\begin{equation}
\label{eq:BL-parity-encoder}
 Q_q(x)=\sum_{i=0}^{\infty}\bigl[T_q^{\circ i}(x)\bigr]_2\,2^i
 \qquad(x\in\Z_2).
\end{equation}
If $z\in\Z_2\setminus\{0\}$, let either $I=\{1,\ldots,r\}$ for some $r\geq1$ or
$I=\Z_{\geq1}$, according as the set of one positions is finite or infinite,
and write
\[
 z=\sum_{\ell\in I}2^{d_\ell},
 \qquad d_1\geq0,
 \qquad d_\ell<d_{\ell+1}\quad(\ell,\ell+1\in I).
\]
Define
\begin{equation}
\label{eq:BL-inverse-formula}
 \Phi_q(z)=-\sum_{\ell\in I}q^{-\ell}2^{d_\ell},
 \qquad \Phi_q(0)=0.
\end{equation}
Then $Q_q,\Phi_q:\Z_2\to\Z_2$ are mutually inverse $2$-adic
isometries, hence solenoidal homeomorphisms, and
\begin{equation}
\label{eq:BL-conjugacy-oriented}
 Q_q\circ T_q=S\circ Q_q,
 \qquad
 T_q=\Phi_q\circ S\circ\Phi_q^{-1}.
\end{equation}
Thus $Q_q$ is the chronological parity encoder and
$\Phi_q=Q_q^{-1}$ is the conjugacy.  This is the odd-$(a,b)=(q,1)$
specialization of Bernstein--Lagarias, not a prime-$q$ statement
\cite[Eqs.~(1.3), (1.5), (1.6), (4.1), and~(4.2)]{BernsteinLagarias}.
\end{proposition}

\begin{proof}
For odd $q$, $H_0$ maps $2\Z_2$ to $\Z_2$ and $H_1$ maps
$1+2\Z_2$ to $\Z_2$.  Thus the branch selected by the current parity is
defined at every iterate, every digit in \cref{eq:BL-parity-encoder} lies in
$\{0,1\}$, and the displayed binary series converges.  If
$x\equiv y\pmod{2^N}$, induction through their common parity
branches gives
\[
 T_q^{\circ i}(x)\equiv T_q^{\circ i}(y)
 \pmod{2^{N-i}}
 \qquad(0\leq i\leq N),
\]
so the first $N$ digits of $Q_q(x)$ and $Q_q(y)$ agree.  Conversely, if
those $N$ digits agree and contain $m$ ones, the same-branch affine
calculation gives
\[
 T_q^{\circ N}(x)-T_q^{\circ N}(y)
 =\frac{q^m}{2^N}(x-y).
\]
The left side lies in $\Z_2$ and $q$ is a $2$-adic unit, hence
$x\equiv y\pmod{2^N}$.  Therefore
\begin{equation}
\label{eq:BL-isometry-congruences}
 x\equiv y\pmod{2^N}
 \quad\Longleftrightarrow\quad
 Q_q(x)\equiv Q_q(y)\pmod{2^N},
\end{equation}
which proves that $Q_q$ is an isometry onto its image.

The finite sum in \cref{eq:BL-inverse-formula} is immediate; in the infinite
case the series converges because
$|q^{-\ell}2^{d_\ell}|_2=2^{-d_\ell}\to0$.  Writing
$z=\delta+2w$ with $\delta\in\{0,1\}$ and separating the first one when
$\delta=1$ gives the exact recursions
\begin{equation}
\label{eq:BL-Phi-recursion}
 \Phi_q(2w)=2\Phi_q(w),
 \qquad
 \Phi_q(1+2w)=\frac{2\Phi_q(w)-1}{q}.
\end{equation}
The first value is even and the second is odd.  Applying the corresponding
branch of $T_q$ therefore gives
\[
 T_q(\Phi_q(z))=\Phi_q(Sz).
\]
Iterating shows that the parity digits of $\Phi_q(z)$ are exactly the digits
of $z$, so $Q_q\Phi_q=\operatorname{id}_{\Z_2}$.  Hence $Q_q$ is
surjective; together with \cref{eq:BL-isometry-congruences}, it is bijective
and $\Phi_q=Q_q^{-1}$.  Dropping the first parity digit in
\cref{eq:BL-parity-encoder} proves $Q_qT_q=SQ_q$, and substitution of the
inverse gives the second conjugacy formula.  The congruence equivalence for
every $N$ proves that both inverse maps are solenoidal isometries and hence
homeomorphisms.
\end{proof}

Bernstein--Lagarias Appendix B is broader still.  Given arbitrary solenoidal
bijections $V_0,V_1:\Z_2\to\Z_2$, it treats
\[
 U_{V_0,V_1}(x)=
 \begin{cases}
 V_0(x/2),&x\equiv0\pmod2,\\
 V_1((x-1)/2),&x\equiv1\pmod2,
 \end{cases}
\]
and proves that its chronological parity encoder is solenoidal bijective and
$U=Q^{-1}SQ$, with a converse classification
\cite[App.~B, Thms.~B.1--B.2]{BernsteinLagarias}.  The shortened $T_q$
has $V_0(y)=y$ and $V_1(y)=qy+(q+1)/2$; the latter is an affine isometry of
$\Z_2$ for every odd $q$, with inverse
$y\mapsto q^{-1}(y-(q+1)/2)$.  This proves the exact inclusion in that
class; it does not identify Appendix B's full class with general Hydras.

\begin{theorem}[reversed purely periodic overlap]
\label{thm:BL-numen-periodic-overlap}
Let $q$ be odd, let $L\geq1$, and let
$\boldsymbol\varepsilon=(\varepsilon_0,\ldots,\varepsilon_{L-1})$ be a
chronological binary block.  If $m\geq1$, write the positions of its $m$ ones
as
\[
 0\leq v_0<\cdots<v_{m-1}<L;
\]
when $m=0$, every position list and sum below is empty.  Set
\begin{equation}
\label{eq:chronological-and-reversed-codes}
 e=\sum_{i=0}^{L-1}\varepsilon_i2^i,
 \qquad
 \operatorname{rev}_L(e)
   =\sum_{i=0}^{L-1}\varepsilon_i2^{L-1-i}.
\end{equation}
Put
\begin{equation}
\label{eq:BL-periodic-rational}
 R_{q,\boldsymbol\varepsilon}
 :=\frac{\displaystyle
         \sum_{k=0}^{m-1}q^{m-k-1}2^{v_k}}
        {2^L-q^m}\in\Q.
\end{equation}
Then
\[
 \Phi_q\!\left(\frac{e}{1-2^L}\right)
 =\iota_2\!\left(R_{q,\boldsymbol\varepsilon}\right)\in\Z_2.
\]
If, in addition, $q$ is an odd prime, so that the edition's global
$\chi_q:\Z_2\to\Z_q$ is available, then
\begin{equation}
\label{eq:BL-numen-overlap}
 \chi_q\!\left(
       \frac{\operatorname{rev}_L(e)}{1-2^L}\right)
 =\iota_q\!\left(R_{q,\boldsymbol\varepsilon}\right)\in\Z_q.
\end{equation}
Thus the two values have the same rational representative
$R_{q,\boldsymbol\varepsilon}$ after reversal.  The reversal in
\cref{eq:BL-numen-overlap} is mandatory.
\end{theorem}

\begin{proof}
If $m=0$, every displayed value is zero, so assume first that $m\geq1$.
The one positions of the purely periodic $2$-adic integer
$e/(1-2^L)$ are
\[
 d_{rm+k+1}=v_k+rL
 \qquad(r\geq0,\ 0\leq k<m).
\]
Substitution into \cref{eq:BL-inverse-formula} and a $2$-adically
convergent geometric sum give
\begin{align*}
 \Phi_q\!\left(\frac{e}{1-2^L}\right)
 &=-\sum_{r\geq0}\sum_{k=0}^{m-1}
       q^{-(rm+k+1)}2^{v_k+rL}\\
 &=-\left(\sum_{k=0}^{m-1}q^{-(k+1)}2^{v_k}\right)
       \sum_{r\geq0}\left(\frac{2^L}{q^m}\right)^r\\
 &=\iota_2\!\left(\frac{\displaystyle
         \sum_{k=0}^{m-1}q^{m-k-1}2^{v_k}}
        {2^L-q^m}\right),
\end{align*}
because $|2^L/q^m|_2=2^{-L}<1$.  This proves
the asserted $\Phi_q$ evaluation with the rational representative defined in
\cref{eq:BL-periodic-rational}.

Now put $h_0(x)=x/2$ and $h_1(x)=(qx+1)/2$.  Induction on the chronological
block gives
\begin{equation}
\label{eq:chronological-affine-return}
 h_{\varepsilon_{L-1}}\circ\cdots\circ h_{\varepsilon_0}(x)
 =\frac{q^m x+
        \displaystyle\sum_{k=0}^{m-1}q^{m-k-1}2^{v_k}}
       {2^L}.
\end{equation}
To see the induction with every index fixed, for $0\leq h\leq L$ put
\[
 m_h=\sum_{i=0}^{h-1}\varepsilon_i,
 \qquad
 C_h=\sum_{\substack{0\leq k<m\\v_k<h}}
       q^{m_h-k-1}2^{v_k}.
\]
The claim after the first $h$ chronological branches is
\[
 h_{\varepsilon_{h-1}}\circ\cdots\circ h_{\varepsilon_0}(x)
 =\frac{q^{m_h}x+C_h}{2^h}.
\]
At $h=0$ the left side denotes the empty composite
$\operatorname{id}_{\Q}$, so the claim is the identity.  For
$0\leq h<L$, assume the claim at level $h$.  If $\varepsilon_h=0$, applying
$h_0$ changes
only the denominator from $2^h$ to $2^{h+1}$, while
$m_{h+1}=m_h$ and $C_{h+1}=C_h$.  If $\varepsilon_h=1$, then
\[
 h_1\!\left(\frac{q^{m_h}x+C_h}{2^h}\right)
 =\frac{q^{m_h+1}x+qC_h+2^h}{2^{h+1}}.
\]
Here $m_{h+1}=m_h+1$; multiplication of $C_h$ by $q$ raises every old
exponent by one, and the new odd position $v_{m_h}=h$ contributes the term
$q^0 2^h$.  Thus the new numerator is exactly $q^{m_{h+1}}x+C_{h+1}$.
At $h=L$ this is \cref{eq:chronological-affine-return}.  Solving its affine
fixed-point equation gives $R_{q,\boldsymbol\varepsilon}$; its denominator is
nonzero because an odd
integer power cannot equal $2^L$.

Under Siegel's convention, the displayed word
$H_{j_0}\circ\cdots\circ H_{j_{L-1}}$ applies $j_{L-1}$ first.  It realizes
the chronological block precisely when
\[
 (j_0,\ldots,j_{L-1})
 =(\varepsilon_{L-1},\ldots,\varepsilon_0),
\]
whose little-endian code is exactly $\operatorname{rev}_L(e)$.  For prime
$q$ and $m\geq1$, its linear part is $q^m/2^L$, of $q$-adic norm $q^{-m}<1$.
The repeated-code theorem therefore evaluates its rising numen at its unique
fixed point, namely $\iota_q(R_{q,\boldsymbol\varepsilon})$.  The $m=0$ case
is again the centered
zero word.  This proves \cref{eq:BL-numen-overlap}.

The reversal is genuinely necessary in the universal formula, not merely a
different way to name the same argument.  Take $q=3$, $L=2$, and
$\boldsymbol\varepsilon=(0,1)$.  Then $e=2$ and
$\operatorname{rev}_2(e)=1$, while the formulas already proved give
\[
 \Phi_3\!\left(\frac{2}{1-4}\right)=\iota_2(2),
 \qquad
 \chi_3\!\left(\frac{1}{1-4}\right)=\iota_3(2),
 \qquad
 \chi_3\!\left(\frac{2}{1-4}\right)=\iota_3(1).
\]
Indeed, $h_1\circ h_0$ is $x\mapsto(3x+2)/4$ with fixed point $2$, whereas
$h_0\circ h_1$ is $x\mapsto(3x+1)/4$ with fixed point $1$.
Thus deleting $\operatorname{rev}_L$ makes the asserted identity false.
\end{proof}

\begin{sourceissue}[the unrestricted B\"ohm--Sontacchi Proposition 4 iterate
identity is false]
\label{status:Bohm-Sontacchi-Proposition-4}
For $q=3$, the right side of \cref{eq:chronological-affine-return} is the
formal branch-composition expression obtained by specializing
B\"ohm--Sontacchi Proposition~1.  Their printed Proposition~4 displays this
expression as $u^n(x)$ while quantifying over every $x\in\Z$ and every run
tuple $(n_0,\ldots,n_m)$.  That unrestricted left side is false.  Take
$x=0$, $m=1$, and $n_0=n_1=1$.  Then $n=3$ and $v_0=1$, so the printed right
side is $1/4$, whereas $u^3(0)=0$.  The same right side is correctly the
formal value
\[
 f\circ g\circ f(0)=\frac14,
 \qquad f(x)=\frac{x}{2},\quad g(x)=\frac{3x+1}{2}.
\]
Consequently this edition uses the formal-composition identity and identifies
it with $u^n(x)$ only after proving that the word is the actual parity
itinerary, as in \cref{prop:Tq-correct-fixed-point}.
\end{sourceissue}

The fixed point $R_{3,\boldsymbol\varepsilon}$ of that correct formal
composition is the candidate printed in B\"ohm--Sontacchi Proposition~5
\cite[p.~263]{BohmSontacchi}.  Their statement literally prints $v_m=n$;
after the explicit edition renaming $n=L$, it prints
$v_{i-1}<v_i$ for $1\leq i<m$ and $v_m=L$, but does not visibly print
$v_{m-1}<v_m$.  For $m\geq1$, the last inequality follows from the formal run
construction, where $v_m-v_{m-1}=1+n_m>0$; it is a reconstruction, not a
silent correction of Proposition~5.  The unnumbered paragraph immediately
after its proof explicitly says the candidate may lie on a cycle of length
$L$ or a proper divisor of $L$ \cite[p.~264]{BohmSontacchi}.  Thus the
independently recovered historical formula is a return-time identity, not an
exact-period classification.

\begin{corollary}[global nonidentity]
\label{cor:BL-numen-global-nonidentity}
For every odd prime $q$, the paired values with common rational representatives
in
\cref{thm:BL-numen-periodic-overlap} do not identify the global maps
$\Phi_q$ and $\chi_q$.  They have types
\[
 \Phi_q:\Z_2\to\Z_2,
 \qquad
 \chi_q:\Z_2\to\Z_q,
\]
The two values have distinct rational representatives:
\begin{equation}
\label{eq:BL-numen-value-separation}
 \Phi_q(1)=\iota_2\!\left(-\frac1q\right),
 \qquad
 \chi_q(1)=\iota_q\!\left(\frac12\right),
 \qquad
 -\frac1q\ne\frac12\quad\text{in }\Q.
\end{equation}
The map $Q_q$ must likewise retain its role as the inverse chronological
parity encoder in \cref{eq:BL-conjugacy-oriented}; it cannot be substituted
for $\Phi_q$ by a notation change.
\end{corollary}

\begin{proof}
Formula \cref{eq:BL-inverse-formula} at the single one position $d_1=0$
gives the first embedded value.  The one-letter numen is
$\chi_q(1)=H_1(0)=\iota_q(1/2)$.  Equality of the rational representatives
would force $q=-2$, contrary to oddness.  The displayed domains, codomains,
and inverse relation prove the remaining typing assertions.
\end{proof}

Laarhoven--de Weger give a graph-theoretic presentation of the same classical
conjugacy.  Their symbol $\Phi$ is the chronological graph isomorphism
satisfying $T=\Phi^{-1}\sigma\Phi$, so in the notation fixed here it plays
the role of $Q_3$, not $\Phi_3$.  Their odd-$an+b$ analogue is exact, whereas
their $p$-ary discussion uses ``appropriately chosen'' coefficients and does
not cover arbitrary Hydras
\cite[Thms.~3.1 and~5.1 and Sec.~5.2]{LaarhovenDeWeger}.

\subsection{Branch correctness for \texorpdfstring{$T_q$}{Tq}}

The shortened map admits a particularly clean correctness argument in
$\Q_2$.  Extend the two affine formulas in \cref{eq:Tq-branches} to all of
$\Q_2$.

\begin{lemma}[one correct branch]
\label{lem:one-correct-Tq-branch}
Let $q$ be odd and $x\in\Z_2$.  Exactly one of $H_0(x),H_1(x)$ belongs to
$\Z_2$, namely $H_{[x]_2}(x)$.  If a wrong branch is applied, the result has
$2$-adic valuation $-1$.  If $v_2(y)<0$, then
\[
 v_2(H_0(y))=v_2(H_1(y))=v_2(y)-1.
\]
\end{lemma}

\begin{proof}
If $x$ is even, $x/2\in\Z_2$ and $qx+1$ is odd, so
$(qx+1)/2\notin\Z_2$.  If $x$ is odd, $x/2\notin\Z_2$ and $qx+1$ is even,
so $(qx+1)/2\in\Z_2$.  This proves the first two assertions.  If
$v_2(y)<0$, then $v_2(qy)=v_2(y)<0=v_2(1)$; the ultrametric inequality with
unequal valuations gives $v_2(qy+1)=v_2(y)$.  Division by $2$ lowers either
valuation by one.
\end{proof}

\begin{proposition}[formal fixed points in $\Z_2$ are genuine]
\label{prop:Tq-correct-fixed-point}
Let $\mathbf j$ be a finite binary word.  If $z\in\Z_2$ and
$H_{\mathbf j}(z)=z$, then every branch in the right-to-left evaluation of
$H_{\mathbf j}(z)$ is correct.  Consequently
\[
 T_q^{\circ|\mathbf j|}(z)=z.
\]
\end{proposition}

\begin{proof}
Start with $z\in\Z_2$ and evaluate the rightmost branch first.  If any branch
is wrong, \cref{lem:one-correct-Tq-branch} says the intermediate value leaves
$\Z_2$ with negative valuation; every later branch decreases that valuation
further.  The final value then cannot equal $z\in\Z_2$.  Hence no branch was
wrong, and the formal composition is the genuine iterate with the reversed
chronological itinerary, as in \cref{eq:chronological-composition}.
\end{proof}

\subsection{The checked \texorpdfstring{$T_q$}{Tq}
periodic-point correspondence}

Assume now that $q$ is an odd prime, so that the numen
$\chi_q:\Z_2\to\Z_q$ is defined everywhere by
\cref{eq:chi-q-adic-series}.

\begin{theorem}[Siegel's periodic-word correspondence for $T_q$]
\label{thm:Tq-periodic-word-CP}
Let $n\geq1$, let $L=\lambda_2(n)$, and let $\mathbf j$ be the shortest
least-significant-digit-first binary word for $n$.  Then
\begin{equation}
\label{eq:chi-B2}
 \chi_q(B_2(n))
 =\frac{\chi_q(n)}{1-M_q(n)},
 \qquad
 M_q(n)=\frac{q^{\#_1(n)}}{2^L}.
\end{equation}
The right side is a rational number.  If it belongs to $\Z_2$, it is a
periodic point of the $2$-adic extension of $T_q$; if it belongs to $\Z$, it
is a periodic integer point of $T_q$.

Conversely, let $x$ be a nonzero periodic integer point of exact period $L$.
If $\mathbf j$ is the length-$L$ word obtained by reversing one complete
chronological itinerary of $x$ and $n=\Dig(\mathbf j)$, then
\begin{equation}
\label{eq:padded-cycle-code}
 x=\chi_q\bigl(B_{2,L}(n)\bigr),
 \qquad B_{2,L}(n)=\frac{n}{1-2^L}.
\end{equation}
At least one cyclic phase of the cycle has a reversed word whose final digit
is $1$; for that phase $L=\lambda_2(n)$ and
$B_{2,L}(n)=B_2(n)$.  Other phases may require the padded length-$L$ code.
\end{theorem}

\begin{proof}
The word contains at least one $1$, hence
$|M_q(n)|_q=q^{-\#_1(n)}<1$.  Apply
\cref{thm:repeated-code-fixed-point}; in the scalar case it gives exactly
\cref{eq:chi-B2}.  If the fixed point lies in $\Z_2$,
\cref{prop:Tq-correct-fixed-point} turns the formal fixed-point equation into
a genuine periodicity equation.

Conversely, let $x\ne0$ have period $L$ and let
$(\varepsilon_0,\ldots,\varepsilon_{L-1})$ be its chronological parity
itinerary.  Put
$\mathbf j=(\varepsilon_{L-1},\ldots,\varepsilon_0)$, so
$H_{\mathbf j}(x)=T_q^{\circ L}(x)=x$.  The word contains a $1$, since an
all-zero cycle would satisfy $x/2^L=x$ and hence $x=0$.  The affine fixed
point is unique because
$M_q(\mathbf j)=q^{\#_1(\mathbf j)}/2^L\ne1$: an odd power of the odd
integer $q$ cannot equal a positive power of $2$.  Therefore
$x=(1-M_q(\mathbf j))^{-1}\chi_q(\mathbf j)$, which is
$\chi_q(B_{2,L}(n))$ by \cref{thm:repeated-code-fixed-point}.  Rotating the
chronological starting point rotates the word.  One can choose a rotation
whose high digit is nonzero and thereby use the shortest-word abbreviation
$B_2(n)$, but that rotation generally changes $x$ to another point of its
cycle.
\end{proof}

\begin{sourceissue}[shortest words do not encode every cycle phase]
The stronger version in the 2007 source overclaims that every point in the
numen image of a cycle is represented by $B_2(n)$ for the shortest positive
word of $n$.  For $T_3$, the point $2$ on the cycle $1\leftrightarrow2$ is
encoded by the length-two word $(1,0)$; its high zero is part of the repeated
block, even though the ordinary integer represented by that word is $1$.
The synchronized code is $B_{2,2}(1)=1/(1-4)=-1/3$, not
$B_2(1)=1/(1-2)=-1$.  The padded formula
\cref{eq:padded-cycle-code} preserves the phase.  The source also writes a
truncation modulus $2^m$ for $m$ repeated length-$L$ blocks where the typed
modulus is $2^{mL}$; taking the correct subsequence repairs the limit.
\end{sourceissue}

\begin{sourceissue}[the centered zero orbit]
The 2026 Correspondence Principle is displayed for every periodic
$z\in\OK$ but requires a preimage in
$(\Z_p\setminus\N_0)\cap\Q$.  A centered Hydra has the periodic point $0$,
encoded by the all-zero word, whose digit-ring element is $0\in\N_0$.
No general argument in the source produces a non-eventually-zero rational
code with numen value $0$.  The earlier $T_q$ theorem correctly treats
nonzero periodic points.  The general statement therefore needs either an
explicit exception for $0$, inclusion of the code $0$, or an additional
surjectivity assertion about the zero fibre.
\end{sourceissue}

\subsection{The general 2026 Correspondence Principle}

We now record the latest source theorem without upgrading its proof status.
Let $H:\OK\to\OK$ be a proper, centered, exact-integral $\Lambda$-Hydra and
$p=|\OK/\Lambda|$.  Normalize a finite-place additive valuation by
$v(K^\times)=\Z$ and write
\[
 v(\Lambda)=\min\{v(\lambda):0\ne\lambda\in\Lambda\}.
\]
The source assumes:
\begin{enumerate}[label=\textup{(H\arabic*)}]
\item for each integer $n\geq1$ there is a place $\ell_n$ of $K$ such that
      $\|M_H(n)\|_{\ell_n}<1$;
\item there is a nontrivial nonarchimedean place $v$ such that
      $\|r_j\|_v>1$ for every branch and $v(\Lambda)\geq1$.
\end{enumerate}

\begin{status}[General CP (source)]
Under (H1)--(H2), arXiv:2601.17030v3 asserts:
\begin{enumerate}[label=\textup{(\roman*)}]
\item an integral point is periodic if and only if it is the value of $X_H$
      at a rational, non-eventually-zero base-$p$ digit input;
\item if $K$ is a number field and an irrational digit input has numen value
      $x\in\OK$, then
      $|H^{\circ m}(x)|_w\to\infty$ as $m\to\infty$ for every nontrivial
      archimedean absolute value $|\cdot|_w$.
\end{enumerate}
The source labels its proof ``Sketch.''  The repeated-word affine calculation
needed for the forward direction is proved in full in
\cref{thm:repeated-code-fixed-point}.  The converse still depends on the
source's completion-level ``correctness'' argument, and the divergence clause
depends on an identification of fibres not established by the displayed
hypotheses.  Accordingly, this edition records (i)--(ii) as source claims,
not as reconstructed theorems.
\end{status}

\begin{sourceissue}[defects inside the 2026 proof sketch]
The sketch contains five points which matter mathematically.
\begin{enumerate}[label=\textup{(\alph*)}]
\item It writes $M_q$ in a theorem about a general $H$; the typed symbol is
      $M_H$.
\item It writes a quotient $(1-M_H^m)/(1-M_H)$ although $M_H$ is generally
      an endomorphism.  The typed expression is
      $(I-M_H^m)(I-M_H)^{-1}$, and invertibility must be proved.
\item The proof of the divergence clause invokes a hypothesis ``(iii)'', but
      the theorem lists only (i) and (ii).
\item From the periodic-point correspondence, existence of a rational code
      for a periodic value does not imply that every code in the same numen
      fibre is rational.  Such an inference requires injectivity or a precise
      fibre theorem, neither of which is part of the displayed statement.
\item In dimension greater than one, $\|r_j\|_v>1$ is only a statement about
      the \emph{largest} expansion of $r_j$.  It does not imply
      $|r_jx|_v>|x|_v$ for every nonzero $x$, which is the kind of
      branch-detection estimate used by the sketch.  A sufficient typed
      replacement is a lower expansion bound
      $\inf_{x\ne0}|r_jx|_v/|x|_v>1$, equivalently
      $\|r_j^{-1}\|_v<1$ for invertible $r_j$, together with the required
      translation/integrality control.  In one scalar dimension the upper
      and lower factors coincide, but the general endomorphism statement
      cannot use the operator supremum alone.
\end{enumerate}
These are not stylistic complaints; they delimit what has and has not been
proved by the sketch.
\end{sourceissue}

\subsection{Weak Collatz versus full Collatz}

For $T_3$, the \emph{weak Collatz conjecture} in Siegel's contour paper is a
periodic-orbit assertion: the only integer cycles are the known ones (with
the sign and state-space convention stated in that paper).  The full Collatz
conjecture additionally requires every positive integer orbit to enter the
positive $1\leftrightarrow2$ cycle.  A characterization of periodic points,
even a complete one, does not rule out divergent positive trajectories.
The contour equivalence in \cref{sec:contour-mellin} is therefore kept under
the name \emph{weak}; it is not promoted to an equivalence with the full
conjecture.

\section{Set-series, dreamcatchers, and singularity conservation}
\label{sec:dreamcatchers}

Siegel's 2019 paper encodes invariant sets by holomorphic generating series
and then retains only their dominant radial boundary data.  The resulting
$\ell^2$ object is called a \emph{dreamcatcher}.  This is a different
analytic route from the numen: the numen codes branch compositions, whereas
the permutation operator below acts on coefficient functions and detects
sets invariant under inverse image.

\subsection{Set-series and the historical Hydra class}

For $V\subseteq\N_0$, define
\begin{align}
 \varsigma_V(w)&=\sum_{v\in V}w^v, && |w|<1,\label{eq:ordinary-set-series}\\
 \psi_V(z)&=\varsigma_V(e^{2\pi iz})
            =\sum_{v\in V}e^{2\pi ivz},&&\Im z>0.\label{eq:Fourier-set-series}
\end{align}
The coefficients are the indicator function $\ind{n\in V}$.  The 2019 paper
also uses exponential and Dirichlet set-series.  Its displayed Dirichlet
series $\sum_{v\in V}v^{-s}$ requires either $0\notin V$ or an explicit
treatment of the term $v=0$.

The Hydra used in that paper is the following one-dimensional historical
specialization.  It is not the controlling 2026 definition.

\begin{definition}[2019 $\varrho$-Hydra]
Fix $\varrho\geq2$.  A 2019 $\varrho$-Hydra is a surjection
$H:\N_0\to\N_0$ with
\begin{equation}
\label{eq:2019-hydra}
 H(n)=\frac{a_jn+b_j}{d_j}
 \quad(n\equiv j\pmod\varrho),
\end{equation}
where $a_j,d_j\geq1$, $b_j\geq0$, $\gcd(a_j,d_j)=1$, and each assigned
branch is integer-valued and nonnegative.  Set
\begin{equation}
\label{eq:mu-j}
 \mu_j=\frac{\varrho a_j}{d_j}\in\N.
\end{equation}
It is \emph{prime} when $\varrho$ is prime and every $a_j$ is $1$ or a
prime.  It is \emph{regulated} when $\mu_j=1$ for at least one $j$.
\end{definition}

The integrality of \cref{eq:mu-j} follows by comparing
$H(\varrho m+j)$ with $H(j)$, and coprimality then gives
$d_j\mid\varrho$.  Surjectivity is an extra hypothesis of the 2019 analytic
framework; it is not part of the 2026 word formalism.

\subsection{Permutation operators}

Let $\ell^\infty(\N_0)$ carry its supremum norm.  The coefficient pullback
is the linear map
\[
 \mathsf P_H:\ell^\infty(\N_0)\longrightarrow\ell^\infty(\N_0),
 \qquad (\mathsf P_Hc)_n=c_{H(n)}.
\]
Because the 2019 Hydra $H$ is surjective,
\[
 \|\mathsf P_Hc\|_\infty
 =\sup_{n\geq0}|c_{H(n)}|
 =\sup_{m\geq0}|c_m|=\|c\|_\infty;
\]
thus $\mathsf P_H$ is an isometric injection (not necessarily a surjection).
For $c\in\ell^\infty(\N_0)$, both
\[
 \mathcal G_{\mathbb D}c(w)=\sum_{n\geq0}c_nw^n
 \quad(|w|<1),
 \qquad
 \mathcal G_{\mathbb H}c(z)=\sum_{n\geq0}c_ne^{2\pi inz}
 \quad(\Im z>0)
\]
converge normally on compact subsets of their displayed domains.  The two
analytic operators are therefore the conjugates of this specified
coefficient map on the ranges of the injective generating maps:
\[
 \mathcal Q_H:\operatorname{ran}(\mathcal G_{\mathbb D})
 \longrightarrow\operatorname{ran}(\mathcal G_{\mathbb D}),
 \qquad
 \mathscr Q_H:\operatorname{ran}(\mathcal G_{\mathbb H})
 \longrightarrow\operatorname{ran}(\mathcal G_{\mathbb H}).
\]
\begin{align}
 \mathcal Q_H(\mathcal G_{\mathbb D}c)
   &=\mathcal G_{\mathbb D}(\mathsf P_Hc)
     =\sum_{n\geq0}c_{H(n)}w^n,\label{eq:ordinary-permutation}\\
 \mathscr Q_H(\mathcal G_{\mathbb H}c)
   &=\mathcal G_{\mathbb H}(\mathsf P_Hc)
     =\sum_{n\geq0}c_{H(n)}e^{2\pi inz}.
     \label{eq:Fourier-permutation}
\end{align}
For set-series this simply records inverse image:
\begin{equation}
\label{eq:Q-inverse-image}
 \mathcal Q_H\varsigma_V=\varsigma_{H^{-1}(V)},
 \qquad
 \mathscr Q_H\psi_V=\psi_{H^{-1}(V)}.
\end{equation}
Thus $H^{-1}(V)=V$ if and only if the corresponding set-series is a fixed
point.

Let $\xi_m=e^{2\pi i/m}$.  Splitting the input index modulo $\varrho$ and
using
\[
 H(\varrho n+j)=\mu_jn+H(j)
\]
gives, for every
$\psi\in\operatorname{ran}(\mathcal G_{\mathbb H})$, the finite branch
formula
\begin{equation}
\label{eq:Fourier-permutation-branch}
 (\mathscr Q_H\psi)(z)
 =\sum_{j=0}^{\varrho-1}\frac1{\mu_j}
   \sum_{k=0}^{\mu_j-1}\xi_{\mu_j}^{-kH(j)}
   e^{-2\pi i(b_j/a_j)z}
   \psi\left(\frac{\varrho z+k}{\mu_j}\right).
\end{equation}
To verify the formula, expand the right side.  The inner root-of-unity sum
selects coefficients whose indices are congruent to $H(j)$ modulo $\mu_j$;
the exponential prefactor converts the surviving exponent back to
$\varrho n+j$.

\begin{proposition}[fixed coefficient functions]
\label{prop:permutation-fixed-coefficients}
For a coefficient function $c:\N_0\to\C$, define its coefficient pullback by
\[
 (\mathsf P_Hc)_n=c_{H(n)}.
\]
Then $\mathsf P_Hc=c$ if and only if
\begin{equation}
\label{eq:coeff-orbit-constant}
 c_n=c_{H(n)}\qquad(n\geq0).
\end{equation}
Equivalently, $c$ is constant on every total orbit class of $H$.
\end{proposition}

\begin{proof}
By definition, $\mathsf P_Hc=c$ is exactly
\cref{eq:coeff-orbit-constant}.
Iterating that identity makes $c$ constant along forward orbits; if two
points have a common forward image, their coefficients agree.  Conversely,
constancy on orbit classes implies \cref{eq:coeff-orbit-constant}.
\end{proof}

\begin{sourceissue}[``basis'' in the Permutation Operator Theorem]
The source says that the set-series of all orbit classes form a basis of the
fixed-point space.  \Cref{prop:permutation-fixed-coefficients} gives the exact
coefficient-level statement.  If infinitely many orbit classes have
nonzero coefficients, the resulting decomposition is an infinite sum, not
a finite Hamel-basis expansion.  A topological basis claim requires a named
function-space topology and a convergence proof.
\end{sourceissue}

\subsection{Edgepoints and normalization}

Let $\mathbb H_{+i}=\{z\in\C:\Im z>0\}$ and
$\psi\in\mathcal A(\mathbb H_{+i})$.

\begin{definition}[edgepoint and virtual residue]
An $x\in\R$ is an \emph{edgepoint} of $\psi$ when the finite limit
\[
 \operatorname{VRes}_x[\psi]
 =\lim_{y\downarrow0}y\psi(x+iy)
\]
exists.  It is a \emph{virtual pole} when that limit is nonzero.  Write
\[
 P(\psi)=\{x\in\R:x\text{ is an edgepoint of }\psi\}.
\]
\end{definition}

For a Fourier set-series, the elementary estimate
\[
 |\psi_V(x+iy)|\leq\sum_{n\geq0}e^{-2\pi ny}
 =\frac1{1-e^{-2\pi y}}
\]
implies
\begin{equation}
\label{eq:no-super-simple-growth}
 \lim_{y\downarrow0}y^{1+\delta}\psi_V(x+iy)=0
 \qquad(x\in\R,\ \delta>0).
\end{equation}
Thus the virtual residue records the largest possible radial growth order
for a digital Fourier series.

\begin{sourceissue}[ordinary versus virtual residues]
If $\psi$ has an ordinary simple pole at $x$ with complex residue $A$, then
$y\psi(x+iy)\to A/i=-iA$.  The source's definition of a ``classical virtual
pole'' subtracts $\operatorname{VRes}_x[\psi]/(z-x)$ and thereby omits this
factor of $i$.  This edition retains the term \emph{virtual residue} for the
radial normalization above but does not identify it numerically with the
ordinary complex residue.
\end{sourceissue}

\begin{proposition}[cross-section densities give the virtual residue]
\label{prop:cross-section-vres}
Let $\alpha\geq1$, $\beta\in\Z$, and assume that every limit
\[
 d_k=\lim_{N\to\infty}\frac1N
 |V\cap(k+\alpha\Z)\cap[0,N)|
 \qquad(0\leq k<\alpha)
\]
exists.  Then $\beta/\alpha$ is an edgepoint of $\psi_V$ and
\begin{equation}
\label{eq:VRes-cross-section-density}
 \operatorname{VRes}_{\beta/\alpha}[\psi_V]
 =\frac1{2\pi}\sum_{k=0}^{\alpha-1}
 d_k\xi_\alpha^{\beta k}.
\end{equation}
In particular, if $V$ has natural density $d(V)$, then
$\operatorname{VRes}_0[\psi_V]=d(V)/(2\pi)$.
\end{proposition}

\begin{proof}
Put $a_n=\ind{n\in V}\xi_\alpha^{\beta n}$.  Splitting a Ces\`aro sum by
residue class gives
\[
 \lim_{N\to\infty}\frac1N\sum_{n=0}^{N-1}a_n
 =\sum_{k=0}^{\alpha-1}d_k\xi_\alpha^{\beta k}=:L.
\]
The elementary Abel implication for a bounded sequence now gives
\[
 (1-r)\sum_{n\geq0}a_nr^n\longrightarrow L
 \qquad(r\uparrow1).
\]
For completeness, if $A_N=\sum_{n<N}(a_n-L)$, then $A_N=o(N)$ and summation
by parts gives
$(1-r)\sum_{n\geq0}(a_n-L)r^n=(1-r)^2\sum_{N\geq1}A_Nr^{N-1}\to0$;
split the last sum at a fixed index after bounding $|A_N|/N$ on its tail.
Finally set $r=e^{-2\pi y}$.  Since
$y/(1-e^{-2\pi y})\to1/(2\pi)$,
\[
 y\psi_V(\beta/\alpha+iy)
 =\frac{y}{1-r}(1-r)\sum_{n\geq0}a_nr^n
 \longrightarrow\frac{L}{2\pi},
\]
which proves the assertion.  Only the Abel implication is used; no Tauberian
converse is invoked.
\end{proof}

Dreamcatchers therefore contain residue-class density data, not merely a
list of analytic singularities.

\subsection{The Square Sum Lemma}

Equip subsets of $\R/\Z$ with counting measure.  The result below is stated
in Siegel's paper, which credits the proof idea to accepted MathOverflow
answer 333801 by user \emph{fedja} \cite{FedjaMO333801}.  The argument below
is instead this edition's finite-subset Bessel proof; it is not attributed to
the answerer.  It avoids the source proof's interchange of infinite sums
before summability has been established.

\begin{lemma}[Square Sum Lemma]
\label{lem:square-sum}
Let $V\subseteq\N_0$ and put $P'=P(\psi_V)\cap[0,1)$.  Then
\begin{equation}
\label{eq:square-sum}
 \sum_{t\in P'}|\operatorname{VRes}_t[\psi_V]|^2
 \leq\frac1{(2\pi)^2}.
\end{equation}
In particular, the virtual-residue function belongs to $\ell^2(P')$ and
has at most countable nonzero support.
\end{lemma}

\begin{proof}
Let $a_n=\ind{n\in V}$, let $r=e^{-2\pi y}$, and put on bounded sequences
the semidefinite inner product
\[
 \langle u,v\rangle_y=(1-r)\sum_{n\geq0}u_n\overline{v_n}r^n.
\]
For $t\in\R/\Z$, set $e_t(n)=e^{-2\pi int}$.  Then
\begin{align*}
 \langle e_s,e_t\rangle_y
 &=\frac{1-r}{1-re^{2\pi i(t-s)}}
   \longrightarrow\ind{s=t},\\
 \langle a,e_t\rangle_y
 &=(1-r)\psi_V(t+iy)
   \longrightarrow2\pi\operatorname{VRes}_t[\psi_V],\\
 \|a\|_y^2&=(1-r)\sum_{n\geq0}a_nr^n\leq1.
\end{align*}
Fix a finite $F\subseteq P'$.  The Gram matrix
$G_y=(\langle e_s,e_t\rangle_y)_{s,t\in F}$ tends to the identity and is
invertible for small $y$.  The finite-dimensional projection inequality is
\[
 b_y^*G_y^{-1}b_y\leq\|a\|_y^2,
 \qquad b_y=(\langle a,e_t\rangle_y)_{t\in F}.
\]
Passing to the limit yields
$(2\pi)^2\sum_{t\in F}|\operatorname{VRes}_t[\psi_V]|^2\leq1$.
Taking the supremum over finite $F$ gives \cref{eq:square-sum}.
\end{proof}

\subsection{Dreamcatchers and branch-invariant carriers}

\begin{definition}[dreamcatcher over $T$]
Let $\psi\in\mathcal A(\mathbb H_{+i})$, and let
$T\subseteq P(\psi)$ be $1$-periodic.  The function $\psi$ \emph{admits a
dreamcatcher over $T$} when
\begin{equation}
\label{eq:dreamcatcher}
 R_{\psi,T}(x)=\operatorname{VRes}_x[\psi]
 \quad(x\in T)
\end{equation}
belongs to $\ell^2(T/\Z)$.  Following the source, denote the class of such
functions by $\mathcal D_1(T)$.
\end{definition}

The carrier $T$ is \emph{$H$-branch invariant} when
\begin{equation}
\label{eq:H-branch-invariant}
 \frac{\varrho T+k}{\mu_j}\subseteq T
 \quad(0\leq j<\varrho,\ 0\leq k<\mu_j).
\end{equation}
This ensures that every virtual residue appearing after applying a branch in
\cref{eq:Fourier-permutation-branch} is defined on the same carrier.

\begin{lemma}[Virtual Residue Formula]
\label{lem:VRF}
Let $T$ be $H$-branch invariant and let
$\psi\in\mathcal D_1(T)\cap\operatorname{ran}(\mathcal G_{\mathbb H})$.
At every $t\in T$, the function $\mathscr Q_H\psi$ has an edgepoint and
\begin{equation}
\label{eq:VRF}
 \operatorname{VRes}_t[\mathscr Q_H\psi]
 =\frac1\varrho\sum_{j=0}^{\varrho-1}
 e^{-2\pi i(b_j/a_j)t}
 \sum_{k=0}^{\mu_j-1}\xi_{\mu_j}^{-kH(j)}
 R_{\psi,T}\left(\frac{\varrho t+k}{\mu_j}\right).
\end{equation}
\end{lemma}

\begin{proof}
Multiply \cref{eq:Fourier-permutation-branch} by $y$ and substitute
$z=t+iy$.  The exponential prefactor converges to its value at $t$.  For
$A_{j,k}(z)=(\varrho z+k)/\mu_j$, the imaginary part of $A_{j,k}(t+iy)$ is
$\varrho y/\mu_j$, hence
\[
 y\psi(A_{j,k}(t+iy))
 =\frac{\mu_j}{\varrho}
   \left(\frac{\varrho y}{\mu_j}\right)
   \psi\left(A_{j,k}(t)+i\frac{\varrho y}{\mu_j}\right).
\]
The factor $1/\mu_j$ in \cref{eq:Fourier-permutation-branch} cancels the
$\mu_j$ here, leaving $1/\varrho$.  The sum is finite, so its limit is the
sum of the limits.
\end{proof}

\begin{definition}[dreamcatcher operator]
For $R\in\ell^2(T/\Z)$ let $\mathcal E_{H:T}R:T\to\C$ denote the raw finite
expression
\begin{equation}
\label{eq:dreamcatcher-operator}
 (\mathcal E_{H:T}R)(x)
 =\frac1\varrho\sum_{j=0}^{\varrho-1}
 e^{-2\pi i(b_j/a_j)x}
 \sum_{k=0}^{\mu_j-1}\xi_{\mu_j}^{-kH(j)}
 R\left(\frac{\varrho x+k}{\mu_j}\right).
\end{equation}
Here $R(y)$ means the value on the class $y+\Z$; branch invariance ensures
that every argument belongs to $T$.  Define the maximal graph domain
\[
 \operatorname{Dom}(\mathfrak Q_{H:T})
 =\{R\in\ell^2(T/\Z):
     \mathcal E_{H:T}R\text{ is $1$-periodic on $T$ and }
     [\mathcal E_{H:T}R]\in\ell^2(T/\Z)\}.
\]
The dreamcatcher operator is the linear operator
\[
 \mathfrak Q_{H:T}:\operatorname{Dom}(\mathfrak Q_{H:T})
 \subseteq\ell^2(T/\Z)\longrightarrow\ell^2(T/\Z),
 \qquad R\longmapsto[\mathcal E_{H:T}R].
\]
\end{definition}

\begin{theorem}[Singularity Conservation Law]
\label{thm:SCL}
Let $T$ be $H$-branch invariant, let
$\psi\in\mathcal D_1(T)\cap\operatorname{ran}(\mathcal G_{\mathbb H})$ be
fixed by $\mathscr Q_H$, and put $R=R_{\psi,T}$.  Suppose
$R\in\operatorname{Dom}(\mathfrak Q_{H:T})$.  Then
\begin{equation}
\label{eq:SCL}
 \mathfrak Q_{H:T}R=R.
\end{equation}
\end{theorem}

\begin{proof}
Since $\mathscr Q_H\psi=\psi$, the left side of
\cref{eq:VRF} is $\operatorname{VRes}_t[\psi]=R(t)$.  Its right side is
$\mathcal E_{H:T}R(t)$, whose descended class is
$\mathfrak Q_{H:T}R$.  Equality holds pointwise on $T$ and hence in
$\ell^2(T/\Z)$.
\end{proof}

\begin{sourceissue}[undefined class and operator domain]
The Virtual Residue Formula and SCL in the source assume
$\psi\in\mathcal D_2(T)$, but only $\mathcal D_1(T)$ is defined.  The intended
symbol appears to be $\mathcal D_1(T)$.  More seriously, the proof that
\cref{eq:dreamcatcher-operator} preserves $\ell^2(T/\Z)$ treats the inclusion
in \cref{eq:H-branch-invariant} as saying that each affine branch map is a
bijection of $T$.  A bijection of $\R$ whose image sends $T$ into $T$ need
not map $T$ onto itself.  The typed SCL above therefore uses the exact
maximal graph domain and makes membership an explicit hypothesis for the
particular dreamcatcher.  It does not assert that the graph domain is all of
$\ell^2(T/\Z)$ or that the operator is bounded.  Either conclusion needs
two-sided branch invariance or a separate finite-multiplicity norm estimate.
\end{sourceissue}

A function $\psi$ is \emph{tame} in the source when every rational number is
an edgepoint.  In that case $T=\Q$ is available and, by
\cref{lem:square-sum}, the rational virtual residues of a Fourier set-series
form an element of $\ell^2(\Q/\Z)$.

\subsection{The Finite Dreamcatcher Theorem}

The later half of the 2019 paper studies
$\ker(I-\mathfrak Q_H)$ purely on $\ell^2(\Q/\Z)$.  Its central statement is
the following.

\begin{status}[Finite Dreamcatcher Theorem, source theorem]
\label{status:FDT}
Let $H$ be a prime, regulated 2019 Hydra.  If
\[
 R\in\ker(I-\mathfrak Q_{H:\Q})\subseteq\ell^2(\Q/\Z)
\]
and $\supp R$ is finite modulo $\Z$, then
\begin{equation}
\label{eq:FDT}
 R(x)=R(0)\ind{x\in\Z}
 \qquad(x\in\Q).
\end{equation}
The source proof proceeds through the Off-$\varrho$ Theorem, an auxiliary
functional equation, two support propositions, and the Off-$\varrho$
Infinity Theorem.  This edition records the result with that proof status;
it does not use the earlier flawed $\ell^2(T)$-invariance argument as an
independent proof of the finite-support theorem.
\end{status}

The result says that a finitely supported conserved singularity pattern has
no nonintegral rational frequency.  Combining it with the SCL gives the
paper's first weak rationality consequence.

\begin{status}[Weak Rationality Theorem 1, source consequence]
If $V\subseteq\N_0$ is rational in the sense that its ordinary generating
function is rational, and if $H^{-1}(V)=V$ for a prime regulated Hydra, then
$V$ is finite or cofinite in $\N_0$.
\end{status}

Indeed, rationality makes the virtual-pole support finite modulo $\Z$;
invariance gives a permutation-operator fixed point; SCL and FDT force the
only possible dominant rational boundary frequency to be $0$.  The remaining
step uses the eventual periodicity of a $0$--$1$ sequence with rational
generating function.

\begin{sourceissue}[results not promoted from the 2019 paper]
Three later claims are kept out of the proved chain.
\begin{enumerate}[label=\textup{(\roman*)}]
\item The proof of Weak Rationality Theorem 2 asserts that naturally
      measurable subsets of $\N_0$ form a $\sigma$-algebra and that natural
      density is countably additive.  Both assertions are false in general,
      so that theorem is not canonized as proved.
\item Theorem 12 assumes branch invariance of both the rational edgepoints
      and their rational complement.  The source itself calls this double
      requirement the method's ``Achilles heel''; the edition retains the
      result only with that hypothesis if cited.
\item The assertion that a limit of dreamcatchers is the dreamcatcher of the
      limiting function is explicitly Conjecture 1.  A nearby corollary also
      uses an asymptotic fixed-point premise not present in its statement.
\end{enumerate}
\end{sourceissue}

\section{Dirichlet series, Perron inversion, and the contour programme}
\label{sec:contour-mellin}

The paper designated Part I.5 moves the fixed-point numerator
\[
 (1-M_q(n))x-\chi_q(n)
\]
into the coefficients of a Dirichlet series.  Perron inversion can then
recover that numerator by a vertical contour integral.  The coefficient
extraction is sound; the advertised equivalence with every periodic point
inherits the shortest-word defect isolated in
\cref{eq:padded-cycle-code}.

\subsection{The three Dirichlet series}

Let $q$ be an odd prime and $x\in\C$.  With $M_q(n)$ defined using the
shortest binary word of $n\geq1$, put
\begin{align}
 \zeta_{M_q}(s)&=\sum_{n=1}^{\infty}\frac{M_q(n)}{n^s},
 \label{eq:zeta-Mq}\\
 \zeta_{\chi_q}(s)&=\sum_{n=1}^{\infty}\frac{\chi_q(n)}{n^s},
 \label{eq:zeta-chiq}\\
 F_q(s,x)&=\sum_{n=1}^{\infty}
 \frac{(1-M_q(n))x-\chi_q(n)}{n^s}\notag\\
 &=\bigl(\zeta(s)-\zeta_{M_q}(s)\bigr)x-
   \zeta_{\chi_q}(s).
 \label{eq:Fq}
\end{align}
The source writes
\begin{equation}
\label{eq:sigma-q}
 \sigma_q=\log_2\left(\frac{q+1}{2}\right)
\end{equation}
for the common absolute-convergence threshold obtained from its summatory
estimates.  All contour identities below are first asserted on a line
$\Re s=c>\sigma_q$.

\subsection{Dyadic summation}

The digit recursions imply exact block sums.  Write
\[
 A_N=\sum_{0\leq n<2^N}\chi_q(n),
 \qquad
 B_N=\sum_{0\leq n<2^N}M_q(n),
\]
with $M_q(0)=1$.  Separating even and odd indices gives
\begin{align}
 A_{N+1}&=\frac{q+1}{2}A_N+2^{N-1},\label{eq:A-recurrence}\\
 B_{N+1}&=\frac{q+1}{2}B_N+\frac12.
 \label{eq:B-recurrence}
\end{align}
Solving these recurrences yields
\begin{align}
 A_N&=
 \begin{cases}
 N2^N/4,&q=3,\\[2mm]
 \displaystyle\frac{2^N}{q-3}
 \left[\left(\frac{q+1}{4}\right)^N-1\right],&q\ne3,
 \end{cases}
 \label{eq:chi-block-sum}\\
 B_N&=\frac{q\left(\frac{q+1}{2}\right)^N-1}{q-1}.
 \label{eq:M-block-sum}
\end{align}

\begin{proof}
For \cref{eq:A-recurrence}, use
\[
 \chi_q(2m)+\chi_q(2m+1)
 =\frac{q+1}{2}\chi_q(m)+\frac12
\]
and sum over $0\leq m<2^N$.  For \cref{eq:B-recurrence}, the usual
$M_q(2m)=M_q(m)/2$ fails only at $m=0$, because the shortest-word convention
sets $M_q(0)=1$.  Hence the even sum is
$1+\frac12(B_N-1)$ and the odd sum is $qB_N/2$.  The initial values are
$A_0=0$ and $B_0=1$; elementary solution of the recurrences gives the two
closed forms.
\end{proof}

These identities explain the appearance of $\sigma_q$.  A guaranteed
absolute-convergence half-plane for a linear combination is the intersection
of the component half-planes and therefore uses the maximum of their
abscissae, not the source's printed ``minimum.''  The exact abscissa of a
particular linear combination can be smaller when leading terms cancel, so
no equality of abscissae is inferred for arbitrary complex $x$.

\subsection{Perron extraction of one coefficient}

Let
\[
 a_n(x)=(1-M_q(n))x-\chi_q(n).
\]
For $0<a<1$, $n\geq3$, and $c>\sigma_q$, the Perron difference kernel
isolates $a_n(x)$:
\begin{theorem}[coefficient contour identity]
\label{thm:coefficient-contour}
Under the preceding hypotheses,
\begin{equation}
\label{eq:coefficient-contour}
 a_n(x)=\operatorname{PV}\left\{
 \frac1{2\pi i}\int_{c-i\infty}^{c+i\infty}
 \frac{(n+a)^s-(n+a-1)^s}{s}F_q(s,x)\,ds
 \right\}.
\end{equation}
\end{theorem}

\begin{proof}
For a Dirichlet series $F(s)=\sum_{m\geq1}a_mm^{-s}$ absolutely convergent
on $\Re s=c$, Perron's formula gives
\[
 \operatorname{PV}\frac1{2\pi i}
 \int_{c-i\infty}^{c+i\infty}F(s)X^s\frac{ds}{s}
 =\sum_{m<X}a_m
\]
when $X$ is not an integer.  Apply this with $X=n+a$ and with
$X=n+a-1$.  Since $0<a<1$, the two index sets differ only by $m=n$.
Subtracting gives \cref{eq:coefficient-contour}.
\end{proof}

Consequently the integral vanishes exactly when
\begin{equation}
\label{eq:contour-fixed-numerator}
 (1-M_q(n))x=\chi_q(n).
\end{equation}
This is a coefficient theorem.  Its dynamical content is precisely whatever
fixed-point correspondence is valid for the shortest word representing
$n$.

\begin{status}[the printed contour ``weak Collatz'' theorem]
The source states that periodicity is equivalent to the vanishing of
\cref{eq:coefficient-contour} for some $n\geq3$, with separate exclusions:
\begin{enumerate}[label=\textup{(\roman*)}]
\item for $q=3$, $x\in\Z\setminus\{-1,0,1\}$;
\item for $q=5$, $x\in\Z\setminus\{-1,0\}$;
\item for $q\geq7$, apparently $x\in\Z\setminus\{0\}$, although the
      printed set is malformed.
\end{enumerate}
Each part needs its own contour line $c>\sigma_3$, $c>\sigma_5$, or
$c>\sigma_q$; the outer quantifiers in the printed theorem do not state
this correctly.
\end{status}

\begin{sourceissue}[the contour equivalence misses padded cycle words]
The claimed biconditional is not valid for every listed periodic phase.
For example, $2$ is a periodic point of $T_3$ and is not among the source's
exclusions.  If a shortest word fixed $2$, branch correctness would force it
to be a whole number of repetitions of the reversed itinerary $(1,0)$.
Every such word has a high terminal zero and therefore is not shortest.
Equivalently, $2=\chi_3(B_{2,2}(1))$ uses the padded length-two block, while
no shortest positive block yields \cref{eq:contour-fixed-numerator} with
$x=2$.  Thus \cref{thm:coefficient-contour} is canonical, but the full
periodicity equivalence is retained only as a source claim.  A repaired
phase-sensitive contour formulation must index the block length separately
from its ordinary digit sum.
\end{sourceissue}

\subsection{Corrected functional equations for \texorpdfstring{$f_q$ and
$F_q$}{fq and Fq}}

Extend the coefficient sequence by $a_0(x)=0$, and set
\begin{equation}
 \label{eq:Fq-C}
 C_q(x)=\frac{(q-2)x+1}{2}.
\end{equation}
The digit recursions determine both the ordinary and Dirichlet generating
functions without an analytic-continuation premise.

\begin{proposition}[coefficient and ordinary generating-function recursions]
\label{prop:Fq-coefficient-recursion}
For $n\geq1$ and $n\geq0$, respectively,
\begin{equation}
 \label{eq:Fq-coefficient-recursion}
 a_{2n}(x)=\frac{a_n(x)+x}{2},
 \qquad
 a_{2n+1}(x)=\frac q2a_n(x)-C_q(x).
\end{equation}
Consequently, in the formal power-series ring $\C[[z]]$, and analytically
for $|z|<1$,
\begin{equation}
 \label{eq:fq-corrected-functional-equation}
 f_q(z,x):=\sum_{n=1}^{\infty}a_n(x)z^n
 =\frac{qz+1}{2}f_q(z^2,x)
  +\frac{xz^2-((q-2)x+1)z}{2(1-z^2)}.
\end{equation}
\end{proposition}

\begin{proof}
For $n\geq1$, the shortest-word convention and
\cref{eq:chi-even,eq:chi-odd} give
\[
 M_q(2n)=\frac12M_q(n),\qquad
 M_q(2n+1)=\frac q2M_q(n),
\]
and the odd identities remain valid at $n=0$ because $M_q(0)=1$ and
$\chi_q(0)=a_0(x)=0$.  Substitution in
$a_n(x)=(1-M_q(n))x-\chi_q(n)$ proves
\cref{eq:Fq-coefficient-recursion}.  Splitting the formal series into its
positive even indices and its nonnegative odd indices gives
\begin{align*}
 \sum_{n\geq1}a_{2n}(x)z^{2n}
 &=\frac12f_q(z^2,x)+\frac{xz^2}{2(1-z^2)},\\
 \sum_{n\geq0}a_{2n+1}(x)z^{2n+1}
 &=\frac{qz}{2}f_q(z^2,x)-\frac{C_q(x)z}{1-z^2}.
\end{align*}
Their sum is \cref{eq:fq-corrected-functional-equation}.  If the shortest
word of $n$ has length $L$, then
$M_q(n)\leq(q/2)^L$ and
$\chi_q(n)\leq\frac12\sum_{j=0}^{L-1}(q/2)^j$.
Thus $a_n(x)$ has polynomial growth in $n$, so the formal identity
converges absolutely for $|z|<1$.
\end{proof}

\begin{theorem}[corrected Dirichlet functional equation]
\label{thm:Fq-corrected-functional-equation}
Let $q$ be an odd prime, $x\in\C$, and
$\Re s>\sigma_q$.  Then every series below is absolutely convergent and
\begin{equation}
\label{eq:Fq-functional-equation}
\begin{split}
 F_q(s,x)={}&
 \frac{((q-1)x+1)-2^s((q-2)x+1)}
      {2(2^s-2^{\sigma_q})}\,\zeta(s)\\
 &+\frac{q/2}{2^s-2^{\sigma_q}}
 \sum_{m=1}^{\infty}\left(-\frac12\right)^m
 \binom{s+m-1}{m}F_q(s+m,x).
\end{split}
\end{equation}
The series over $m$ converges locally uniformly in that half-plane.
\end{theorem}

\begin{proof}
Put $\alpha=(q+1)/2=2^{\sigma_q}$.  Positivity of $M_q(n)$ and
$\chi_q(n)$, together with \cref{eq:chi-block-sum,eq:M-block-sum}, gives
\[
 \sum_{1\leq n<2^N}|a_n(x)|
 \leq |x|\bigl(2^N+B_N\bigr)+A_N.
\]
For $q>3$ the right side is $O_{q,x}(\alpha^N)$; for $q=3$ it is
$O_x(N2^N)$.  Dyadic subdivision therefore proves absolute convergence of
$F_q(s,x)$ for $\Re s>\sigma_q$.

Split the Dirichlet series by parity.  The even recursion gives
\[
 \sum_{n\geq1}\frac{a_{2n}(x)}{(2n)^s}
 =2^{-s-1}\bigl(F_q(s,x)+x\zeta(s)\bigr).
\]
For the odd part, the binomial series at $1/(2n)$ gives
\begin{equation}
 \label{eq:Fq-odd-shift-expansion}
 \sum_{n\geq1}\frac{a_n(x)}{(2n+1)^s}
 =2^{-s}\sum_{m=0}^{\infty}\left(-\frac12\right)^m
   \binom{s+m-1}{m}F_q(s+m,x).
\end{equation}
To justify the interchange exactly, let $K$ be a compact subset of
$\Re s>\sigma_q$, put
$u=\inf_{s\in K}\Re s$, and choose
$R\geq\max(1,\sup_{s\in K}|s|)$.  From the product formula,
\begin{align*}
 \left|\binom{s+m-1}{m}\right|
 &\leq R\prod_{j=1}^{m-1}\left(1+\frac{R-1}{j+1}\right)\\
 &\leq Re^{R-1}(m+1)^{R-1}.
\end{align*}
Consequently the absolute value of the double series is bounded uniformly
on $K$ by
\[
 Re^{R-1}\left(\sum_{m\geq0}2^{-m}(m+1)^{R-1}\right)
 \left(\sum_{n\geq1}|a_n(x)|n^{-u}\right)<\infty.
\]
This proves both interchanges and local uniformity.  Since
\[
 \sum_{n\geq0}(2n+1)^{-s}=(1-2^{-s})\zeta(s)
\]
the odd recursion in \cref{eq:Fq-coefficient-recursion} and
\cref{eq:Fq-odd-shift-expansion} yield
\begin{align*}
 F_q(s,x)={}&2^{-s-1}\bigl(F_q(s,x)+x\zeta(s)\bigr)\\
 &+\frac q2\,2^{-s}\sum_{m=0}^{\infty}
 \left(-\frac12\right)^m\binom{s+m-1}{m}F_q(s+m,x)\\
 &-C_q(x)(1-2^{-s})\zeta(s).
\end{align*}
Move the $m=0$ and even copies of $F_q(s,x)$ to the left, multiply by
$2^s$, and use
$x/2+C_q(x)=((q-1)x+1)/2$.  This is
\cref{eq:Fq-functional-equation}.
\end{proof}

\begin{sourceissue}[the printed functional equation is false for
\texorpdfstring{$x\ne0$}{x not equal to zero}]
The source substitutes $z/(1-z)$ where the first even correction in
$f_q(z^2,x)$ is $z^2/(1-z^2)$.  Its resulting Dirichlet equation has
zeta coefficient
\[
 -\frac{2^s((q-1)x+1)-1}{2(2^s-2^{\sigma_q})}.
\]
As real $s\to+\infty$, the left side tends to
$a_1(x)=-((q-2)x+1)/2$.  The shifted series is bounded in this limit:
the binomial identity and the estimate
$F_q(s+m,x)=a_1(x)+O(2^{-s-m})$ give
\[
 \sum_{m\geq1}\left(-\frac12\right)^m
 \binom{s+m-1}{m}F_q(s+m,x)
 =a_1(x)\bigl((2/3)^s-1\bigr)+O((2/3)^s).
\]
The printed right side therefore tends to $-((q-1)x+1)/2$.  Thus the
printed identity fails for every $x\ne0$;
\cref{eq:Fq-functional-equation} is the corrected identity.  No
continuation or singularity claim is inherited from the false equation.
A fresh continuation proof from the corrected equation is given below.
The denominator lattice alone still does not establish an actual pole,
because the numerator can vanish.
The coefficient contour identity is unaffected: it uses only the original
absolutely convergent Dirichlet series on $\Re s>\sigma_q$.
\end{sourceissue}

\subsection{Meromorphic continuation from the corrected equation}

Part I.5 uses its shift equation to continue $F_q$, and the public Web II
exposition also reports a meromorphic continuation with a leftward
half-lattice of singularities
\cite{SiegelPartIHalf,SiegelWebII}.  Because the printed equation is false,
the conclusion must be rebuilt from \cref{eq:Fq-functional-equation}.
Put
\begin{align}
 b_m(s)&=\left(-\frac12\right)^m\binom{s+m-1}{m},\label{eq:Fq-bm}\\
 A_q(s,x)&=\frac{((q-1)x+1)-2^s((q-2)x+1)}2,\label{eq:Fq-Aqx}\\
 S_q(s,x)&=\sum_{m=1}^{\infty}b_m(s)F_q(s+m,x).
 \label{eq:Fq-Sqx}
\end{align}

\begin{lemma}[normal convergence of the shifted tail]
\label{lem:Fq-normal-shifted-tail}
For every compact $K\subset\C$, there are constants $C_K,A_K>0$ such
that
\[
 |b_m(s)|\leq C_K2^{-m}(m+1)^{A_K}
 \qquad(s\in K,\ m\geq1).
\]
If $m$ is large enough that $s+m$ lies in the original
absolute-convergence half-plane for every $s\in K$, then
\[
 F_q(s+m,x)=a_1(x)+O_K(2^{-m}).
\]
Hence the tail of \cref{eq:Fq-Sqx} converges normally on $K$.
\end{lemma}

\begin{proof}
Take $R\geq\max(1,\sup_{s\in K}|s|)$.  The product estimate used in the
proof of \cref{thm:Fq-corrected-functional-equation} gives
\[
 \left|\binom{s+m-1}{m}\right|
 \leq Re^{R-1}(m+1)^{R-1},
\]
which proves the first assertion.  Choose $r>\sigma_q$ and put
$D_r=\sum_{n\geq1}|a_n(x)|n^{-r}<\infty$.  If
$u=\inf_{s\in K}\Re s$ and $m+u\geq r$, then
\begin{align*}
 |F_q(s+m,x)-a_1(x)|
 &\leq\sum_{n\geq2}|a_n(x)|n^{-\Re s-m}\\
 &\leq D_r2^{r-u-m}.
\end{align*}
Multiplication by the first bound leaves a convergent exponential--
polynomial majorant, uniformly on $K$.
\end{proof}

\begin{theorem}[global meromorphic continuation]
\label{thm:Fq-global-meromorphic-continuation}
For every odd prime $q$ and every $x\in\C$, the Dirichlet series
$F_q(\,\cdot\,,x)$ has a unique meromorphic continuation to $\C$.
If
\[
 \Lambda_q=\left\{\sigma_q+\frac{2\pi i k}{\log 2}:k\in\Z\right\},
\]
then its poles are contained in the explicitly generated candidate set
\begin{equation}
 \label{eq:Fq-candidate-poles}
 \operatorname{Poles}F_q(\,\cdot\,,x)\subseteq\mathcal P_q,
 \qquad
 \mathcal P_q=
 \{\lambda-n:\lambda\in\Lambda_q,\ n\in\N_0\}
 \cup\{1-n:n\in\N_0\}.
\end{equation}
The inclusion need not be equality: cancellation can make a candidate
singularity removable.
\end{theorem}

\begin{proof}
For $r\in\N_0$, set
\[
 H_r=\{s\in\C:\Re s>\sigma_q-r\}.
\]
The defining series is holomorphic on $H_0$.  Suppose inductively that it
has been continued meromorphically to $H_r$.  If $s\in H_{r+1}$, then
$s+m\in H_r$ for every $m\geq1$.  On a compact subset of $H_{r+1}$,
only finitely many low-shift terms in \cref{eq:Fq-Sqx} can meet poles of
the existing continuation; by
\cref{lem:Fq-normal-shifted-tail}, the remaining holomorphic tail
converges normally.  Thus $S_q(\,\cdot\,,x)$ is meromorphic on
$H_{r+1}$, and
\begin{equation}
 \label{eq:Fq-continuation-recursion}
 F_q(s,x)=\frac{A_q(s,x)\zeta(s)+(q/2)S_q(s,x)}
                  {2^s-2^{\sigma_q}}
\end{equation}
defines a meromorphic extension to $H_{r+1}$.  It agrees with the existing
function on $H_r$, so the identity theorem gives uniqueness on the overlap.
Since $\bigcup_{r\geq0}H_r=\C$, induction proves the continuation.

At the first step, new candidates can occur only at the pole $1$ of
$\zeta$ and at the zero set $\Lambda_q$ of the denominator.  Each later
shifted term replaces an existing candidate $p$ by $p-m$, $m\geq1$,
while division introduces $\Lambda_q$ again.  Induction therefore gives
the containing set in \cref{eq:Fq-candidate-poles}; it does not supply
nonvanishing of the corresponding numerators.
\end{proof}

\subsection{Corrected principal parts for \texorpdfstring{$q=3$}{q=3}}

Write $L=\log 2$ and use $S_3$ as in \cref{eq:Fq-Sqx}.  Then
\begin{equation}
 \label{eq:F3-corrected-continuation}
 F_3(s,x)=\frac{
 \frac12\bigl((2x+1)-2^s(x+1)\bigr)\zeta(s)
 +(3/2)S_3(s,x)}{2^s-2}.
\end{equation}

\begin{proposition}[the corrected $q=3$ principal parts]
\label{prop:F3-corrected-principal-parts}
Let $x\in\C$ and let $\gamma$ be Euler's constant.
\begin{enumerate}[label=\textup{(\roman*)}]
\item At $s=1$, $F_3$ has the genuine double-pole expansion
\begin{equation}
 \label{eq:F3-principal-part-one}
 F_3(s,x)=-\frac1{4L}(s-1)^{-2}
 +\left(
   \frac{3S_3(1,x)}{4L}-\frac x2-\frac38-\frac{\gamma}{4L}
  \right)(s-1)^{-1}+O(1).
\end{equation}
\item If $s_k=1+2\pi ik/L$ with $k\in\Z\setminus\{0\}$, then
$s_k$ is at most a simple pole and
\begin{equation}
 \label{eq:F3-nonreal-lattice-residue}
 \operatorname{Res}_{s=s_k}F_3(s,x)
 =\frac{3S_3(s_k,x)-\zeta(s_k)}{4L}.
\end{equation}
It is removable precisely when $3S_3(s_k,x)=\zeta(s_k)$.
\item The point $s=0$ is a genuine simple pole and
\begin{equation}
 \label{eq:F3-zero-residue}
 \operatorname{Res}_{s=0}F_3(s,x)=-\frac3{16L}.
\end{equation}
\item Every $-k$, $k\in\N_0$, is at most a simple pole.  With
$R_k=\operatorname{Res}_{s=-k}F_3(s,x)$, where $R_k=0$ for a removable
candidate,
\begin{equation}
 \label{eq:F3-negative-residue-recurrence}
 R_k=\frac{3}{16L(1-2^{k+1})(k+1)}
 +\frac{3}{2(1-2^{k+1})}
  \sum_{j=0}^{k-1}2^j\binom{k}{j}R_j.
\end{equation}
The sum is empty for $k=0$.  In particular, every $R_k$ is independent
of $x$; the recurrence alone does not prove $R_k\ne0$ for all $k$.
\end{enumerate}
\end{proposition}

\begin{proof}
Near $s=1$, put $t=s-1$.  Since $S_3$ is holomorphic there,
\begin{align*}
 \frac{(2x+1)-2^s(x+1)}2
   &=-\frac12-(x+1)Lt+O(t^2),\\
 2^s-2&=2Lt+L^2t^2+O(t^3),\\
 \zeta(s)&=t^{-1}+\gamma+O(t).
\end{align*}
Division gives \cref{eq:F3-principal-part-one}; its nonzero
$(s-1)^{-2}$ coefficient proves that the pole is genuinely double.

At $s_k$ with $k\ne0$, both $S_3$ and $\zeta$ are holomorphic,
the denominator has derivative $2L$, and the coefficient of $\zeta$ in
the numerator equals $-1/2$.  This proves
\cref{eq:F3-nonreal-lattice-residue} and the stated removability
criterion.

At $s=0$, only the $m=1$ term of $S_3$ can inherit the double pole at
$1$ to order one:
\[
 -\frac{s}{2}F_3(s+1,x)=\frac1{8L}s^{-1}+O(1).
\]
All $m\geq2$ terms are holomorphic there.  Since $2^0-2=-1$,
\cref{eq:F3-zero-residue} follows.

Finally, fix $k\geq0$ and put $s=-k+t$.  Inductively assume the
candidates $0,-1,\ldots,-(k-1)$ are at most simple.  For
$1\leq m\leq k$,
\[
 \left(-\frac12\right)^m\binom{-k+m-1}{m}
 =2^{-m}\binom{k}{m};
\]
these terms contribute
$\sum_{j=0}^{k-1}2^{-(k-j)}\binom{k}{j}R_j$ to the residue of $S_3$.
For $m=k+1$,
\[
 \left(-\frac12\right)^{k+1}\binom{s+k}{k+1}
 =-\frac{t}{2^{k+1}(k+1)}+O(t^2),
\]
and \cref{eq:F3-principal-part-one} shows that its contribution is
$1/(2^{k+3}(k+1)L)$.  Terms $m\geq k+2$ are holomorphic.  Therefore
\[
 \operatorname{Res}_{s=-k}S_3(s,x)
 =\frac1{2^{k+3}(k+1)L}
  +\sum_{j=0}^{k-1}2^{-(k-j)}\binom{k}{j}R_j.
\]
The zeta term in \cref{eq:F3-corrected-continuation} is holomorphic at
$-k$, and
$2^{-k}-2=(1-2^{k+1})/2^k$.  Multiplication by
$(3/2)/(2^{-k}-2)$ yields
\cref{eq:F3-negative-residue-recurrence} and completes the induction.
\end{proof}

\begin{sourceissue}[the printed $q=3$ pole claims and recurrence]
The source's $q=3$ residues inherit the false $x$-term and therefore
contain spurious factors depending on $4x+1$.  Its headline assertion that
every nonreal lattice point and every nonpositive integer is an actual pole
does not prove noncancellation.  Moreover, the recurrence displayed in the
corollary is undefined at its stated initial index and conflicts with the
formula reached in its proof.  This edition retains the genuine double pole
at $1$, the genuine simple pole at $0$, the exact Laurent coefficients, the
candidate-pole containment, and the corrected recurrence
\cref{eq:F3-negative-residue-recurrence}; it does not turn unproved
noncancellation into a pole classification.
\end{sourceissue}

\section{\texorpdfstring{$(p,q)$-adic}{(p,q)-adic} Fourier analysis and
Wiener--Tauberian cyclicity}
\label{sec:pq-fourier}

Siegel's phrase \emph{$(p,q)$-adic} means that the domain topology and the
value-field topology come from different primes.  The cleanest theorem in
this strand is a Wiener--Tauberian theorem for continuous functions on
$\Z_p$ with values in a complete nonarchimedean field of residue
characteristic $q\ne p$.  We reconstruct that theorem after fixing the
character and Haar normalizations.

\subsection{The two-place setting}

Let $p$ be prime.  Let $K$ be a complete algebraically closed
nonarchimedean field whose residue characteristic is a prime $q\ne p$.
The strongest headline assumptions in the source also take $K$ to be
spherically incomplete, as is the standard completed algebraic closure
$\C_q$.  Since $p$ is a unit in $K$, every $p^N$ has absolute value $1$ in
$K$.

As an abstract group, the Pontryagin dual of the additive compact group
$\Z_p$ is
\begin{equation}
\label{eq:dual-Zp}
 \Gamma_p=\Z[1/p]/\Z.
\end{equation}
Because $\operatorname{char}K\ne p$ and $K$ is algebraically closed, choose
roots of unity
\begin{equation}
\label{eq:compatible-p-power-roots}
 \zeta_0=1,
 \qquad \operatorname{ord}(\zeta_N)=p^N,
 \qquad \zeta_{N+1}^{p}=\zeta_N
 \quad(N\geq0).
\end{equation}
Such a system can be constructed recursively.  First choose a nontrivial
root $\zeta_1$ of $X^p-1$; it exists because this polynomial has $p$
distinct roots in $K$.  After $\zeta_N$ is chosen with $N\geq1$, take any
root of $X^p-\zeta_N$ in $K$.  Its order divides $p^{N+1}$ and cannot divide
$p^N$, so it has order exactly $p^{N+1}$.

For $t=a/p^N+\Z\in\Gamma_p$ and $z\in\Z_p$, let $z_N$ be the image of $z$
in $\Z/p^N\Z$, choose any integer lift $\widetilde z_N$ of $z_N$, and set
\begin{equation}
\label{eq:K-valued-character-pairing}
 \left\langle \frac{a}{p^N}+\Z,z\right\rangle
 :=\zeta_N^{a\widetilde z_N}\in K^\times.
\end{equation}

\begin{proposition}[the chosen $K$-valued duality]
\label{prop:K-valued-character-duality}
Formula \cref{eq:K-valued-character-pairing} is independent of every chosen
representative and defines a pairing that is a group homomorphism in each
variable,
\[
 \langle\ ,\ \rangle:\Gamma_p\times\Z_p\longrightarrow K^\times.
\]
For each $t$, the map $\chi_t(z)=\langle t,z\rangle$ is locally constant,
and
\begin{equation}
\label{eq:K-character-classification}
 \Gamma_p\xrightarrow{\ \sim\ }
 \operatorname{Hom}_{\mathrm{lc}}\bigl((\Z_p,+),K^\times\bigr),
 \qquad t\longmapsto\chi_t,
\end{equation}
is an isomorphism of groups.
\end{proposition}

\begin{proof}
Changing $a$ or $\widetilde z_N$ by a multiple of $p^N$ does not change
\cref{eq:K-valued-character-pairing}, because $\zeta_N$ has order $p^N$.
Suppose the same $t$ is also represented by $b/p^M+\Z$.  At a common level
$L\geq M,N$,
\[
 ap^{L-N}\equiv bp^{L-M}\pmod{p^L},
 \qquad \zeta_L^{p^{L-N}}=\zeta_N,
 \qquad \zeta_L^{p^{L-M}}=\zeta_M.
\]
Reading both values at level $L$ proves independence of the denominator and
hence of the representative of $t$.  At a common level, addition in either
variable adds exponents modulo $p^L$; this proves biadditivity.  If
$p^Nt=0$, then $\chi_t$ factors through $\Z_p/p^N\Z_p$, so it is locally
constant.

If $\chi_t$ is trivial, evaluation at $1$ in any presentation
$t=a/p^N+\Z$ gives $\zeta_N^a=1$, hence $p^N\mid a$ and $t=0$; thus
\cref{eq:K-character-classification} is injective.  Conversely, let
$\chi:\Z_p\to K^\times$ be a locally constant homomorphism.  Local
constancy at $0$ gives $p^N\Z_p\subseteq\ker\chi$ for some $N$, so $\chi$
factors through the cyclic group $\Z/p^N\Z$.  Put $u=\chi(1)$.  Then
$u^{p^N}=1$.  The $p^N$ distinct powers of $\zeta_N$ exhaust the roots of
$X^{p^N}-1$ in $K$, so $u=\zeta_N^a$ for a unique $a$ modulo $p^N$.
Therefore $\chi(z)=\zeta_N^{a\widetilde z_N}=\chi_{a/p^N+\Z}(z)$, proving
surjectivity and uniqueness.
\end{proof}

The complex-looking notation $e^{2\pi i\{tz\}_p}$ in the source denotes
the element \cref{eq:K-valued-character-pairing}; it is not a complex
exponential living in a $q$-adic field.  This compatible-root construction,
rather than the source's characteristic-zero embedding notation, also works
when $K$ has positive characteristic different from $p$.

There is a unique normalized $K$-valued Haar integral, specified by
\begin{equation}
\label{eq:K-Haar-ball}
 \int_{k+p^N\Z_p}dz=p^{-N}.
\end{equation}
For continuous $f:\Z_p\to K$ it is the following limit in $K$:
\begin{equation}
\label{eq:K-Haar-Riemann}
 \int_{\Z_p}f(z)\,dz
 =\lim_{N\to\infty}\frac1{p^N}\sum_{k=0}^{p^N-1}f(k).
\end{equation}

Indeed, on functions constant modulo $p^N$ the right side of
\cref{eq:K-Haar-Riemann} is $p^{-N}\sum_{k\bmod p^N}f(k)$; refinement from
level $N$ to $N+1$ repeats every value $p$ times.  Since $|p|=1$ in $K$,
these finite-level functionals have norm at most $1$ and extend uniquely from
the uniformly dense locally constant functions to $C(\Z_p,K)$.  The
extension is translation-invariant and sends $1$ to $1$.  Conversely, any
translation-invariant continuous functional sending $1$ to $1$ gives the
same mass to all $p^N$ residue classes, and their $p^N$ disjoint indicators
sum to $1$; their common mass is therefore $p^{-N}$.  Finally, the locally
constant functions $z\mapsto f([z]_{p^N})$, where
$[z]_{p^N}\in\{0,\ldots,p^N-1\}$ is the canonical residue representative,
converge uniformly to $f$.  Applying the continuous integral proves
\cref{eq:K-Haar-Riemann}.  This proves existence and uniqueness as well.

Put $\Gamma_p[p^N]=\{t\in\Gamma_p:p^Nt=0\}$; it consists of the $p^N$
classes $a/p^N+\Z$, $0\leq a<p^N$.

\begin{proposition}[finite character orthogonality]
\label{prop:finite-character-orthogonality}
For $N\geq0$ one has
\begin{align}
 \sum_{t\in\Gamma_p[p^N]}\langle t,x\rangle
 &=p^N\ind{x\in p^N\Z_p}
 &&(x\in\Z_p),\label{eq:finite-orthogonality-primal}\\
 \sum_{k\in\Z/p^N\Z}\langle t-u,k\rangle
 &=p^N\ind{t=u}
 &&(t,u\in\Gamma_p[p^N]).\label{eq:finite-orthogonality-dual}
\end{align}
In the second sum the residue class $k$ is evaluated through any lift in
$\Z_p$.
Consequently,
\begin{equation}
\label{eq:finite-character-indicator}
 \ind{z\equiv k\pmod{p^N}}
 =\frac1{p^N}\sum_{t\in\Gamma_p[p^N]}
   \langle t,z-k\rangle.
\end{equation}
\end{proposition}

\begin{proof}
Writing $x_N$ for any integer representing $x$ modulo $p^N$, the first sum
is
\[
 \sum_{a=0}^{p^N-1}(\zeta_N^{x_N})^a.
\]
If $x\in p^N\Z_p$, every summand is $1$.  Otherwise
$\zeta_N^{x_N}\ne1$, while $(\zeta_N^{x_N})^{p^N}=1$, so the finite
geometric sum is $0$.  This proves
\cref{eq:finite-orthogonality-primal}.  For the second identity write
$t-u=b/p^N+\Z$ and apply the same geometric-sum argument to
$\sum_{k=0}^{p^N-1}(\zeta_N^b)^k$.  Division by $p^N$ is valid because
$\operatorname{char}K\ne p$, and the first identity with $x=z-k$ gives
\cref{eq:finite-character-indicator}.
\end{proof}

\subsection{Fourier transform and inversion}

Let $C(\Z_p,K)$ carry the supremum norm.  Let
\[
 c_0(\Gamma_p,K)=\left\{a:\Gamma_p\to K:
 \text{for every }\epsilon>0,
 \#\{t:|a(t)|\geq\epsilon\}<\infty\right\}
\]
with its supremum norm.  Define
\begin{equation}
\label{eq:pq-Fourier}
 \widehat f(t)=\int_{\Z_p}f(z)\langle-t,z\rangle\,dz.
\end{equation}

\begin{theorem}[nonarchimedean Fourier isometry]
\label{thm:pq-Fourier-isometry}
The Fourier transform extends to an isometric isomorphism
\begin{equation}
\label{eq:Fourier-isometry}
 \mathscr F:C(\Z_p,K)\xrightarrow{\sim}c_0(\Gamma_p,K).
\end{equation}
For every $f\in C(\Z_p,K)$,
\begin{equation}
\label{eq:pq-Fourier-inversion}
 f(z)=\sum_{t\in\Gamma_p}\widehat f(t)\langle t,z\rangle,
\end{equation}
where the series converges uniformly.
\end{theorem}

\begin{proof}
Suppose first that $f$ is constant modulo $p^N$.  For
$t\in\Gamma_p[p^N]$, \cref{eq:K-Haar-Riemann} gives
\begin{equation}
\label{eq:finite-pq-transform}
 \widehat f(t)=\frac1{p^N}\sum_{k\bmod p^N}
 f(k)\langle-t,k\rangle.
\end{equation}
If $t\notin\Gamma_p[p^N]$, choose $M>N$ with
$t\in\Gamma_p[p^M]$.  In the level-$M$ Riemann sum, partitioning residues as
$k=r+p^Nh$ produces, for each $r\bmod p^N$, the factor
\[
 \sum_{h=0}^{p^{M-N}-1}\langle-t,p^Nh\rangle.
\]
This is a geometric sum with nontrivial ratio and ratio raised to
$p^{M-N}$ equal to $1$, so it vanishes.  Thus $\widehat f$ is supported on
$\Gamma_p[p^N]$.  Substitution of \cref{eq:finite-pq-transform} and
\cref{eq:finite-orthogonality-primal} now gives finite inversion.  Conversely,
\cref{eq:finite-orthogonality-dual} shows that every function supported on
$\Gamma_p[p^N]$ is the transform of its finite inverse series.

All chosen characters, $p^N$, and $p^{-N}$ have absolute value $1$ in $K$.
The ultrametric inequality applied to the forward and inverse finite
transforms therefore gives
\[
 \sup_t|\widehat f(t)|\leq\sup_z|f(z)|
 \quad\text{and}\quad
 \sup_z|f(z)|\leq\sup_t|\widehat f(t)|.
\]
Thus the finite transform is isometric.  Locally constant functions are
uniformly dense in $C(\Z_p,K)$, and finitely supported sequences are dense
in $c_0(\Gamma_p,K)$; every finite character set is contained in
$\Gamma_p[p^N]$ for some $N$.  Completion therefore extends the compatible
finite transforms to the isometric bijection \cref{eq:Fourier-isometry}.
For $a\in c_0$, the inverse family
$\sum_ta(t)\langle t,z\rangle$ converges uniformly, because the supremum norm
of every tail is at most the largest omitted $|a(t)|$.  This proves
\cref{eq:pq-Fourier-inversion}.
\end{proof}

The source develops the same conclusion through van der Put expansions.
For $n\geq1$, the relevant basis function is the indicator of the ball
$n+p^{\lambda_p(n)}\Z_p$; every continuous function has a unique uniformly
convergent expansion in these indicators with coefficients tending to zero.
That route supplies explicit formulas relating van der Put and Fourier
coefficients, but the finite-transform completion above exposes the norm
structure used by the Wiener theorem.

\subsection{Convolution algebra}

We use the following conventions throughout this chapter.  For
$a,b\in c_0(\Gamma_p,K)$,
\begin{equation}
\label{eq:c0-convolution}
 (a*b)(t)=\sum_{s\in\Gamma_p}a(s)b(t-s).
\end{equation}
For $f,g\in C(\Z_p,K)$ and $\mu\in C(\Z_p,K)'$,
\begin{align}
 (f*g)(z)&=\int_{\Z_p}f(z-y)g(y)\,dy,
 \label{eq:function-convolution}\\
 (f*\mu)(z)=(\mu*f)(z)
 &=\mu\bigl(y\longmapsto f(z-y)\bigr).
 \label{eq:function-measure-convolution}
\end{align}
The last line is a definition of both displayed notations; it fixes the
minus sign in the translate.  Uniform continuity of $f$ shows that
$z\mapsto f(z-\cdot)$ is continuous into $C(\Z_p,K)$ in the supremum norm,
so $f*\mu\in C(\Z_p,K)$ and
$\|f*\mu\|_\infty\leq\|f\|_\infty\|\mu\|$.
The Fourier--Stieltjes transform uses the same negative character as the
function transform:
\begin{equation}
\label{eq:Fourier-Stieltjes}
 \widehat\mu(t)=\mu\bigl(z\longmapsto\langle-t,z\rangle\bigr).
\end{equation}
In particular $|\widehat\mu(t)|\leq\|\mu\|$, because every chosen character
has supremum norm $1$.

\begin{proposition}[the three convolution identities]
\label{prop:typed-convolution-identities}
The convolution in \cref{eq:c0-convolution} is a well-defined continuous
bilinear map $c_0\times c_0\to c_0$, with
\begin{equation}
\label{eq:c0-convolution-bound}
 \|a*b\|_\infty\leq\|a\|_\infty\|b\|_\infty.
\end{equation}
With the Fourier and Fourier--Stieltjes signs fixed by
\cref{eq:pq-Fourier,eq:Fourier-Stieltjes}, one has
\begin{equation}
\label{eq:product-convolution}
 \widehat{fg}=\widehat f*\widehat g.
\end{equation}
Moreover,
\begin{equation}
\label{eq:function-convolution-Fourier}
 \widehat{f*g}(t)=\widehat f(t)\widehat g(t),
 \qquad
 \widehat{f*\mu}(t)=\widehat f(t)\widehat\mu(t).
\end{equation}
Thus \cref{eq:Fourier-isometry} is an isometric Banach-algebra isomorphism
from pointwise multiplication on $C(\Z_p,K)$ to convolution on
$c_0(\Gamma_p,K)$.
\end{proposition}

\begin{proof}
For fixed $t$, the family $a(s)b(t-s)$ tends to zero off finite subsets:
if $b\ne0$ and $|a(s)b(t-s)|\geq\epsilon$, then
$|a(s)|\geq\epsilon/\|b\|_\infty$, which occurs for only finitely many
$s$.  Completeness and the nonarchimedean criterion for summability therefore
give \cref{eq:c0-convolution}.  The ultrametric inequality gives
\cref{eq:c0-convolution-bound}.  If $b$ has finite support, $a*b$ is a finite
sum of translates of $a$ and hence belongs to $c_0$.  Approximating an
arbitrary $b\in c_0$ by finitely supported $b_n$ and using
\cref{eq:c0-convolution-bound} proves $a*b\in c_0$ and continuity.

If $f$ and $g$ are constant modulo $p^N$, finite inversion and
\cref{eq:finite-orthogonality-dual} give, for $t\in\Gamma_p[p^N]$,
\[
 \widehat{fg}(t)
 =\sum_{s\in\Gamma_p[p^N]}\widehat f(s)\widehat g(t-s).
\]
Both sides vanish outside a sufficiently large finite character group, so
this is \cref{eq:product-convolution} at finite level.  Uniformly approximate
$f,g$ by locally constant functions.  Pointwise multiplication, the Fourier
isometry, and sequence convolution are continuous, so the identity passes to
all $f,g\in C(\Z_p,K)$.

Fourier inversion gives, uniformly in $(z,y)\in\Z_p^2$,
\[
 f(z-y)=\sum_{s\in\Gamma_p}
 \widehat f(s)\langle s,z\rangle\langle-s,y\rangle.
\]
Applying the continuous functional $g(y)\,dy$, respectively $\mu$, term by
term yields
\[
 (f*g)(z)=\sum_s\widehat f(s)\widehat g(s)\langle s,z\rangle,
 \qquad
 (f*\mu)(z)=\sum_s\widehat f(s)\widehat\mu(s)\langle s,z\rangle.
\]
The coefficient families lie in $c_0$ because $\widehat f\in c_0$ and
$\widehat g,\widehat\mu$ are bounded.  Uniqueness in Fourier inversion proves
\cref{eq:function-convolution-Fourier}.
\end{proof}

For $s\in\Gamma_p$, write
\[
 (\tau_sa)(t)=a(t+s).
\]

\begin{proposition}[finite translates and the closed principal ideal]
\label{prop:finite-translates-principal-closure}
For every $a\in c_0(\Gamma_p,K)$,
\begin{align}
 \tau_sa&=a*\delta_{-s}\qquad(s\in\Gamma_p),
 \label{eq:translate-pointmass}\\
 \Span_K\{\tau_sa:s\in\Gamma_p\}
 &=a*c_{00}(\Gamma_p,K),\label{eq:translate-span-c00}\\
 \overline{\Span_K\{\tau_sa:s\in\Gamma_p\}}
 &=\overline{a*c_0(\Gamma_p,K)}.\label{eq:translate-ideal-closure}
\end{align}
The closures are taken in the supremum norm.  Thus the finite translate span
is generally not the algebraic principal ideal $a*c_0$; its closure is the
closed principal ideal generated by $a$.
\end{proposition}

\begin{proof}
For $t\in\Gamma_p$,
\[
 (a*\delta_{-s})(t)
 =\sum_{u\in\Gamma_p}a(u)\delta_{-s}(t-u)
 =a(t+s),
\]
which proves \cref{eq:translate-pointmass}.  Every element of $c_{00}$ is a
finite linear combination of point masses, so bilinearity gives
\cref{eq:translate-span-c00} in both directions.

Let $T_a:c_0(\Gamma_p,K)\to c_0(\Gamma_p,K)$ be $T_a(b)=a*b$.
The convolution estimate gives
$\|T_a(b)\|_\infty\leq\|a\|_\infty\|b\|_\infty$, hence $T_a$ is continuous.
For every $b\in c_0$, choose $b_N\in c_{00}$ with
$\|b_N-b\|_\infty\to0$.  Then $T_a(b_N)\to T_a(b)$, so
$a*c_0\subseteq\overline{a*c_{00}}$.  The reverse inclusion before taking
closures follows from $c_{00}\subseteq c_0$.  This proves
\cref{eq:translate-ideal-closure}.
\end{proof}

\begin{remark}
The distinction cannot be removed even for the convolution unit.  If
$a=\delta_0$, then its finite translate span is
$c_{00}(\Gamma_p,K)$, whereas
$\delta_0*c_0(\Gamma_p,K)=c_0(\Gamma_p,K)$.  The group $\Gamma_p$ is
infinite, so these two spaces are unequal.
\end{remark}

\begin{theorem}[Wiener--Tauberian theorem for continuous $(p,q)$-adic functions]
\label{thm:pq-WTT-function}
For $\chi\in C(\Z_p,K)$ the following are equivalent.
\begin{enumerate}[label=\textup{(\roman*)}]
\item $1/\chi\in C(\Z_p,K)$;
\item $\widehat\chi$ is invertible in the convolution algebra
      $c_0(\Gamma_p,K)$;
\item $\Span_K\{\tau_s\widehat\chi:s\in\Gamma_p\}$ is dense in
      $c_0(\Gamma_p,K)$;
\item $\chi(z)\ne0$ for every $z\in\Z_p$.
\end{enumerate}
\end{theorem}

\begin{proof}
The transform of $1$ is the convolution unit $\delta_0$.  If (i) holds,
\cref{eq:product-convolution} shows that
$\widehat{1/\chi}$ is a convolution inverse of $\widehat\chi$.  Conversely,
if $b*\widehat\chi=\delta_0$, surjectivity of \cref{eq:Fourier-isometry}
gives $b=\widehat g$ for a unique $g\in C(\Z_p,K)$; injectivity and
\cref{eq:product-convolution} then give $g\chi=1$, so $g=1/\chi$.  This proves
the equivalence of (i) and (ii).

If (ii) holds, then every $a\in c_0$ equals
\[
 a=\widehat\chi*(\widehat\chi^{-1}*a).
\]
Choose $b_N\in c_{00}$ converging to
$\widehat\chi^{-1}*a$ in the supremum norm.  Boundedness of convolution gives
\[
 \widehat\chi*b_N\longrightarrow
 \widehat\chi*(\widehat\chi^{-1}*a)=a.
\]
By \cref{eq:translate-span-c00}, every $\widehat\chi*b_N$ is a finite linear
combination of translates of $\widehat\chi$.  This proves (iii).

Suppose (iii) holds but $\chi(z_0)=0$.  Fourier inversion makes evaluation
at $z_0$ the nonzero continuous functional
\[
 L_{z_0}(a)=\sum_{t\in\Gamma_p}a(t)\langle t,z_0\rangle
 \quad(a\in c_0).
\]
The sum converges and $|L_{z_0}(a)|\leq\|a\|_\infty$ by the same
nonarchimedean tail criterion used in Fourier inversion.
For every $s$, a change of variables gives
$L_{z_0}(\tau_s\widehat\chi)
=\langle-s,z_0\rangle\chi(z_0)=0$.  Thus $L_{z_0}$ vanishes on a dense
subspace and hence on all of
$c_0$, contradicting $L_{z_0}(\delta_0)=1$.  Therefore (iii) implies (iv).
Finally, if $\chi$ is continuous and nowhere zero on compact $\Z_p$, then
$1/\chi$ is continuous (equivalently, $|\chi|$ is bounded away from zero),
so (iv) implies (i).
\end{proof}

\subsection{Measures and the proposed measure theorem}

A $K$-valued measure in the paper is a continuous functional
$\mu\in C(\Z_p,K)'$.  Its Fourier--Stieltjes transform is the precisely
signed functional in \cref{eq:Fourier-Stieltjes}; integral notation writes
the same value as
$\int_{\Z_p}\langle-t,z\rangle\,d\mu(z)$.
Via \cref{eq:Fourier-isometry} and the nonarchimedean duality
$c_0'=\ell^\infty$, this transform is naturally bounded, not necessarily
vanishing at infinity.

Let
\begin{equation}
\label{eq:Dirichlet-kernel-padic}
 D_N(z)=p^N\ind{z\in p^N\Z_p}.
\end{equation}
Then
\begin{equation}
\label{eq:measure-local-density}
 (D_N*\mu)(z)=p^N\mu(z+p^N\Z_p)
\end{equation}
with no sign ambiguity.  Here $\mu(U)=\mu(\ind U)$ for a clopen set $U$.
Indeed, the convention \cref{eq:function-measure-convolution} gives
\[
 (D_N*\mu)(z)
 =\mu\bigl(y\longmapsto p^N\ind{z-y\in p^N\Z_p}\bigr)
 =p^N\mu(z+p^N\Z_p),
\]
because $z-y\in p^N\Z_p$ if and only if
$y\in z+p^N\Z_p$.  The source calls the limit, when it exists, the
pointwise Radon--Nikodym derivative of $\mu$ at $z$.

\begin{status}[measure Wiener--Tauberian claim]
The paper asserts that the translates of $\widehat\mu$ are dense enough to
approximate every element of $c_0(\Gamma_p,K)$ if and only if every existing
limit
\[
 \lim_{N\to\infty}(D_N*\mu)(z)
\]
is nonzero.  Since $\widehat\mu$ need not itself lie in $c_0$, the precise
typed density condition is
\begin{equation}
\label{eq:measure-cyclicity-typed}
 c_0(\Gamma_p,K)
 \subseteq\overline{\Span_K\{\tau_s\widehat\mu:s\in\Gamma_p\}}^{\,\ell^\infty}.
\end{equation}
The proof establishes the forward implication from
\cref{eq:measure-cyclicity-typed} to nonvanishing of every existing local
density limit.
\end{status}

\begin{sourceissue}[the converse for measures is not proved]
The purported second half assumes both a zero local-density limit and the
density of translates, then derives a contradiction.  This is the
contrapositive of the already proved implication; it is not the converse
``no zero limit implies density.''  The proposed condition is also vacuous
for a measure for which no limit in \cref{eq:measure-local-density} exists.
A valid converse requires an annihilator/closed-ideal argument together with
an existence or regularity hypothesis.  The measure equivalence is therefore
not promoted to a theorem in this edition.
\end{sourceissue}

\subsection{The characteristic distribution of a numen}

Suppose $X_{H,\ell}:\Z_p\to K_\ell$ is measurable and satisfies
\cref{eq:numen-FE-adic}.  Push normalized real Haar probability forward by
$X_{H,\ell}$ and call the resulting probability measure $\nu_{H,\ell}$.
For an additive character $\gamma_t$ of $K_\ell$, define
\begin{equation}
\label{eq:numen-characteristic}
 \widehat\nu_{H,\ell}(t)
 =\int_{\Z_p}\gamma_t(-X_{H,\ell}(z))\,dz.
\end{equation}
Partitioning $\Z_p$ into $j+p\Z_p$ and using the numen recursion gives
\begin{equation}
\label{eq:numen-characteristic-FE}
 \widehat\nu_{H,\ell}(t)
 =\frac1p\sum_{j=0}^{p-1}
 \gamma_t(-c_j)\widehat\nu_{H,\ell}(r_j^*t),
\end{equation}
where $r_j^*$ is pullback on the character group:
$\gamma_{r_j^*t}(x)=\gamma_t(r_jx)$.  When $r_j$ is scalar multiplication,
this is the source's shorter notation $r_jt$.

Equivalently, the pushforward measure satisfies
\begin{equation}
\label{eq:self-similar-measure}
 \nu_{H,\ell}=\frac1p\sum_{j=0}^{p-1}(H_j)_*\nu_{H,\ell}.
\end{equation}
Existence follows from the numen and its functional equation.  Uniqueness of
a probability measure solving \cref{eq:self-similar-measure} is a separate
IFS theorem and requires, for example, that all $H_j$ be strict contractions
on a complete metric space.  The latest manual's unconditional uniqueness
wording is not used without such a hypothesis.

\section{Rising-continuity, frames, and products of F-series}
\label{sec:rising-fseries}

The 2022 dissertation, the 2023 paper on variable topologies, and the 2025
paper on products all study series whose target topology depends on the input
digit sequence.  The stable core is: canonical truncation limits,
point-selected completions (frames), digit-multiplicative F-series, and
finite-level Fourier recursions.  Several Banach-algebra and product-of-
distributions conclusions in the sources require repairs recorded below.

\subsection{Rising-continuity}

Let $p$ be prime, let $K$ be a complete nonarchimedean field, and let
$f:\Z_p\to K$.

\begin{definition}[rising-continuity]
The function $f$ is \emph{rising-continuous} when
\begin{equation}
\label{eq:rising-continuity}
 f(z)=\lim_{N\to\infty}f([z]_{p^N})
 \qquad(z\in\Z_p).
\end{equation}
Write $\widetilde C(\Z_p,K)$ for the class of such functions.
\end{definition}

For an arbitrary function $f:\Z_p\to K$, define the source's level
differences by
\begin{align}
 (\nabla_{p^0}f)(z)&=f(0),\label{eq:del-zero}\\
 (\nabla_{p^N}f)(z)
 &=f([z]_{p^N})-f([z]_{p^{N-1}}),
 &&N\geq1,\label{eq:del-positive}
\end{align}
and its extended del quantity by
\begin{equation}
\label{eq:del-norm}
 \|f\|_{\nabla}
 =\sup_{N\geq0}\sup_{z\in\Z_p}|(\nabla_{p^N}f)(z)|_K
 \in[0,\infty].
\end{equation}

\begin{proposition}[the exact domain of the del norm]
\label{prop:del-norm-domain}
If $f\in\widetilde C(\Z_p,K)$, then
\begin{equation}
\label{eq:del-equals-sup}
 \|f\|_{\nabla}=\sup_{z\in\Z_p}|f(z)|_K
\end{equation}
as an equality in $[0,\infty]$.  Consequently
\[
 \widetilde C_b(\Z_p,K)
 :=\{f\in\widetilde C(\Z_p,K):\|f\|_\nabla<\infty\}
\]
is precisely the bounded rising-continuous class, and it is a Banach space
under $\|\cdot\|_\nabla$.
\end{proposition}

\begin{proof}
For every $N$ and $z$, the ultrametric inequality bounds
$|(\nabla_{p^N}f)(z)|_K$ by $\sup_w|f(w)|_K$, including $N=0$.  Conversely,
the finite telescoping identity
\[
 f([z]_{p^N})=\sum_{h=0}^N(\nabla_{p^h}f)(z)
\]
and the ultrametric inequality give
$|f([z]_{p^N})|_K\leq\|f\|_\nabla$.  Rising-continuity and continuity of
the absolute value let $N\to\infty$, proving the reverse inequality in
\cref{eq:del-equals-sup}.  Thus the two finite norms agree.  Bounded
$K$-valued functions are complete under the supremum norm.  If $(f_m)$ is a
uniformly Cauchy sequence of bounded rising-continuous functions with uniform
limit $f$, then for fixed $z$ the ultrametric estimate
\[
 |f(z)-f([z]_{p^N})|_K
 \leq\max\{|f(z)-f_m(z)|_K,
            |f_m(z)-f_m([z]_{p^N})|_K,
            |f_m([z]_{p^N})-f([z]_{p^N})|_K\}
\]
first becomes small by choosing $m$ and then $N$.  Hence $f$ is
rising-continuous, so the displayed subspace is closed and complete.
\end{proof}

Every continuous function is rising-continuous because
$[z]_{p^N}\to z$ in $\Z_p$.  The converse is false: the condition tests one
distinguished convergent sequence at each point, not every convergent
sequence.  For every odd prime $q$, the shortened numen
$\chi_q:\Z_2\to\Z_q$ is the principal example (the preceding definition with
$p=2$ and $K=\Q_q$): it is rising-continuous everywhere but discontinuous at
each nonnegative integer.

Define the truncation of $f$ at level $N$ by
\begin{equation}
\label{eq:function-truncation}
 f_N(z)=f([z]_{p^N})
 =\sum_{k=0}^{p^N-1}f(k)\ind{z\equiv k\pmod{p^N}}.
\end{equation}

\begin{proposition}[pointwise versus uniform truncation]
\label{prop:rising-vs-continuous}
For an arbitrary $f:\Z_p\to K$:
\begin{enumerate}[label=\textup{(\roman*)}]
\item $f_N\to f$ pointwise if and only if $f$ is rising-continuous;
\item $f_N\to f$ uniformly if and only if $f$ is continuous.
\end{enumerate}
\end{proposition}

\begin{proof}
Part (i) is the definition.  Every $f_N$ is locally constant and hence
continuous, so a uniform limit is continuous.  Conversely, a continuous
function on compact $\Z_p$ is uniformly continuous.  Since
$|z-[z]_{p^N}|_p\leq p^{-N}$ uniformly in $z$, uniform continuity gives
$\sup_z|f(z)-f_N(z)|\to0$.
\end{proof}

\subsection{Van der Put paths}

Set $\lambda_p(0)=0$ and
\[
 \lambda_p(n)=1+\lfloor\log_p n\rfloor\qquad(n\geq1).
\]
For $n\geq1$, let $n^-$ be the integer obtained by deleting the highest
nonzero base-$p$ digit of $n$.  Define
\begin{align}
 v_0(z)&=1,\\
 v_n(z)&=\ind{z\equiv n\pmod{p^{\lambda_p(n)}}},\qquad n\geq1,\label{eq:vdp-basis}\\
 c_0(f)&=f(0),\\
 c_n(f)&=f(n)-f(n^-),\qquad n\geq1.
\end{align}

\begin{theorem}[rising van der Put representation]
\label{thm:rising-vdp}
A function $f:\Z_p\to K$ is rising-continuous if and only if
\begin{equation}
\label{eq:rising-vdp}
 f(z)=\sum_{n=0}^{\infty}c_n(f)v_n(z)
\end{equation}
converges at every $z$.  Equivalently, along every truncation path,
\begin{equation}
\label{eq:vdp-path-condition}
 c_{[z]_{p^N}}(f)\,ind{\lambda_p([z]_{p^N})=N}
 \longrightarrow0.
\end{equation}
The expansion is unique.
\end{theorem}

\begin{proof}
At a fixed $z$, exactly one possible new van der Put ball appears when the
$N$th digit is nonzero: its index is $[z]_{p^N}$ and its coefficient is the
difference between two successive effective truncation values.  Therefore
the partial sum through level $N$ telescopes to $f([z]_{p^N})$.  In a
complete nonarchimedean field, the path series converges exactly when its
terms tend to zero, giving \cref{eq:vdp-path-condition}; its sum is $f(z)$
exactly when \cref{eq:rising-continuity} holds.  Evaluation successively on
the nested balls recovers every coefficient, proving uniqueness.
\end{proof}

\begin{corollary}[unique rising continuation]
\label{cor:unique-rising-continuation}
Let $g:\N_0\to K$.  A rising-continuous extension to $\Z_p$ exists if and
only if $g([z]_{p^N})$ converges in $K$ for every $z\in\Z_p$.  When it exists
it is unique and equals
\[
 \widetilde g(z)=\lim_{N\to\infty}g([z]_{p^N}).
\]
\end{corollary}

\begin{proof}
If a rising-continuous extension $f$ exists, then
$g([z]_{p^N})=f([z]_{p^N})\to f(z)$ for every $z$, so the displayed limit is
necessary and determines $f$ uniquely.  Conversely, suppose every displayed
limit exists and use it to define $\widetilde g$.  If $m\in\N_0$, then
$[m]_{p^N}=m$ for all sufficiently large $N$, hence
$\widetilde g(m)=g(m)$.  Therefore
\[
 \widetilde g([z]_{p^N})=g([z]_{p^N})\longrightarrow\widetilde g(z),
\]
which is exactly rising-continuity.  Thus $\widetilde g$ is the required
extension, and the necessity argument already proved uniqueness.
\end{proof}

This is the precise uniqueness condition missing from several functional-
equation statements in the 2026 manual: uniqueness on $\N_0$ does not by
itself imply uniqueness on all infinite digit strings, but rising-continuity
does.

\subsection{Frames: a topology selected at each point}

Let $F$ be a global field and let $\operatorname{Comp}(F)$ denote a chosen
collection of complete valued extensions of $F$.

\begin{definition}[frame]
An $F$-\emph{frame} on a set $X$ is a map
\[
 \mathcal F:X\to\operatorname{Comp}(F).
\]
A function $f$ is compatible with $\mathcal F$ when
$f(x)\in\mathcal F(x)$ for every $x$.  A sequence $f_N$ converges to $f$ in
the frame when
\begin{equation}
\label{eq:frame-convergence}
 f_N(x)\longrightarrow f(x)\quad\text{in }\mathcal F(x)
 \qquad(x\in X).
\end{equation}
\end{definition}

Thus a frame is not one exotic topology on a disjoint union.  It is a rule
for which ordinary completion is to judge convergence at each input.
Pointwise sums and products of compatible functions are meaningful because
both operands at a fixed point lie in the same selected field.

The standard two-place example in the body of the 2023 paper uses complex
convergence on $\N_0$ and $q$-adic convergence on
$\Z_p\setminus\N_0$.  Its abstract describes the assignments in the opposite
order; the body definition controls this edition.

\begin{sourceissue}[an unrestricted frame function need not have a sup norm]
One source proposition declares all functions from compact $X$ to a complete
field to be a Banach space under the supremum norm.  Compactness bounds
continuous functions, not arbitrary functions.  The statement becomes
valid after restricting to bounded functions (and proving completeness
pointwise plus uniformly), or to continuous functions in a fixed target.
The same issue reappears in the dissertation's proposed Banach algebra of all
rising-continuous functions: rising-continuity alone does not imply
boundedness or finiteness of the displayed ``del'' norm.
\end{sourceissue}

\subsection{Digit-multiplicative F-series}

For a digit prefix of fixed length $N$, use the padded count
\begin{equation}
\label{eq:padded-count}
 \#^{(N)}_{p:j}(z)=\#\{0\leq h<N:d_h(z)=j\}.
\end{equation}
This count includes high zero digits.  The ordinary digit count of the
integer $[z]_{p^N}$ silently discards them and is not interchangeable with
\cref{eq:padded-count} when a zero-digit multiplier occurs.

Fix a complete valued field $E$ containing nonzero parameters
$d,q_0,\ldots,q_{p-1}$.  For $z\in\Z_p$, put
\[
 u_N(z)=d^{-N}\prod_{j=0}^{p-1}
 q_j^{\#^{(N)}_{p:j}(z)},
 \qquad
 D_E=\left\{z\in\Z_p:\sum_{N\geq0}u_N(z)
                       \text{ converges in }E\right\}.
\]
Define the $E$-valued F-series on its exact convergence domain by
\begin{equation}
\label{eq:general-F-series}
 S^E_{d;\mathbf q}:D_E\longrightarrow E,
 \qquad S^E_{d;\mathbf q}(z)=\sum_{N=0}^{\infty}u_N(z).
\end{equation}
 For $0\leq k<p$, the digit identity
\[
 \#^{(N+1)}_{p:j}(pz+k)
 =\#^{(N)}_{p:j}(z)+\ind{j=k}
\]
gives the termwise tail identity
\[
 u_{N+1}(pz+k)=\frac{q_k}{d}u_N(z).
\]
Since $q_k/d\ne0$, it follows without an unstated convergence condition that
$pz+k\in D_E$ if and only if $z\in D_E$, and on this common domain
\begin{equation}
\label{eq:F-series-FE}
 S^E_{d;\mathbf q}(pz+k)=1+\frac{q_k}{d}S^E_{d;\mathbf q}(z).
\end{equation}
A frame may select a different completion at different inputs, but an
identity comparing $z$ with $pz+k$ is typed only after both values are placed
in one specified field $E$ (or after a specified identification of their
target fields).

The primitive numen series \cref{eq:partial-F-series} is the affine version:
its $N$th term contains a digit-dependent translation $c_{d_N}$ and the
ordered product of all earlier linear multipliers.  Fix a complete valued
field $E$ and scalars $a_0,\ldots,a_{p-1},b_0,\ldots,b_{p-1}\in E$.  In this
scalar commuting case put
\[
 D_X=\left\{z\in\Z_p:
 \sum_{N=0}^{\infty}b_{d_N(z)}
  \prod_{h=0}^{N-1}a_{d_h(z)}\text{ converges in }E\right\}
\]
and define
\begin{equation}
\label{eq:scalar-affine-F-series}
 X:D_X\longrightarrow E,\qquad
 X(z)=\sum_{N=0}^{\infty}b_{d_N(z)}
 \prod_{j=0}^{p-1}a_j^{\#^{(N)}_{p:j}(z)}.
\end{equation}
For $0\leq k<p$ and $z\in D_X$ one has $pz+k\in D_X$ and
\begin{equation}
\label{eq:scalar-affine-FE}
 X(pz+k)=a_kX(z)+b_k.
\end{equation}
Conversely, if $a_k\ne0$, then $pz+k\in D_X$ implies $z\in D_X$.
Indeed, the zeroth summand at $pz+k$ is $b_k$, and its summand of index
$N+1$ is $a_k$ times the summand of index $N$ at $z$.  Thus convergence in
the forward direction is automatic even when $a_k=0$, while division by
$a_k$ in the converse is valid exactly under the displayed nonvanishing
hypothesis.

For the Collatz numen, the 2023 paper writes the useful frame identity
\begin{equation}
\label{eq:chi3-F-series}
 \chi_3(z)=-\frac12+\frac14
 \sum_{N=0}^{\infty}\frac{3^{\#^{(N)}_{2:1}(z)}}{2^N}.
\end{equation}
On nonnegative integers the right side is read archimedeanly; on
$\Z_2\setminus\N_0$ it is read $3$-adically.  This differing-codomain
identity does not follow from \cref{cor:unique-rising-continuation}, whose
domain and codomain are fixed.  It follows directly from one finite identity.
Indeed, write
\[
 A_N(z)=\frac{3^{\#^{(N)}_{2:1}(z)}}{2^N},
 \qquad S_N(z)=\sum_{h=0}^{N-1}A_h(z),
\]
and let $d_h(z)\in\{0,1\}$ be the $h$th digit.  The $N$-digit numen partial
sum is
\[
 \chi_{3,N}(z)=\frac12\sum_{h=0}^{N-1}d_h(z)A_h(z).
\]
Since $A_{h+1}=(\frac12+d_h)A_h$ and $A_0=1$, telescoping gives the exact
finite relation
\begin{equation}
\label{eq:chi3-finite-frame-identity}
 \chi_{3,N}(z)=\frac14S_N(z)-\frac12+\frac12A_N(z).
\end{equation}
If $z\in\N_0$, its one-count is eventually constant and $A_N(z)\to0$ in
$\R$.  If $z\in\Z_2\setminus\N_0$, its binary expansion has infinitely many
one digits, so $|A_N(z)|_3=3^{-\#^{(N)}_{2:1}(z)}\to0$.  In the first case
$\chi_{3,N}(z)$ is eventually the finite numen value; in the second it
converges in $\Q_3$ by the primitive numen series.  Passing to the limit in
\cref{eq:chi3-finite-frame-identity} proves
\cref{eq:chi3-F-series} separately in the stated completion at every input.

\begin{sourceissue}[pointwise convergence claim in the 2023 paper]
One necessity proof chooses a general nonnatural input and then replaces it
by $z=-1$.  That proves a global obstruction, not the asserted pointwise
necessity for the original digit sequence.  Placewise convergence should be
checked from the actual prefix products, as in
\cref{prop:numen-pointwise-convergence}, rather than from that substitution.
\end{sourceissue}

\subsection{Products and the exact finite recursion}

Let $R$ be a commutative ring with identity, let
$a_{j,k},b_{j,k}\in R$, and let
$X_1,\ldots,X_d:\Z_p\to R$ be degree-one digit-recursive functions
satisfying
\begin{equation}
\label{eq:degree-one-recursion}
 X_j(pz+k)=a_{j,k}X_j(z)+b_{j,k}
 \qquad(1\leq j\leq d,\ 0\leq k<p).
\end{equation}
For a multi-index $\mathbf n=(n_1,\ldots,n_d)\in\N_0^d$, put
\[
 X_{\mathbf n}=\prod_{j=1}^dX_j^{n_j}.
\]
If $\mathbf m\leq\mathbf n$ coordinatewise, define
\begin{align}
 \binom{\mathbf n}{\mathbf m}&=\prod_j\binom{n_j}{m_j},\\
 a_{\mathbf m,k}&=\prod_ja_{j,k}^{m_j},\\
 b_{\mathbf n-\mathbf m,k}&=\prod_jb_{j,k}^{n_j-m_j},\\
 r_{\mathbf m,\mathbf n,k}
 &=\binom{\mathbf n}{\mathbf m}
   a_{\mathbf m,k}b_{\mathbf n-\mathbf m,k}.
\end{align}

\begin{proposition}[multinomial branch recursion]
\label{prop:product-branch-recursion}
For every $k$,
\begin{equation}
\label{eq:product-branch-recursion}
 X_{\mathbf n}(pz+k)
 =\sum_{\mathbf m\leq\mathbf n}
 r_{\mathbf m,\mathbf n,k}X_{\mathbf m}(z).
\end{equation}
\end{proposition}

\begin{proof}
Raise each equation in \cref{eq:degree-one-recursion} to the power $n_j$,
expand by the binomial theorem, multiply over $j$, and collect the term with
multi-index $\mathbf m$.
\end{proof}

This finite algebraic identity is independent of every limiting assumption
in the 2025 paper.

\subsection{Finite Fourier recursion and formal solutions}

Fix through the end of \cref{prop:finite-transform-limit} a complete
nonarchimedean field $K$ of residue
characteristic $q\ne p$, a compatible system of $p$-power roots of unity,
and the normalized $K$-valued Haar integral of
\cref{sec:pq-fourier}.  Assume the branch parameters and the functions under
transformation are $K$-valued.  A point-varying frame by itself is not a
single coefficient field in which Fourier sums can be added.

Let
\begin{equation}
\label{eq:alpha-product}
 \alpha_{\mathbf m,\mathbf n}(t)
 =\frac1p\sum_{k=0}^{p-1}
 r_{\mathbf m,\mathbf n,k}\langle-t,k\rangle,
 \qquad
 \alpha_{\mathbf n}=\alpha_{\mathbf n,\mathbf n}.
\end{equation}
Fix $\mathbf n$ and assume explicitly that
$X_{\mathbf m}\in C(\Z_p,K)$ for every $\mathbf m\leq\mathbf n$.  Then all
Fourier transforms below are the normalized Haar integrals of
\cref{eq:pq-Fourier}.  The substitution
$x=pz+k:\Z_p\xrightarrow{\sim}k+p\Z_p$ gives
\[
 \int_{k+p\Z_p}f(x)\,dx
 =\frac1p\int_{\Z_p}f(pz+k)\,dz
\]
for continuous $f$.  Applying this identity to
$f(x)=X_{\mathbf n}(x)\langle-t,x\rangle$, using
\cref{eq:product-branch-recursion}, and interchanging only finite sums gives
\begin{equation}
\label{eq:product-Fourier-recursion}
 \widehat X_{\mathbf n}(t)
 =\sum_{\mathbf m\leq\mathbf n}
 \alpha_{\mathbf m,\mathbf n}(t)
 \widehat X_{\mathbf m}(pt).
\end{equation}
At finite truncation level the identity is exact before any limiting
integrability claim, but it has a necessary one-level shift.  Put
\[
 X_{j,N}(z)=X_j([z]_{p^N}),\qquad
 X_{\mathbf n,N}(z)=\prod_{j=1}^dX_{j,N}(z)^{n_j}.
\]
Then
\begin{equation}
\label{eq:finite-product-Fourier-recursion}
 \widehat X_{\mathbf n,N+1}(t)
 =\sum_{\mathbf m\leq\mathbf n}
 \alpha_{\mathbf m,\mathbf n}(t)\widehat X_{\mathbf m,N}(pt).
\end{equation}
Indeed, $[pz+k]_{p^{N+1}}=p[z]_{p^N}+k$.  More generally, for
$f_N(z)=f([z]_{p^N})$,
\begin{equation}
\label{eq:finite-truncation-transform}
 \widehat f_N(t)=
 \frac{\ind{p^Nt=0}}{p^N}
 \sum_{r=0}^{p^N-1}f(r)\langle-t,r\rangle.
\end{equation}

At $t=0$, separate the diagonal term in
\cref{eq:product-Fourier-recursion}:
\begin{equation}
\label{eq:formal-mean-equation}
 (1-\alpha_{\mathbf n}(0))\widehat X_{\mathbf n}(0)
 =\sum_{\mathbf m<\mathbf n}
 \alpha_{\mathbf m,\mathbf n}(0)\widehat X_{\mathbf m}(0).
\end{equation}

\begin{proposition}[formal compatibility classification]
\label{prop:formal-compatibility}
Assume the lower-degree transforms are fixed.
\begin{enumerate}[label=\textup{(\roman*)}]
\item If $\alpha_{\mathbf n}(0)\ne1$, the formal recursion has a unique
      value at $0$, and hence a unique value at every $t\in\Gamma_p$ by
      repeated use of \cref{eq:product-Fourier-recursion} until $p^Nt=0$.
\item If $\alpha_{\mathbf n}(0)=1$ and the right side of
      \cref{eq:formal-mean-equation} is nonzero, no formal solution exists.
\item If $\alpha_{\mathbf n}(0)=1$ and that right side is zero, the value
      $\widehat X_{\mathbf n}(0)$ is free and the solutions form a
      one-dimensional affine family.
\end{enumerate}
\end{proposition}

\begin{proof}
\Cref{eq:formal-mean-equation} gives the three alternatives at
$0$.  Every $t\in\Gamma_p$ has $p$-power order, so $p^Nt=0$ for some $N$;
running \cref{eq:product-Fourier-recursion} backward from that terminal value
determines all other values.
\end{proof}

The paper calls the locus $\alpha_{\mathbf n}(0)=1$ a breakdown locus.
\Cref{prop:formal-compatibility} shows why breakdown does not automatically
mean nonexistence: the compatibility sum decides between no solution and an
affine family.

\subsection{M-functions and limits}

For scalar branch multipliers $r_0,\ldots,r_{p-1}\in K^\times$, the associated
M-function is the prefix cocycle
\begin{equation}
\label{eq:M-function}
 M_N(z)=\prod_{h=0}^{N-1}r_{d_h(z)}
 =r_0^N\prod_{j=1}^{p-1}\left(\frac{r_j}{r_0}\right)^{
 \#^{(N)}_{p:j}(z)}.
\end{equation}
It satisfies
\[
 M_{N+L}(z)=M_N(z)M_L(\sh^{\circ N}z).
\]
These cocycles are the proposed majorants for passing from
\cref{eq:finite-truncation-transform} to a framewise distributional limit.

Before stating the exact limit consequence, set
\begin{align*}
 \mathcal S(\Z_p,K)
 &=\{\phi:\Z_p\to K:\phi\text{ is locally constant modulo }p^N
       \text{ for some }N\},\\
 \mathcal S(\Z_p,K)'
 &=\operatorname{Hom}_K(\mathcal S(\Z_p,K),K).
\end{align*}

\begin{proposition}[the exact consequence of transform convergence]
\label{prop:finite-transform-limit}
Fix $\mathbf n$.  Suppose that for every $\mathbf m\leq\mathbf n$ and every
$t\in\Gamma_p$ the limit
\[
 A_{\mathbf m}(t)=\lim_{N\to\infty}
 \widehat X_{\mathbf m,N}(t)
\]
exists in the fixed complete field $K$.  Then
\begin{equation}
\label{eq:limit-product-Fourier-recursion}
 A_{\mathbf n}(t)=
 \sum_{\mathbf m\leq\mathbf n}
 \alpha_{\mathbf m,\mathbf n}(t)A_{\mathbf m}(pt)
 \qquad(t\in\Gamma_p).
\end{equation}
Moreover, every function $A:\Gamma_p\to K$ defines a unique algebraic
distribution $T_A\in\mathcal S(\Z_p,K)'$ by
\begin{equation}
\label{eq:formal-Fourier-distribution}
 \langle T_A,\phi\rangle
 =\sum_{t\in\Gamma_p}\widehat\phi(t)A(-t),
 \qquad \phi\in\mathcal S(\Z_p,K),
\end{equation}
where the sum is finite.  With the Fourier sign of
\cref{eq:pq-Fourier}, this convention gives
$\langle T_A,\langle-t,\cdot\rangle\rangle=A(t)$.  In particular, if
\[
 \langle\mu_{\mathbf m,N},\phi\rangle
 =\int_{\Z_p}X_{\mathbf m,N}(z)\phi(z)\,dz,
\]
then $\mu_{\mathbf m,N}\to T_{A_{\mathbf m}}$ pointwise on
$\mathcal S(\Z_p,K)$ for every $\mathbf m\leq\mathbf n$.
\end{proposition}

\begin{proof}
Take $N\to\infty$ in
\cref{eq:finite-product-Fourier-recursion}.  Its sum over
$\mathbf m\leq\mathbf n$ is finite, so convergence of the displayed terms is
all that is required and yields
\cref{eq:limit-product-Fourier-recursion}.  Every locally constant test
function has a unique finite character expansion
$\phi(z)=\sum_t\widehat\phi(t)\langle t,z\rangle$ by finite Fourier
inversion.  Hence \cref{eq:formal-Fourier-distribution} is a well-defined
$K$-linear functional, and evaluating it on the character
$\langle-t,\cdot\rangle$ proves the final identity and uniqueness.  Finally,
finite Fourier inversion for $\phi$ gives
\[
 \langle\mu_{\mathbf m,N},\phi\rangle
 =\sum_t\widehat\phi(t)\widehat X_{\mathbf m,N}(-t),
\]
where the sum is finite.  Termwise passage to the limit is therefore exact
and yields \cref{eq:formal-Fourier-distribution} with
$A=A_{\mathbf m}$.
\end{proof}

\begin{status}[conditional analytic theorem in the 2025 paper]
The paper explicitly introduces a Main Limit Assumption before asserting
that the formal Fourier data are realized by frame-quasi-integrable products.
\Cref{prop:finite-transform-limit} isolates exactly what follows if all of
the required finite-transform limits exist in one fixed completion.

The later Main Limit Lemma proposes sufficient conditions in terms of strong
frame summability of $M_N$, uniform control of shifted $X$, and
$(\mathcal F,M)$-type decay.  The printed proof does not establish those
sufficient conditions, so the ensuing base case, induction, power theorem,
and main product corollary remain conditional.
\end{status}

\subsection{Test functions, distributions, and what the product means}

More generally, for a commutative coefficient ring $R$ with identity, let
\begin{align}
 \mathcal S(\Z_p,R)
 &=\{\phi:\Z_p\to R:\phi\text{ is locally constant modulo }p^N
       \text{ for some }N\},\\
 \mathcal S(\Z_p,R)'&=\operatorname{Hom}_R(\mathcal S(\Z_p,R),R).
\end{align}
The second module is the algebraic space of $(p,R)$-adic distributions.  If
$\mu$ is a distribution and $\phi$ a test function, multiplication by a test
function is unambiguous:
\[
 \langle\phi\mu,\psi\rangle
 =\langle\mu,\phi\psi\rangle
 \qquad\bigl(\psi\in\mathcal S(\Z_p,R)\bigr).
\]
Here the distribution occupies the first slot and the test function the
second.  This convention is used throughout the subsection.
Products of two arbitrary distributions are not defined by this rule.

\begin{sourceissue}[the product algebra is transported, not intrinsic]
The 2025 paper defines, for its designated family,
\[
 (X^m(z)\,dz)(X^n(z)\,dz):=X^{m+n}(z)\,dz
\]
and defines the corresponding ``convolution powers'' to be the recursively
constructed formal Fourier data.  The exact formal object is the polynomial
algebra $R[T]$, with $T^n$ labelled by the symbol $[X^n(z)\,dz]$ and with
$T^mT^n=T^{m+n}$.  This construction exists without realizing any symbol as
a distribution.

Suppose separately that distributions
$\mu_n\in\mathcal S(\Z_p,R)'$ have actually been constructed for the named
symbols, and define the $R$-linear realization map
\[
 \Phi:R[T]\longrightarrow\mathcal S(\Z_p,R)',
 \qquad \Phi\!\left(\sum_na_nT^n\right)=\sum_na_n\mu_n.
\]
The rule
\begin{equation}
\label{eq:transported-product-rule}
 \mu_m\mathbin\odot\mu_n:=\mu_{m+n}
\end{equation}
descends to a well-defined multiplication on $\operatorname{im}\Phi$ if and
only if $\ker\Phi$ is an ideal of $R[T]$.  Indeed, well-defined descent is
equivalent to
$\Phi(f)=\Phi(f')\Rightarrow\Phi(fg)=\Phi(f'g)$ for every $g$, which is
equivalent to closure of $\ker\Phi$ under multiplication by every
polynomial.  When this holds, the realized algebra is canonically
$R[T]/\ker\Phi$; when $\Phi$ is injective it is the polynomial algebra
itself.  If $\ker\Phi$ is not an ideal, the declaration
\cref{eq:transported-product-rule} depends on the chosen polynomial
representative and is not a product on the image.

The paper also warns that ordinary convolution of its Fourier data generally
diverges.  The formal product above is therefore not a proof that arbitrary
distributions possess pointwise products, nor that standard Fourier
convolution realizes the declared multiplication.
\end{sourceissue}

\subsection{Failure and repair ledger for the product paper}

The following defects affect theorem status rather than notation alone.
\begin{enumerate}[label=\textup{(\arabic*)}]
\item Haar probability of a residue class modulo $p^N$ is printed in two
      places as $1/p$; it is $p^{-N}$.  The corresponding finite character
      expansion needs the same correction.
\item In the archimedean Wiener algebra, absolutely summable Fourier
      coefficients must carry the $\ell^1$ norm
      $\sum_t|\widehat f(t)|$, not the coefficient supremum norm printed in
      the paper.
\item The Main Limit Lemma repeatedly infers that a positive sequence tending
      to zero is monotone on a tail, and that a tail sum is
      $O$ of its first term.  Neither follows from convergence or from a root
      bound.  The base case also uses the false estimate $Na_N=O(a_N)$.
\item The localization $R[(1-a_0)^{-1}]$ can be quotiented by an ideal $I$
      only after requiring $I$ to be disjoint from every power of
      $1-a_0$; the weaker printed condition $1-a_0\notin I$ does not suffice.
      A fraction field of the quotient additionally requires a domain, for
      example a prime ideal.  Multivariable polynomial localizations are not
      Dedekind domains in dimension greater than one.
\item A final corollary says a breakdown index gives no solution for any
      initial data, contradicting the compatible affine-family case already
      proved in the paper and restated in
      \cref{prop:formal-compatibility}.
\end{enumerate}

A repair programme must first correct Haar and Wiener normalization, impose
domain-safe localization hypotheses, replace the false monotonicity steps by
a genuine monotone majorant or regular-variation estimate, and prove rising
continuity of the proposed limits.  Only then can the analytic product and
inversion conclusions be promoted from conditional claims.

\section{Synthesis: terminology, dependencies, and exact theorem status}
\label{sec:synthesis}

This section is the control panel for the edition.  It does not add another
layer of analogy to the preceding mathematics.  It identifies the same
objects in Siegel's terminology and in standard mathematical language,
records the hypotheses needed to pass from one object to the next, and marks
the precise point at which each proved chain stops.  A status assigned below
belongs to a particular statement in a particular witness; it is not a
judgment on an entire paper or on the programme as a whole.

\subsection{Authorship and editorial boundary}

\begin{status}[Authorship boundary]
The mathematical corpus edited in this volume is the work of Maxwell C.
Siegel.  In particular, the source constructions and the names
\emph{Hydra map}, \emph{numen}, \emph{dreamcatcher},
\emph{$(p,q)$-adic analysis}, \emph{rising-continuity}, and
\emph{F-series} remain attributed to Siegel.  The standard-language column
in the crosswalk below explains those objects; it does not rename them and
does not transfer their authorship.

Typed recastings with explicit source crosswalks, corrected quantifiers,
padded-word conventions,
replacement hypotheses, counterexamples to printed assertions, and proofs
explicitly supplied in this edition are editorial work.  They are described
as reconstructions or repairs, not retroactively attributed to Siegel as
source text.  Conversely, a correction supplied by this edition does not
change authorship of the underlying corpus.
\end{status}

\subsection{Terminology crosswalk}

The third column gives a standard description of the typed source object.
The fourth records an identification which must \emph{not} be made.

\begingroup\small
\setlength{\tabcolsep}{3pt}
\begin{longtable}{@{}p{0.18\textwidth}p{0.23\textwidth}p{0.28\textwidth}p{0.23\textwidth}@{}}
\toprule
Siegel term & Typed object in this edition & Standard-language description & Exactness boundary\\
\midrule
\endfirsthead
\toprule
Siegel term & Typed object in this edition & Standard-language description & Exactness boundary\\
\midrule
\endhead
Hydra map & A residue-class assignment of affine branches
$H_j=(r_j,c_j)$ on $\mathcal O_K$, admissible on the assigned classes & A
piecewise-affine arithmetic dynamical system with a finite digit alphabet &
The 2019 surjective $\varrho$-Hydra is a historical specialization; it is not
identical to the 2026 global-field definition.\\

Scalar integer Hydra & The specialization with $D\in\Z_{\geq2}$,
$K=\F=\Q$, $\OK=\Z$, $\Lambda=D\Z$, and rational-affine branches integral
on their assigned classes & Exactly a fixed-modulus Matthews generalized
Collatz presentation, hence an rcwa mapping in Kohl's mapping monoid; see
\cref{prop:scalar-hydra-matthews,prop:scalar-hydra-rcwa} & The displayed
modulus need not be minimal, an rcwa mapping need not be bijective, and rcwa
membership supplies none of the Hydra-specific centeredness, properness, or
branch-recognition predicates.  Exact-integrality of a displayed
$D$-presentation is equivalent to Matthews relatively prime type only together
with $D=d_{\mathrm{Mat}}(H)$.\\

Affine lattice morphism &
$\Aff_\F(K)=\operatorname{GL}_\F(K)\ltimes K$, acting by
$x\mapsto rx+c$ & An affine $\F$-linear automorphism of the affine space
underlying $K$ & The source phrase does not imply that such a map preserves
$\mathcal O_K$, $\Lambda$, or any lattice; each lattice condition is stated
separately.\\

Exact-integral Hydra &
$H_j(x)\in\mathcal O_K$ exactly when $x\in\Lambda_j$, for
$x\in\mathcal O_K$ & A branch recognition or branch-correctness condition &
Mere admissibility only says that the selected branch maps its selected class
into $\mathcal O_K$; the stronger field-wide predicate on all $x\in K$ is
not substituted.\\

Hydra word &
$H_{\mathbf j}=H_{j_0}\circ\cdots\circ H_{j_{L-1}}$ & A word in an affine
semigroup, or a symbolic cylinder word & The rightmost map acts first;
chronological itineraries must therefore be reversed under this convention.\\

$M_H$ & The linear part of $H_{\mathbf j}$ & A multiplicative linear cocycle
or monoid representation & In the global-field theory it is generally an
endomorphism, not a scalar.\\

Numen $X_H$ & The translation part
$H_{\mathbf j}(x)=M_H(\mathbf j)x+X_H(\mathbf j)$ & An affine translation
cocycle, equivalently $H_{\mathbf j}(0)$ for every affine word & Centeredness
is what makes this word cocycle descend through appended zero digits.  For a
noncentered Hydra the digit-recursive solution also contains the terminal
zero-branch fixed point.\\

Primitive F-series & The truncation expansion of $X_H$ along the digits of
$z$ & A digit-recursive affine series & Its infinite value has no meaning
until a codomain completion and a convergence condition are specified.\\

Correspondence Principle & Repeated symbolic code $\mapsto$ affine fixed
point, followed by branch correctness & A periodic symbolic-coding theorem
for a piecewise-affine map & A formal fixed point of a branch composition is
not automatically a periodic point of the piecewise map.\\

Classical parity conjugacy &
$Q_q:\Z_2\to\Z_2$ is the chronological parity encoder and
$\Phi_q=Q_q^{-1}$ conjugates the binary shift to $T_q$ & The
Bernstein--Lagarias solenoidal conjugacy, with the B\"ohm--Sontacchi return
formula at $q=3$ & For odd prime $q$, neither classical map is the cross-adic
numen $\chi_q:\Z_2\to\Z_q$.  Their purely periodic values have the same
rational representative under separate canonical embeddings only after
reversing the chronological word; see
\cref{thm:BL-numen-periodic-overlap,cor:BL-numen-global-nonidentity}.\\

Set-series & $\varsigma_V(w)=\sum_{v\in V}w^v$ and
$\psi_V(z)=\sum_{v\in V}e^{2\pi ivz}$ & Ordinary and Fourier generating
functions of an indicator sequence & The Dirichlet set-series needs a
separate convention when $0\in V$.\\

Permutation operator &
$\mathsf P_H:\ell^\infty(\N_0)\to\ell^\infty(\N_0)$,
$(\mathsf P_Hc)_n=c_{H(n)}$, together with its conjugates on the named
generating-map ranges & An isometric coefficient pullback when the 2019
Hydra $H$ is surjective & It is not the transfer operator which sums over
inverse branches, and the isometric injection $\mathsf P_H$ need not be
surjective.\\

Edgepoint and virtual residue &
$\operatorname{VRes}_x\psi=\lim_{y\downarrow0}y\psi(x+iy)$ when finite & A
radial first-order boundary coefficient & It differs by a factor from the
ordinary complex residue of a simple pole.\\

Dreamcatcher & The trace
$R_{\psi,T}(t)=\operatorname{VRes}_t[\psi]$ in $\ell^2(T/\Z)$ & A
square-summable boundary-singularity profile on a specified periodic carrier
$T\subseteq P(\psi)$ & It is not the original holomorphic function and need
not determine that function.\\

Singularity Conservation Law & The fixed-point identity
$\mathfrak Q_{H:T}R_{\psi,T}=R_{\psi,T}$ on the maximal graph domain defined
by the raw finite expression $\mathcal E_{H:T}$ & A typed boundary-value
intertwining identity & One-sided branch invariance makes the raw expression
defined; it does not prove that $\mathfrak Q_{H:T}$ is an everywhere-defined
bounded operator on $\ell^2(T/\Z)$.\\

$(p,q)$-adic analysis & Functions on $\Z_p$ with values in a complete field
of residue characteristic $q\ne p$ & Harmonic analysis with different
domain and codomain places & It is not analysis over a single object carrying
simultaneously the $p$- and $q$-adic topologies.\\

Rising-continuity &
$f(z)=\lim_N f([z]_{p^N})$; on the bounded subspace
$\widetilde C_b(\Z_p,K)$ the del norm equals the supremum norm & Sequential
continuity along the one canonical truncation path at each point, with an
exact Banach subspace after boundedness is imposed & It is weaker than
topological continuity and does not itself imply boundedness.\\

Frame & A point-indexed choice $z\mapsto\mathcal F(z)$ of a complete valued
extension in which to read a limit & A field of pointwise convergence
topologies, or a bundle of chosen completions & It is not one normed target
space unless additional uniform structure is imposed.\\

M-function &
$M_N(z)=\prod_{h<N}r_{d_h(z)}$ & A multiplicative digit cocycle or prefix
product & Decay of $M_N(z)$ does not by itself make its norm sequence
monotone.\\

Quasi-integrability & In one fixed complete field $K$, pointwise limits
$A_{\mathbf m}(t)=\lim_N\widehat X_{\mathbf m,N}(t)$; each resulting function
$A:\Gamma_p\to K$ acts on locally constant tests by the finite sum
\cref{eq:formal-Fourier-distribution} & An algebraic distributional
realization of specified limiting Fourier data & A merely formal solution of
the Fourier recursion does not establish the required finite-transform
limits.\\

Breakdown locus & The resonance equation
$\alpha_{\mathbf n}(0)=1$ & A zero-mode resonance locus for the formal
Fourier recurrence & It does not mean automatic nonexistence: the lower-term
compatibility sum decides between no solution and an affine family.\\

Algebra of distributions & The declared rule
$(X^m dz)(X^n dz):=X^{m+n}dz$ on designated symbols, realized by an
$R$-linear map $\Phi:R[T]\to\mathcal S(\Z_p,R)'$ & A polynomial algebra
transported from powers of the named F-series & Multiplication descends to
$\operatorname{im}\Phi$ exactly when $\ker\Phi$ is an ideal; even then it is
not an intrinsic product of arbitrary distributions or ordinary convergent
Fourier convolution.\\
\bottomrule
\end{longtable}
\endgroup

\subsection{The typed dependency and result map}

There are four principal chains.  Each arrow below has the hypotheses shown
in its row; no arrow is supplied merely by similarity of notation.

\begingroup\small
\setlength{\tabcolsep}{3pt}
\begin{longtable}{@{}p{0.18\textwidth}p{0.34\textwidth}p{0.39\textwidth}@{}}
\toprule
Chain & Input and indispensable hypothesis & Output, and the next permitted
step\\
\midrule
\endfirsthead
\toprule
Chain & Input and indispensable hypothesis & Output, and the next permitted
step\\
\midrule
\endhead
Number-field state space & A number field $K$, its full Minkowski map
$\iota_K:K\to\R^{r_1}\times\C^{r_2}$, and an integral basis of $\OK$ &
$\iota_K(\OK)$ is a full lattice and hence locally finite; the
preperiodic-or-escaping dichotomy then applies to self-maps of this full
image.  A single archimedean projection is not substituted; see
\cref{prop:minkowski-lattice,prop:preperiodic-or-escaping}.\\

Digit compactum & An integer $p\geq2$, the discrete alphabet
$\mathcal A_p$, and the compatible residue system
$\varprojlim_N\Z/p^N\Z$ & The digit map is a homeomorphism
$\mathcal A_p^{\N_0}\simeq\Z_p^{\mathrm{dig}}$, and Haar probability is the
uniform product law with cylinder mass $p^{-N}$; see
\cref{prop:digit-compactum-haar}.\\

Affine foundation & Admissible branches
$H_j=(r_j,c_j)$ and the affine semidirect-product law & The piecewise map
$H$ and the exact decomposition
$H_{\mathbf j}(x)=M_H(\mathbf j)x+X_H(\mathbf j)$; see
\cref{prop:affine-concatenation}.\\

Scalar prior-art boundary & $K=\F=\Q$, $\OK=\Z$, standard classes modulo a
displayed $D\in\Z_{\geq2}$, and rational-affine branches integral on their assigned
classes & Unique integers $m_i\ne0$ and $\rho_i\equiv im_i\pmod D$ give
$H_i(x)=(m_ix-\rho_i)/D$; the resulting map is rcwa and
$d_{\mathrm{Mat}}(H)=\operatorname{Mod}(H)\mid D$.  Exact-integrality is
equivalent to $\gcd(m_0\cdots m_{D-1},D)=1$ for the displayed presentation;
it forces both least moduli to equal $D$ and hence makes the map Matthews
relatively-prime.  Conversely, Matthews relatively prime type gives this
exact-integral presentation only when $D=d_{\mathrm{Mat}}(H)$.
Membership in Kohl's group
$\operatorname{RCWA}(\Z)$ requires a separate global bijectivity proof; see
\cref{prop:scalar-hydra-matthews,prop:scalar-hydra-rcwa}.\\

Affine recentering & An integral translation $\tau_a$ and the induced
quotient permutation $\Lambda_{\pi_a(j)}=\Lambda_j+a$ & The conjugated branch
on $\Lambda_j$ is $\tau_a^{-1}H_{\pi_a(j)}\tau_a$; its zero-class branch is
centered exactly when $H_{\pi_a(0)}(a)=a$.  A fixed point of $H_0$ suffices
without another branch equation only when $a\in\Lambda$; see
\cref{prop:translation-conjugacy-cosets}.  Properness after translation is
likewise the exact condition that $I-r_{\pi_a(0)}$ be invertible.\\

 Finite numen & A specified terminal value $a$ satisfying
 $(I-r_0)a=c_0$; for a centered proper Hydra this forces $a=0$ & A unique
 digit-recursive function on $\N_0$ with $f(0)=a$ and
 $f(pn+j)=r_j f(n)+c_j$; see
 \cref{thm:digit-recursive-characterization}.\\

Infinite numen & A named completion $K_\ell$ and pointwise prefix-product
decay, almost-everywhere geometric-mean contraction, or uniform contraction
& A rising, measurable, or continuous extension, respectively, satisfying
the same digit equations; see
\cref{prop:numen-pointwise-convergence,thm:numen-ae-extension,thm:numen-uniform-contraction}.\\

Rational digit input & The canonical CRT identification
$\Z_p^{\mathrm{dig}}\simeq\prod_{\ell\mid p}\Z_\ell$ and the diagonal copy of
$\Q$ & The intersection is exactly
$\Z_{(p)}=\{a/b:\gcd(b,p)=1\}$, and these are exactly the eventually periodic
base-$p$ digit strings; see \cref{prop:rational-eventually-periodic}.\\

Periodic code & A repeated word $\mathbf j$ and invertibility of
$I-M_H(\mathbf j)$ (obtained, for example, from
$\|M_H(\mathbf j)\|_\ell<1$) & The unique \emph{formal} fixed point
$(I-M_H(\mathbf j))^{-1}X_H(\mathbf j)$; see
\cref{thm:repeated-code-fixed-point}.\\

Classical conjugacy comparison & An odd integer $q$, the chronological
parity encoder $Q_q$, and a length-$L$ parity block with code $e$ and reversed
code $\operatorname{rev}_L(e)$ & The typed isometric conjugacy
$Q_qT_q=SQ_q$ for every odd $q$ and the common rational representative
$R_{q,\boldsymbol\varepsilon}$ of \cref{eq:BL-periodic-rational} under the
separate embeddings; when $q$ is prime, its cross-adic numen evaluation is
\cref{eq:BL-numen-overlap}.  The latter equality is neither a global map
identity nor a composite-$q$ extension of $\chi_q$.\\

Genuine periodicity & The preceding formal fixed point lies in the integral
state space, and branch correctness is proved & A periodic point of the
piecewise map.  This implication is proved for $T_q$ by
\cref{prop:Tq-correct-fixed-point}; the general 2026 implication remains at
source-sketch status.\\

Dreamcatcher input & A function
$\psi\in\mathcal D_1(T)\cap\operatorname{ran}(\mathcal G_{\mathbb H})$, a
branch-invariant carrier $T$, and the coefficient pullback
$\mathscr Q_H$ on its declared range & The Virtual Residue Formula
\cref{eq:VRF}; cross-section densities provide an independent exact Abelian
criterion for its radial limits by \cref{prop:cross-section-vres}.\\

Dreamcatcher conservation & The preceding data and a verified
$R_{\psi,T}$-membership in the maximal graph domain of
$\mathfrak Q_{H:T}$ & The fixed-point identity
$\mathfrak Q_{H:T}R_{\psi,T}=R_{\psi,T}$ in \cref{thm:SCL}; no all-space
boundedness is inferred.\\

Finite rational obstruction & Prime and regulated 2019 Hydra plus finite
support modulo $\Z$ & The Finite Dreamcatcher conclusion
$R=R(0)\ind{\Z}$; rational invariant sets are then finite or cofinite.\\

Contour coefficient & The exact coefficients
$a_n(x)=(1-M_q(n))x-\chi_q(n)$ and a line
$c>\sigma_q$ of absolute convergence & Perron inversion recovers one
coefficient by \cref{eq:coefficient-contour}; the integral vanishes exactly
when that particular word numerator vanishes.\\

Dirichlet recursion & The coefficient recursions
\cref{eq:Fq-coefficient-recursion} and $\Re s>\sigma_q$ & The corrected
ordinary and Dirichlet generating-function identities
\cref{eq:fq-corrected-functional-equation,eq:Fq-functional-equation}.\\

Dirichlet continuation & The corrected shift equation and the normal-tail
estimate \cref{lem:Fq-normal-shifted-tail} & Unique meromorphic
continuation to $\C$, the candidate-pole containment
\cref{eq:Fq-candidate-poles}, and the corrected $q=3$ principal parts.
No actual pole is inferred merely from a denominator zero.\\

Contour dynamics & A synchronized block length and the $T_q$ branch-
correctness theorem & A phase-sensitive periodicity statement using
$B_{2,L}(n)$; shortest-word data alone do not encode every phase.\\

Function Fourier theory & $q\ne p$, a specified tower
$\zeta_{N+1}^p=\zeta_N$ of primitive roots, normalized $K$-valued Haar
integration, and a complete algebraically closed nonarchimedean target & The
explicit pairing $\Gamma_p\times\Z_p\to K^\times$, classification of locally
constant $K$-valued characters, both finite orthogonality identities, the
isometry $C(\Z_p,K)\simeq c_0(\Gamma_p,K)$, and the precisely signed
product--convolution identities; see
\cref{prop:K-valued-character-duality,prop:finite-character-orthogonality,%
thm:pq-Fourier-isometry,prop:typed-convolution-identities}.\\

Function cyclicity & The preceding Banach-algebra isomorphism & The four-way
Wiener--Tauberian equivalence in \cref{thm:pq-WTT-function}.  Its dense
finite translate span is
$\widehat\chi*c_{00}(\Gamma_p,K)$; only after closure does it equal the
closed principal ideal
$\overline{\widehat\chi*c_0(\Gamma_p,K)}$.\\

Measure cyclicity & A bounded Fourier--Stieltjes transform and existing local
density limits & Density of translates implies nonvanishing of each existing
limit.  The reverse arrow is not proved.\\

Rising Banach domain & A complete nonarchimedean target $K$ and the bounded
rising-continuous class $\widetilde C_b(\Z_p,K)$ & The del norm equals the
supremum norm and makes this exact subspace Banach; unrestricted
rising-continuity still permits infinite del quantity; see
\cref{prop:del-norm-domain}.\\

Product algebra, finite level & Degree-one affine digit recursions in a common
commutative coefficient ring; for analytic transforms, one fixed
nonarchimedean field $K$ and continuous $K$-valued products & The
multinomial branch identity \cref{eq:product-branch-recursion}, the continuous
Fourier recursion, the universally exact truncation recursion, and the
zero-mode compatibility classification \cref{prop:formal-compatibility}.\\

Product distributions & For every $\mathbf m\leq\mathbf n$ and
$t\in\Gamma_p$, convergence of
$\widehat X_{\mathbf m,N}(t)$ in one fixed complete field $K$ & The limiting
functions satisfy the exact Fourier recursion and define unique algebraic
distributions on $\mathcal S(\Z_p,K)$; the truncation distributions converge
pointwise on that test space by \cref{prop:finite-transform-limit}.  The
printed Main Limit sufficient conditions do not establish this hypothesis.\\
\bottomrule
\end{longtable}
\endgroup

In compact form, the verified central spine is
\[
 \{H_j=(r_j,c_j)\}
 \longrightarrow (M_H,X_H)
 \longrightarrow X_H\!\mid_{\N_0}
 \longrightarrow X_{H,\ell}
 \longrightarrow (I-M_H(\mathbf j))^{-1}X_H(\mathbf j),
\]
where the fourth arrow requires a place and convergence, and the last value
becomes a periodic point only after branch correctness.  The three analytic
branches are
\[
 \psi_V\longrightarrow\operatorname{VRes}\psi_V
 \longrightarrow\mathfrak Q_{H:T}R,
 \qquad
 (M_q,\chi_q)\longrightarrow F_q\longrightarrow a_n(x),
\]
and
\[
 C(\Z_p,K)\xrightarrow{\ \mathscr F\ }c_0(\Gamma_p,K),
 \qquad
 (X_j)_j\longrightarrow X_{\mathbf n,N}\longrightarrow
 \widehat X_{\mathbf n,N}.
\]
Within these three displayed analytic branches, the unresolved steps are the
everywhere-defined boundedness of the dreamcatcher operator, the
shortest-word contour biconditional, the converse measure cyclicity criterion,
and source-sufficient hypotheses for the general product limit.  The exact
maximal graph domain and the exact consequence of pointwise finite-transform
convergence are proved.  The separate source-status table below records
additional unproved or false claims outside those displayed branches.

\subsection{Result-status table}

The five codes below are mutually exclusive as editorial dispositions.
``Proof checked'' means that the source proof and its named dependencies were
read against the typed statement used here; it does not mean computer formal
verification.  When this edition gives a complete proof or changes a
hypothesis and proves the changed statement, the stronger code
\textbf{R} is used instead.

\begin{description}[leftmargin=2.8em,labelwidth=1.8em,style=multiline]
\item[\textbf{R}] reconstructed and proved in this edition;
\item[\textbf{S}] source theorem with its proof checked against the typed
      hypotheses;
\item[\textbf{C}] conditional, conjectural, or formal statement not promoted
      to an unconditional theorem;
\item[\textbf{U}] superseded witness retained for mathematical history;
\item[\textbf{W}] withdrawn witness or statement invalid as printed.
\end{description}

\begingroup\small
\setlength{\tabcolsep}{3pt}
\begin{longtable}{@{}p{0.07\textwidth}p{0.27\textwidth}p{0.18\textwidth}p{0.39\textwidth}@{}}
\toprule
Code & Result or witness & Source & Exact disposition\\
\midrule
\endfirsthead
\toprule
Code & Result or witness & Source & Exact disposition\\
\midrule
\endhead
\textbf{R} & Preperiodic-or-escaping dichotomy & 2601.17030v3 &
Reconstructed and proved in \cref{prop:preperiodic-or-escaping} with the
necessary locally finite hypothesis replacing bare uniform discreteness.\\

\textbf{R} & Full Minkowski lattice foundation & Standard number-field
foundation supplied by this edition & The full map, its domain and codomain,
the integral-basis matrix, nonzero discriminant, and local finiteness are
proved in \cref{prop:minkowski-lattice}; no one-embedding substitute is used.\\

\textbf{R} & Digit compactum and Haar product law & 2412.02902v1 and
2601.17030v3, rebuilt & The explicit map from digit sequences to the inverse
limit, its inverse, homeomorphism property, and every cylinder mass are
proved in \cref{prop:digit-compactum-haar}, including composite $p$.\\

\textbf{R} & Rational points and eventually periodic digits & Edition-supplied
arithmetic foundation & The compatible CRT isomorphism, diagonal
intersection $\Q\cap\Z_p^{\mathrm{dig}}=\Z_{(p)}$, and both directions of
eventual periodicity are proved in
\cref{prop:rational-eventually-periodic}.\\

\textbf{R} & Affine word calculus and repeated-word fixed point &
2601.17030v3, reorganized & Proved in
\cref{prop:affine-concatenation,prop:repeated-word}; endomorphism order and
invertibility are explicit.\\

\textbf{R} & Exact scalar prior-art identification & Matthews's generalized
Collatz form and Kohl's rcwa definition & The fixed-displayed-modulus
equivalence, the typed inclusion in the rcwa mapping monoid, divisibility of
the displayed modulus by the minimal modulus
$d_{\mathrm{Mat}}(H)=\operatorname{Mod}(H)$, the exact equivalence between
scalar exact-integrality and the conjunction of Matthews relatively-prime type
with a displayed least modulus, equality of the minimal modulus with the
displayed modulus under exact-integrality, and the
separate bijectivity condition for
$\operatorname{RCWA}(\Z)$ are proved in
\cref{prop:scalar-hydra-matthews,prop:scalar-hydra-rcwa}.\\

\textbf{R} & Translation conjugacy and zero-class centering &
2601.17030v3, corrected & The induced residue-class permutation, every
conjugated branch, preservation of admissibility and exact-integrality, and
the necessary-and-sufficient centering equation are proved in
\cref{prop:translation-conjugacy-cosets}.\\

\textbf{R} & Digit recursion and the finite numen & 2007.15936v3,
2412.02902v1, 2601.17030v3 & Reconstructed and proved in
\cref{thm:digit-recursive-characterization}, including the noncentered
terminal-value alternative.\\

\textbf{R} & Pointwise, almost-everywhere, and uniform rising extension &
2412.02902v1, 2601.17030v3 & Proved under the displayed prefix-product,
geometric-mean, and uniform contraction hypotheses in
\cref{prop:numen-pointwise-convergence,thm:numen-ae-extension,thm:numen-uniform-contraction}.\\

\textbf{R} & Phase-sensitive $T_q$ correspondence & 2007.15936v3 & The
affine identity, $2$-adic branch correctness, and padded length-$L$ converse
are proved in \cref{thm:Tq-periodic-word-CP}.\\

\textbf{R} & Classical conjugacy and reversed periodic overlap &
B\"ohm--Sontacchi (1978), Bernstein--Lagarias (1996), and
Laarhoven--de Weger (2012) & The parity map and inverse conjugacy are typed
and proved for every odd $q$ in \cref{prop:BL-conjugacy-typed}.  The two
adic values are proved to be the images of one rational representative under
their separate canonical embeddings.  The mandatory reversal, the $q=3$
specialization to the B\"ohm--Sontacchi Proposition~5 candidate, its
period-dividing boundary, and global nonidentity with $\chi_q$ are proved in
\cref{thm:BL-numen-periodic-overlap,cor:BL-numen-global-nonidentity}.\\

\textbf{R} & Square Sum Lemma, Virtual Residue Formula, and typed SCL &
1909.09733v3 & The square-sum proof-idea attribution is bound to accepted
MathOverflow answer 333801 by user \emph{fedja} \cite{FedjaMO333801}; this
edition's finite-subset Bessel proof remains editorial and is not attributed
to the answerer.  VRF/SCL are proved with the operator-domain premise stated.\\

\textbf{R} & Exact coefficient pullback & 1909.09733v3, rebuilt & The map
$\mathsf P_H:\ell^\infty\to\ell^\infty$ is proved to be an isometric
injection from surjectivity, and the two generating-series operators are
defined only as its conjugates on their declared ranges.\\

\textbf{R} & Cross-section densities and virtual residues & 1909.09733v3,
rebuilt & The residue-class Ces\`aro limits, summation-by-parts Abel passage,
and exact factor $1/(2\pi)$ are proved in
\cref{prop:cross-section-vres}; no Tauberian converse is used.\\

\textbf{R} & Dyadic sums and coefficient contour identity &
2111.07883v2 & Recalculated and proved in
\cref{eq:chi-block-sum,eq:M-block-sum,thm:coefficient-contour}.\\

\textbf{R} & Nonarchimedean Fourier isometry and function
Wiener--Tauberian theorem & 2208.11082v2 & Reconstructed from finite Fourier
duality.  The compatible-root pairing, character classification, both finite
orthogonality formulas, the isometry, the typed convolution morphisms, and
the cyclicity theorem are proved in
\cref{prop:K-valued-character-duality,prop:finite-character-orthogonality,%
thm:pq-Fourier-isometry,prop:typed-convolution-identities,%
thm:pq-WTT-function}.\\

\textbf{R} & Rising van der Put theorem and unique rising continuation &
2412.02902v1 & Proved by the exact telescoping path argument in
\cref{thm:rising-vdp,cor:unique-rising-continuation}.\\

\textbf{R} & Del norm on the bounded rising class & 2412.02902v1, corrected &
The level operators are defined for every $N$, the extended del quantity is
proved equal to the supremum quantity on rising-continuous functions, and
$\widetilde C_b(\Z_p,K)$ is proved complete in
\cref{prop:del-norm-domain}.\\

\textbf{R} & Product branch recursion and formal compatibility trichotomy &
2507.13358v1 & Proved at the finite/formal level in
\cref{prop:product-branch-recursion,prop:formal-compatibility}; no limiting
distribution is inferred.\\

\textbf{R} & Fixed-target F-series domains and the $\chi_3$ frame identity &
2307.00436v1 and 2412.02902v1, rebuilt & Padded digit counts give the exact
convergence domains $D_E$ and $D_X$, every forward/converse domain
implication, and the scalar and affine branch equations in one complete
field.  The two-place $\chi_3$ identity follows instead from the finite
telescoping formula \cref{eq:chi3-finite-frame-identity}; no disputed
pointwise-necessity argument is used.\\

\textbf{R} & Pointwise finite-transform limits and algebraic distributions &
2507.13358v1, reconstructed & Under the exact fixed-$K$ pointwise convergence
hypothesis, the limiting Fourier recursion, unique functional on locally
constant tests, and pointwise convergence of the truncation distributions are
proved in \cref{prop:finite-transform-limit}.\\

\textbf{S} & Finite Dreamcatcher Theorem & 1909.09733v3 & Source theorem;
the Off-$\varrho$ theorem, auxiliary functional equation, support lemmas,
Off-$\varrho$ Infinity theorem, and final dependency chain were checked.  It
is retained with the source-proof status recorded in \cref{status:FDT}.\\

\textbf{S} & Weak Rationality Theorem 1 & 1909.09733v3 & Source consequence;
its dependency on rational generating functions, SCL, finite virtual-pole
support, and the Finite Dreamcatcher Theorem was checked.  This disposition
does not extend to Weak Rationality Theorem 2.\\

\textbf{R} & Corrected coefficient, ordinary generating-function, and
Dirichlet functional equations for $F_q$ & 2111.07883v2, corrected & The
printed ordinary generating-function derivation has a wrong first even
term.  The recursions and both corrected identities are proved directly in
\cref{prop:Fq-coefficient-recursion,thm:Fq-corrected-functional-equation}
on the proved initial absolute-convergence half-plane
$\Re s>\sigma_q$.\\

\textbf{R} & Global continuation and corrected $q=3$ principal parts &
2111.07883v2 and Web II, rebuilt & Normal convergence and strip induction
prove \cref{thm:Fq-global-meromorphic-continuation}; direct Laurent
calculation proves \cref{prop:F3-corrected-principal-parts}.  The pole set
is contained in the stated candidate set; equality is not asserted.\\

\textbf{C} & General global-field Correspondence Principle and divergence
clause & 2601.17030v3 & Conditional/source formulation only.  The source
labels the proof a sketch, and the zero fibre, correctness, and irrational-
fibre steps are not established by its displayed hypotheses.\\

\textbf{C} & Theorem 12 rational-edgepoint conclusion & 1909.09733v3 & May
be cited only with branch invariance of both rational edgepoints and rational
non-edgepoints; the source itself identifies this as the method's Achilles
heel.\\

\textbf{C} & Identification of a limiting dreamcatcher with the virtual
residue of the limiting function & 1909.09733v3 & Explicitly Conjecture 1,
not a theorem.\\

\textbf{W} & Printed $F_q$ functional equation and its dependent
singularity chain & 2111.07883v2 & The printed identity fails for every
$x\ne0$, as certified by the $s\to+\infty$ limit in Chapter 7.  Its
residue formulas contain spurious $x$-terms, and its actual-pole assertions
lack noncancellation proofs.  The continuation theorem and every retained
principal part are instead re-derived from the corrected equation.\\

\textbf{C} & Self-similar numen measure uniqueness & 2601.17030v3 &
Existence of the pushforward and its functional equation are
typed; uniqueness requires a separate contraction or IFS hypothesis.\\

\textbf{C} & Main Limit, product distributions, power theorem, and inversion
& 2507.13358v1 & Formal or conditional.  Exact pointwise convergence of every
required finite transform in one fixed $K$ would produce the limiting
recursion and algebraic distributions by
\cref{prop:finite-transform-limit}; the printed sufficient-condition proof
and induction do not establish that convergence.\\

\textbf{C} & Declared product/convolution algebra of named distributions &
2507.13358v1 & The symbolic polynomial algebra exists formally.  After the
designated distributions exist, its multiplication descends to their realized
image if and only if the realization kernel is an ideal.  It is not ordinary
distribution multiplication or convergent convolution.\\

\textbf{U} & Parity-vector Fourier predecessor & 2002.12762v1 & Preserved as
a historical stage but superseded by the numen formulations; v2 is the
withdrawal notice declaring the work obsolete.\\

\textbf{U} & Original 2111.07883v1 & 2111.07883v1 & Materially different
title and subject; superseded for the logical work by the current contour
paper v2, not silently merged with it.\\

\textbf{U} & Earlier general Hydra conventions & pre-2601 corpus & The 2026
manual v3 controls the general name and word formalism.  The specialized
2019 $\varrho$-Hydra remains the correct historical object for the
dreamcatcher theorems.\\

\textbf{U} & Earlier witnesses of the 2026 manual & 2601.17030v1--v2 & v3 is
the controlling witness and includes the correction to Proposition 7.6;
earlier text is retained only for version history.\\

\textbf{W} & Number-ring functional-equation paper & 2111.07882 v1/v2 & v2
withdraws the work for an error in the quotient-ring section.  It is not used
to certify a general number-ring theorem.\\

\textbf{W} & Weak Rationality Theorem 2 proof & 1909.09733v3 & Invalid as
printed: naturally measurable sets do not form a $\sigma$-algebra and
natural density is not a countably additive measure on such a class.\\

\textbf{W} & Shortest-word contour biconditional for every listed cycle
phase & 2111.07883v2, relying on 2007.15936 & Invalid as printed; the point
$2$ of the $T_3$ cycle already requires a padded block.  The coefficient
identity survives.\\

\textbf{W} & Measure Wiener--Tauberian equivalence & 2208.11082v2 & The
proof gives density $\Rightarrow$ no existing zero local density and then
repeats its contrapositive; the converse is not proved and can be vacuous.\\

\textbf{W} & Banach space of all compatible frame functions & 2307.00436v1
& Invalid without boundedness or continuity: compactness of the domain does
not bound an arbitrary function.\\

\textbf{W} & Banach algebra of all rising-continuous functions &
2412.02902v1 & Invalid as defined: rising-continuity does not imply
boundedness or finite del norm.  This edition instead introduces and proves
the exact Banach space $\widetilde C_b(\Z_p,K)$ in
\cref{prop:del-norm-domain}; the unrestricted source claim remains invalid.\\

\textbf{W} & Main Limit proof and resulting unconditional product theorems &
2507.13358v1 & Invalid as a proof: decay is repeatedly treated as eventual
monotonicity; root bounds are used to claim a tail is the order of its first
term; and $Na_N=O(a_N)$ is asserted.\\

\textbf{W} & Simultaneous universal and exact integrality clauses &
2601.17030v3 & Inconsistent for $p\geq2$ and incompatible with $T_3$.
Both source clauses quantify over $\mathcal O_K$.
\Cref{eq:branch-admissibility,eq:exact-integrality} preserves that domain and
separates the two mathematically usable predicates.\\

\textbf{W} & Pointwise digit-density convergence criterion &
2601.17030v3 & Invalid as printed: the source uses a limsup for frequencies
of contractive branches and a liminf for expansive branches, reversing the
inequalities needed for an upper Lyapunov bound.  The corrected bound is
\cref{eq:density-lyapunov-upper-bound}.\\

\textbf{W} & Multidimensional repeated-block identity and resulting
Correspondence Principle & 2412.02902v1 & Invalid as printed when coordinate
digit lengths differ.  Repetition of one synchronized vector block of common
length $N$ gives coordinate $n_i/(1-p^N)$, not
$n_i/(1-p^{\lambda_p(n_i)})$.  The multidimensional correspondence built on
the latter identity is not promoted.\\
\bottomrule
\end{longtable}
\endgroup

\subsection{Source-by-source repair ledger}

This ledger gives the minimum repair actually used by the edition.  An entry
of ``not promoted'' means that no repaired theorem is smuggled into a later
argument.

\begingroup\small
\setlength{\tabcolsep}{3pt}
\begin{longtable}{@{}p{0.16\textwidth}p{0.35\textwidth}p{0.41\textwidth}@{}}
\toprule
Witness & Source defect or unstable point & Editorial treatment in this
edition\\
\midrule
\endfirsthead
\toprule
Witness & Source defect or unstable point & Editorial treatment in this
edition\\
\midrule
\endhead
1909.09733v3 & The paper alternates between the defined class
$\mathcal D_1(T)$ and an undefined $\mathcal D_2(T)$; its operator proof
treats one-sided branch inclusion as bijectivity; its virtual residue is not
the ordinary complex residue; the Permutation Operator Theorem says
 ``basis'' without a topology. & Declare $\mathcal D_1(T)$ as the edition's
 explicit source-level choice and intersect it with the declared generating-
operator range wherever $\mathscr Q_H$ is used; define the maximal graph
domain and state its membership premise in SCL; retain the radial
normalization; prove the exact coefficient pullback and Abelian cross-section
residue formula; and replace the Hamel-basis wording by the coefficient-orbit
statement.  Exclude Weak Rationality Theorem 2 and retain the limit
identification as Conjecture 1.\\

2007.15936v3 & The shortest binary word loses a high zero that belongs to a
periodic block; one repeated-block truncation is written modulo $2^m$ rather
than $2^{mL}$.  Its general integer-domain Hydra definition also requires
the selected branch value to be a \emph{nonnegative} integer iff the input
has the assigned residue, which fails for negative inputs (including the
negative fixed point of $T_3$). & Keep the block length in
$B_{2,L}(n)=n/(1-2^L)$, reverse the chronological itinerary explicitly, and
prove correctness in $\Z_2$.  The shortest notation $B_2(n)$ is used only
for a phase whose high block digit is nonzero.  On domain $\Z$, branch
recognition means \emph{integer} iff assigned residue; nonnegativity belongs
only to a separately restricted state space such as $\N_0$.\\

B\"ohm--Sontacchi (1978) & Printed Proposition~4 identifies a formal affine
branch composition with $u^n(x)$ while quantifying independently over every
$x\in\Z$ and every run tuple.  This is false: for
$x=0$, $m=1$, and $n_0=n_1=1$, its right side is $1/4$ whereas
$u^3(0)=0$.  The same right side is correctly
$f\circ g\circ f(0)=1/4$. & Derive the formal-composition identity from
their Proposition~1 and identify it with an iterate only after the branch
word is proved to be the actual parity itinerary.  Recover the
Proposition~5 fixed-return candidate independently from the edition's
branch-correctness theorem; retain the unnumbered paragraph after
Proposition~5 on p.~264 as the exact period-dividing locator; and mark the
last position inequality as a reconstruction from the formal run data, not
as text printed in Proposition~5.\\

2111.07883v2 & The contour theorem inherits the shortest-word phase defect;
its common $c>\sigma_q$ quantifier is not correctly specialized to the
$q=3,5$ parts; the $q\geq7$ excluded set is malformed; one abscissa paragraph
says ``minimum''; the derivation of $f_q(z^2,x)$ uses $z/(1-z)$ in place of
$z^2/(1-z^2)$; and the resulting false Dirichlet identity underlies the
later singularity discussion. & Prove the Perron coefficient identity on
$c>\sigma_q$ and state each specialization with its own line.  Re-derive
the coefficient, ordinary generating-function, and Dirichlet recursions;
prove the corrected identity first on $\Re s>\sigma_q$; prove normal
convergence and continue it strip by strip; compute the corrected $q=3$
principal parts; and treat the remaining lattice as candidates unless
noncancellation is separately proved.\\

2208.11082v2 & The body restatement omits some headline hypotheses; a radius
in the differentiation discussion tends in the wrong direction; and the
measure WTT proves only one implication.  Its notation also asks for an
embedding $\overline{\Q}\hookrightarrow K$, which is unavailable in positive
characteristic even though the stated field class allows that case. & Carry
the full nonarchimedean setting and residue-characteristic restriction;
construct only the required compatible $p$-power roots directly in $K$;
prove representative independence, character classification, and finite
orthogonality; fix $(f*\mu)(z)=\mu(y\mapsto f(z-y))$ so the local-density
coset is exact; use balls shrinking to a point; reconstruct the function WTT
from finite Fourier duality; retain only density $\Rightarrow$ nonzero
existing local densities for measures.\\

2307.00436v1 & The abstract reverses the complex and $q$-adic parts of the
standard frame; one $X_{p,q}$ display swaps $p$ and $q$; a convergence proof
replaces the chosen input by $-1$; and all compatible functions are declared
Banach under a possibly infinite sup norm. & The body frame definition and
the original $S_{p,q}$ summand control; define the exact fixed-target domains
$D_E$ and $D_X$ and test convergence on the actual digit path; prove the
branch equations only on those domains; and derive the two-place $\chi_3$
identity from its finite telescoping formula.  Normed claims are restricted to
a bounded or continuous class.\\

2412.02902v1 & The definition of rising-continuity is pointwise, but a
footnote asserts all such functions are bounded and the later del-norm Banach
algebra uses that assertion.  In the multidimensional block construction,
coordinates of unequal digit lengths are repeated with their individual
denominators even though the vector word has one synchronized common length.
& Retain the pointwise definition and the valid truncation/van der Put
theory.  The prefix-free example
$f(p^k)=q^{-k}$, with $f=0$ elsewhere, is rising-continuous and unbounded in
$\Q_q$.  Define every del operator, prove that the del and supremum quantities
agree on the rising-continuous class, and prove completeness of the exact
bounded subspace $\widetilde C_b(\Z_p,K)$.  For a common block length $N$, use
$n_i/(1-p^N)$ in every coordinate; the printed multidimensional
Correspondence Principle is not retained without a new proof from that
corrected synchronization.\\

2507.13358v1 & Haar mass and the finite character factor are printed as
$1/p$ instead of $p^{-N}$; the archimedean Wiener norm is printed as a
coefficient supremum instead of $\ell^1$; one finite recursion double-counts
the diagonal; the Main Limit proof uses false monotonicity and tail
estimates; localization and fraction-field hypotheses are insufficient; the
paper's own warning says ordinary convolution usually diverges. & Correct
the finite normalizations and use $\mathbf m<\mathbf n$ after separating the
diagonal.  Retain the multinomial and formal compatibility algebra.  Prove
the exact fixed-$K$ consequence of pointwise finite-transform convergence,
including the algebraic distribution and pointwise test-space limit.  Require
ideals disjoint from the localization set and a domain before taking a
fraction field.  Classify the source's limiting product/inversion results as
conditional, and make the named transported multiplication descend to a
realized image only under the exact kernel-ideal condition.\\

2601.17030v3 & The Hydra definition simultaneously imposes universal branch
integrality and an incompatible iff branch condition; properness is used as
though it implied centeredness; the source's ``affine lattice morphisms'' are
typed as all affine $\F$-linear automorphisms and need not preserve a lattice;
scalar division notation is used for an endomorphism; bare uniform
discreteness is used as local finiteness; the
pointwise density criterion reverses liminf and limsup; and the
Correspondence proof is a sketch with an unlisted hypothesis, an unresolved
zero orbit, and an unsupported fibre inference.  Its preliminary
ultrametric proposition says equality holds iff the summands have unequal
norms; unequal norms are sufficient, not necessary. & Separate admissibility
from exact-integrality; call $\Aff_\F(K)$ the affine automorphism group and
state lattice preservation separately; track the quotient permutation under
translation and use $H_{\pi_a(0)}(a)=a$ as the exact centering condition;
retain the zero-branch fixed point in the noncentered recursion; use
$(I-M)^{-1}$; impose local finiteness for the escape dichotomy; use the
corrected Lyapunov upper bound; and prove only the affine repeated-word
calculation.  The false biconditional for ultrametric equality is not used.
Record the general Correspondence Principle and divergence statement at
source-sketch status; do not infer self-similar-measure uniqueness without
contraction.\\
\bottomrule
\end{longtable}
\endgroup

The two withdrawn logical works require no mathematical repair inside this
volume.  The witness arXiv:2002.12762v1 is retained only as a superseded
predecessor and its v2 record is the withdrawal notice.  The witness
arXiv:2111.07882 is retained only as
a historical record of the number-ring programme; its quotient-ring theorem
is not imported through another paper.

\subsection{Exact coverage and explicit nonclaims}

The edition provides exact mathematical coverage of the following objects.
\begin{enumerate}[label=\textup{(\arabic*)}]
\item The full Minkowski embedding of a number-field integer ring as a locally
      finite lattice and the resulting exact preperiodic-or-escaping
      dichotomy.
\item Residue-class affine dynamics, affine-automorphism typing, the
      translation-induced coset permutation and exact centering condition,
      the distinction between admissibility and branch recognition, affine
      word composition, and the ordered cocycles $M_H$ and $X_H$; together
      with the exact scalar fixed-modulus equivalence to Matthews's form and
      the typed inclusion into Kohl's rcwa mapping monoid.
\item The base-$p$ digit compactum as an explicit inverse limit and product
      probability space, including composite $p$; its compatible CRT
      decomposition; and the exact equivalence between diagonal rational
      points and eventually periodic digit strings.
\item The finite digit recursion for the numen, its rising extension in a
      specified completion, and pointwise, almost-everywhere, and uniform
      convergence criteria with their hypotheses displayed.
\item The repeated-code fixed-point identity, the complete branch-correctness
      argument for shortened $T_q$, and the phase-sensitive padded-word
      periodic correspondence; the complete classical parity
      encoder/conjugacy typing; the valid formal-composition identity behind
      B\"ohm--Sontacchi's printed Proposition~4; the independently recovered
      Proposition~5 return candidate; and the reversed purely periodic common
      rational representative under the separate Bernstein--Lagarias and
      numen embeddings.
\item Set-series, the exact isometric coefficient pullback and its declared
      generating-map ranges, edgepoints, the Abelian cross-section density
      formula, the Square Sum Lemma, the Virtual Residue Formula, the maximal
      graph-domain Singularity Conservation Law, and the source-status Finite
      Dreamcatcher Theorem.
\item The Dirichlet series $\zeta_{M_q}$, $\zeta_{\chi_q}$, and $F_q$; exact
      dyadic block sums; Perron extraction of one coefficient; the exact
      coefficient recursions; and the corrected ordinary and Dirichlet
      functional equations in $\Re s>\sigma_q$; unique meromorphic
      continuation to $\C$; a rigorous candidate-pole set; and corrected
      $q=3$ principal parts and negative-integer residue recurrence.
\item The explicit compatible-root $K$-valued character pairing on $\Z_p$,
      classification and finite orthogonality, normalized nonarchimedean Haar
      and Fourier analysis, the isometric Fourier transform, precisely signed
      function/sequence/function--measure convolution, and the function
      Wiener--Tauberian equivalence, together with the one proved direction
      of the measure claim.
\item Rising-continuity, canonical truncations, the exact del operators and
      bounded Banach subspace, van der Put path expansions, point-selected
      frames, fixed-target digit-multiplicative and affine F-series, and the
      finite proof of the two-place $\chi_3$ identity.
\item Multinomial product recursions over a common coefficient ring, exact
      continuous and finite Fourier recursions, the formal compatibility
      trichotomy at a breakdown index, the exact distributional consequence
      of fixed-field finite-transform convergence, and the kernel-ideal
      condition for a transported polynomial product to descend.
\end{enumerate}

It makes none of the following claims.
\begin{enumerate}[label=\textup{(N\arabic*)}]
\item It does not prove the full Collatz conjecture, the weak Collatz
      conjecture, or a classification of all integer cycles.
\item It does not claim that the scalar integer Hydra is a new class distinct
      from Matthews's generalized Collatz presentation or Kohl's rcwa mapping
      monoid.  It does not identify an rcwa mapping with an rcwa permutation,
      equate a displayed Hydra modulus with the minimal rcwa modulus in
      general or without the scalar exact-integrality hypothesis, or infer
      Hydra-specific predicates from rcwa membership.
\item It does not identify $Q_q$, $\Phi_q$, and $\chi_q$ as global maps,
      suppress the mandatory chronological-word reversal, promote a return
      under $T_q^L$ to exact period $L$, or use Bernstein--Lagarias's
      composite-odd-$q$ conjugacy to enlarge the prime-$q$ codomain theorem
      for $\chi_q$.  It does not use B\"ohm--Sontacchi's false unrestricted
      Proposition~4 identity to replace a formal branch word by $u^L$ without
      first proving that the word is the actual parity itinerary.
\item It does not claim that every phase of a cycle is represented by the
      shortest positive digit word, or that a periodic-point theorem rules
      out divergent positive trajectories.
\item It does not prove the full 2026 global-field Correspondence Principle,
      its divergence clause, general injectivity of a numen, or a theorem
      identifying every point of a numen fibre.
\item It does not claim that the 2019 natural-density argument is a measure-
      theoretic proof, that rational edgepoint carriers are automatically
      branch invariant, or that a limit of dreamcatchers is automatically the
      dreamcatcher of the limiting holomorphic function.
\item It does not retain any continuation, residue, or pole conclusion
      merely because it was derived from the source's false printed $F_q$
      equation.  The newly proved global continuation does not imply that
      every candidate lattice point is an actual pole, does not establish a
      complete pole classification or the source's asserted
      hyperexponential growth, and does not justify shifting the Perron
      contour left of the absolute-convergence half-plane.
\item It does not prove the converse Wiener--Tauberian theorem for arbitrary
      measures, nor uniqueness of the self-similar numen measure without an
      independent regularity or contraction hypothesis.
\item It does not treat every rising-continuous or frame-compatible function
      as bounded, and it does not assert the dissertation's Banach-algebra
      theorem on the unrestricted class.
\item It does not define a general product of distributions, identify the
      formal convolution powers with ordinarily convergent convolution, or
      promote the 2025 Main Limit, product, measure, and inversion conclusions
      beyond their conditional status.  The edition proves what follows from
      actual pointwise finite-transform convergence; it does not claim the
      source's defective sufficient-condition argument establishes that
      premise.
\item It does not promote the dissertation's multidimensional
      Correspondence Principle: its repeated vector block uses
      coordinatewise digit lengths where one common synchronized block length
      is required.
\item It does not revive the withdrawn quotient-ring paper or merge materially
      different arXiv versions into one theorem.
\end{enumerate}

These exclusions are part of the mathematical result.  They identify the
exact frontier of the reconstructed corpus: the affine and finite-level
structures are strong and explicit; several infinite-level analytic
passages remain conditional; and the celebrated dynamical conjecture is a
target of the programme, not a conclusion of this edition.

\section{Canonical source register}
\label{sec:source-register}

This register fixes the bibliographic identity, version policy, and evidential
role of every source used to define Siegel's canonical research corpus.  A
``lineage'' below is an arXiv identifier, not a claim that every version under
that identifier is the same mathematical text.  Likewise, an author-site
``last modified'' date is provenance metadata, not a rule of mathematical
precedence.

\subsection{Authority order}

Precedence is decided claim by claim, in the following order.

\begin{enumerate}[label=\textup{(\arabic*)}]
\item An explicit authorial correction or forward convention controls the
      object it names.  Thus 2601.17030v3 controls the definitions of Hydra
      maps and numens: it warns that earlier definitions, including the
      dissertation's, may be inconsistent and declares its definitions the
      standard going forward.  For frames, F-series, and rising-continuity,
      2307.00436v1 explicitly takes precedence where it conflicts with the
      dissertation.
\item Subject to (1), the latest substantive, non-withdrawn arXiv version
      controls the theorem stated by that logical work.  A metadata-only
      withdrawal revision contains no replacement mathematics.
\item The dissertation supplies the largest continuous exposition and all
      material not replaced by a later specific source.  Its three public
      witnesses are collated rather than silently identified.
\item A journal-layout PDF is a same-work witness.  It is not presumed to be
      textually identical to the arXiv source, and typography alone does not
      make it a later mathematical correction.
\item An author-site post controls its own expository wording and any material
      unique to that post.  It does not supersede an arXiv theorem merely
      because the page reports a later modification date.
\item A materially distinct earlier version remains part of the critical
      corpus.  A withdrawn work remains part of the historical corpus, but it
      cannot certify a current theorem.
\end{enumerate}

This order fixes source selection; it does not certify truth.  A controlling
witness may still contain a typographical error, an omitted hypothesis, or a
proof sketch, each of which is recorded without alteration.

\subsection{The ten arXiv lineages}

Submission dates below are UTC dates in the archived arXiv metadata.  The
similarity figures are derived TeX-line routing metrics, not altered source
witnesses.  The reproducible procedure is the following.  For the first TeX
file in each version profile, 
\path{scripts/profile_version_deltas.py} reads UTF-8 text with replacement of
decoding errors, splits it into lines, deletes the suffix beginning at an
unescaped percent sign, replaces each nonempty whitespace run by one space,
trims the result, and discards empty lines.  If the resulting line sequences
have lengths $N_1,N_2$ and the matching blocks returned by Python's
\texttt{difflib.SequenceMatcher} with \texttt{autojunk=False} contain $M$
lines, the reported value is
\[
 \frac{2M}{N_1+N_2}.
\]
The inputs are
\path{public_records/maxwell_siegel_arxiv_2026-08-16.json} and
\path{dossiers/_structural/LATEX_PROFILE_INDEX.jsonl}; exact compared paths,
source SHA-256 hashes, line counts, edit counts, and full-precision values are
in \path{dossiers/_structural/VERSION_DELTAS.jsonl}, whose SHA-256 is
\nolinkurl{40520bb04b54100bbc3a640f752c9fba0fbba4378545b3ad8b01dbb459d22647}.
This derived diagnostic only orders collation work.  It neither changes a
source nor proves semantic equivalence, semantic difference, or permission to
merge witnesses.

\begingroup\small
\setlength{\tabcolsep}{3pt}
\begin{longtable}{@{}p{0.12\textwidth}p{0.31\textwidth}p{0.25\textwidth}p{0.25\textwidth}@{}}
\toprule
Lineage & Work & Witness history & Canonical treatment\\
\midrule
\endfirsthead
\toprule
Lineage & Work & Witness history & Canonical treatment\\
\midrule
\endhead
1909.09733 & \emph{Conservation of Singularities in Functional Equations Associated to Collatz-Type Dynamical Systems; or, Dreamcatchers for Hydra Maps}
& v1, 20 Sep. 2019; v2, 4 Oct. 2019; v3, 19 Oct. 2019.  Similarities: v1--v2, $0.6041$; v2--v3, $0.8736$.
& Use v3; retain v1 and v2 as substantial same-title witnesses.  This is the earlier unit-disk/profinite dreamcatcher branch, not a chapter-version of the dissertation's central $(p,q)$-adic branch.\\

2002.12762 & \emph{Fourier Analysis of the Parity-Vector Parameterization of the Generalized Collatz $px+1$ maps}
& v1, 24 Feb. 2020, substantive paper; v2, 14 Dec. 2023, metadata-only withdrawal.
& Withdrawn as obsolete and under the wrong subject classification.  Preserve v1 only as a parity/Fourier precursor and historical witness.\\

2007.15936 & \emph{The Collatz Conjecture \& Non-Archimedean Spectral Theory: Part I---Arithmetic Dynamical Systems and Non-Archimedean Value Distribution Theory}
& v1, 19 Jul. 2020, titled \emph{Syracuse Random Variables and the Periodic Points of Collatz-type maps}; v2, 28 Mar. 2023; v3, 24 Apr. 2023.  Similarities: v1--v2, $0.0248$; v2--v3, $0.9790$.
& Use v3 for the shortened-$qx+1$ numen and Correspondence Principle.  Treat v1 as a materially distinct predecessor, not a minor draft.  Collate the journal-layout PDF and Web installments I--II.\\

2111.07882 & \emph{Functional Equations Associated to Collatz-Type Maps on Integer Rings of Algebraic Number Fields}
& v1, 10 Nov. 2021, substantive paper; v2, 16 Apr. 2022, metadata-only withdrawal.
& Withdrawn for an error in Section 3 concerning a quotient of a ring of algebraic integers modulo an ideal.  Historical only; do not use it to certify the number-field extension.\\

2111.07883 & \emph{The Collatz Conjecture \& Non-Archimedean Spectral Theory: Part I.5--How To Write The (Weak) Collatz Conjecture As A Contour Integral}
& v1, 10 Nov. 2021, titled \emph{A $p$-Adic Characterization of the Periodic Points of a Class of Collatz-Type Maps on the Integers}; v2, 15 Dec. 2023.  Similarity $0.0551$.
& Use v2 for the Dirichlet/Mellin--Perron branch.  Preserve v1 separately as a substantially different foundational witness to the Correspondence Principle.\\

2208.11082 & \emph{The Collatz Conjecture \& Non-Archimedean Spectral Theory: Part II---$(p,q)$-Adic Fourier Analysis and Wiener's Tauberian Theorem}
& v1, 18 Aug. 2022, titled \emph{Wiener Tauberian Theorems and the Value Distributions of $(p,q)$-adic Functions}; v2, 20 Jun. 2025.  Similarity $0.2870$.
& Use v2; preserve v1 as a materially distinct shorter witness.  Collate Web installment III.\\

2307.00436 & \emph{Infinite Series Whose Topology of Convergence Varies From Point to Point}
& v1, 1 Jul. 2023.
& Use v1 for frames, F-series, rising-continuity, and variable topology.  It expressly takes precedence over conflicting dissertation terminology.\\

2412.02902 & \emph{$(p,q)$-adic Analysis and the Collatz Conjecture}
& v1, 3 Dec. 2024; arXiv witness of the 2022 University of Southern California dissertation.
& Use the TeX as the machine-readable backbone, collated with the August and corrected October 2022 site PDFs.  Later specific sources override it only within their declared scope.\\

2507.13358 & \emph{Algebras of $p$-Adic Distributions Induced by Pointwise Products of F-Series}
& v1, 1 Jul. 2025.
& Use v1.  The March 2025 Kansas talk is a precursor and expository witness, not a replacement source.\\

2601.17030 & \emph{The Hydra Map and Numen Formalisms for Collatz-Type Problems}
& v1, 19 Jan. 2026; v2, 29 Jan. 2026; v3, 12 Feb. 2026.  Similarities: v1--v2, $0.9691$; v2--v3, $0.9935$.
& Use v3 as the present definition authority and global-field manual.  The v3 metadata states that minor errors in Proposition 7.6 were rectified.\\
\bottomrule
\end{longtable}
\endgroup

The changed titles and direct collation of 2007.15936v1, 2111.07883v1, and
2208.11082v1 require separate witness records.  The routing values above only
prioritize that collation: neither a shared arXiv identifier nor a line metric
is a semantic merge instruction.

\subsection{The Web installments and the paper-series numbering}

The author site's four-part numbering is not the paper-series numbering.  To
prevent false citations, this edition always prefixes the site objects with
\emph{Web installment}.

\begingroup\small
\setlength{\tabcolsep}{3pt}
\begin{longtable}{@{}p{0.12\textwidth}p{0.30\textwidth}p{0.20\textwidth}p{0.31\textwidth}@{}}
\toprule
Site label & Exact title & Page dates & Relation and status\\
\midrule
\endfirsthead
\toprule
Site label & Exact title & Page dates & Relation and status\\
\midrule
\endhead
Web I & \emph{The Collatz Conjecture \& Non-Archimedean Spectral Theory--Part I--The Numen of a Hydra Map}
& Published 2022-08-09 21:45:22 UTC; page reports modification 2024-07-11 21:24:00 UTC.
& Expository witness for the construction half of paper Part I, arXiv:2007.15936v3.  It is not an additional arXiv Part I.\\

Web II & \emph{The Collatz Conjecture \& Non-Archimedean Spectral Theory--Part II--The Correspondence Principle}
& Published 2022-08-10 20:55:28 UTC; page reports modification 2022-10-11 19:55:07 UTC.
& Expository witness for the Correspondence-Principle half of \emph{paper Part I}.  Its final Dirichlet/Perron aside is a precursor to paper Part I.5, arXiv:2111.07883v2.  It is not paper Part II.\\

Web III & \emph{The Collatz Conjecture \& Non-Archimedean Spectral Theory--Part III--$(p,q)$-adic Fourier Analysis \& Tauberian Spectral Theory}
& Published 2022-08-12 22:51:30 UTC; archived page reports modification 2025-12-26 22:12:20 UTC.
& Evolving site exposition overlapping \emph{paper Part II}, arXiv:2208.11082v2.  Unique material may be used after theorem-level collation; its modification date alone confers no authority.\\

Web IV & \emph{The Collatz Conjecture \& Non-Archimedean Spectral Theory--Part IV--A Fourier Analysis of the Numen of the Shortened $qx+1$ map}
& Published 2022-08-15 19:14:38 UTC; archived page reports modification 2025-12-26 22:13:00 UTC.
& Site-only numen/Fourier application, chiefly parallel to the dissertation's one-dimensional Fourier material.  There is no corresponding arXiv paper Part IV in this corpus.\\
\bottomrule
\end{longtable}
\endgroup

The archived Web I and Web II objects are current snapshots acquired on
23 August 2026.  WordPress's public endpoint did not expose their revision
histories without authentication.  Consequently their published and modified
timestamps do not identify the contents of every intervening revision.

\subsection{Dissertation and supplemental author-site witnesses}

The dissertation has three distinct public witnesses:

\begin{enumerate}[label=\textup{(\alph*)}]
\item the August 2022 author-site PDF, 476 pages;
\item the October 2022 author-site PDF explicitly labelled corrected,
      478 pages; and
\item arXiv:2412.02902v1, submitted 3 December 2024, 467 pages, with an
      extractable TeX source.
\end{enumerate}

Their normalized texts and pagination differ.  The arXiv TeX is therefore an
editable base, not evidence that the explicitly corrected October PDF is a
byte-for-byte or proposition-for-proposition predecessor.  Any difference
used in the edition must be collated at the relevant passage.

The remaining author-site witnesses have the following roles.

\begin{description}[style=nextline]
\item[Part I journal-layout PDF.]
Same-work layout witness for 2007.15936; use for publication pagination and
typographic variants only after comparison with v3.

\item[Infinite-series journal-layout PDF.]
Same-work layout witness for 2307.00436; it does not silently replace the
available TeX.

\item[\emph{Arithmetic Properties of F-series} talk notes.]
March 2025 Emporia State University talk; supplementary precursor to
2507.13358.

\item[Episode 1 script.]
June 2024 companion exposition for the dissertation and the $(p,q)$-adic
programme; useful orientation, not a primary theorem source.

\item[$p$-adic $L$-functions seminar and Fourier-series notes.]
Siegel-authored background with high contextual relevance.  These are
supplementary notes, not parts of the Collatz research lineage.
\end{description}

\subsection{Corpus boundary}

The core edition includes the ten arXiv lineages, all materially distinct
versions identified above, the four Web installments, the three dissertation
witnesses, the two same-work layout witnesses, and the research-specific talk
and companion script.  The $p$-adic $L$-function seminar and Fourier-series
notes are Siegel-authored apparatus, not core research works.  Authorship and
core status are separate questions.  Siegel's generic
course notes in introductory real analysis, linear algebra, vector calculus,
complex variables, partial and ordinary differential equations, and
probability remain Siegel-authored works, but they are outside the Collatz,
Hydra, numen, dreamcatcher, and F-series research corpus.

External supporting literature is not part of Siegel's corpus, but it is used
with exact citations to establish prior art, standard theorems, and comparison
maps.

\subsection{External prior-art comparison ledger}

The following sources are outside Siegel's corpus.  They are present because
an exact critical edition must distinguish Siegel's added structure from an
older identical special case, and must distinguish a shared value formula
from an identity of maps.  Each row records the strongest comparison proved
in this edition and the boundary beyond which the citation may not be used.

\begingroup\small
\setlength{\tabcolsep}{3pt}
\begin{longtable}{@{}p{0.19\textwidth}p{0.23\textwidth}p{0.28\textwidth}p{0.23\textwidth}@{}}
\toprule
Object & Primary source & Exact proved overlap & Nonidentity boundary\\
\midrule
\endfirsthead
\toprule
Object & Primary source & Exact proved overlap & Nonidentity boundary\\
\midrule
\endhead
Scalar integer Hydra at displayed modulus $D\in\Z_{\geq2}$ &
Matthews, Sec.~2 \cite{MatthewsGeneralized} & Admissibility is equivalent to
unique integers $m_i\ne0$ and $\rho_i\equiv im_i\pmod D$ with
$H_i(x)=(m_ix-\rho_i)/D$; see
\cref{prop:scalar-hydra-matthews}.  Exact-integrality is equivalent to
the displayed-presentation condition
$\gcd(m_0\cdots m_{D-1},D)=1$. & Matthews attaches relatively prime type to
the least-modulus presentation.  Thus exact-integrality is equivalent to
Matthews relatively prime type only together with
$D=d_{\mathrm{Mat}}(H)$.  A nonminimal refinement may destroy
exact-integrality.  Neither comparison identifies Siegel's non-scalar or
global-field structure or adds centeredness or properness.\\

Residue-class-wise affine mapping & Kohl, Defs.~1.1 and~1.4 and Rem.~1.2
\cite{KohlRCWA} & Every scalar integer Hydra in the preceding row is rcwa,
and $d_{\mathrm{Mat}}(H)=\operatorname{Mod}(H)\mid D$; under
exact-integrality both least moduli equal $D$; see
\cref{prop:scalar-hydra-rcwa}. & The group
$\operatorname{RCWA}(\Z)$ contains only bijective rcwa maps.  The shortened
Collatz map is rcwa but is not injective; a Hydra is therefore not
automatically an rcwa group element.\\

Formal affine word and $q=3$ return candidate & B\"ohm--Sontacchi,
Prop.~1; printed Props.~4--5; unnumbered paragraph after Prop.~5 on p.~264
\cite{BohmSontacchi} & Proposition~1 specializes to the valid formal
composition in \cref{eq:chronological-affine-return}; at $q=3$, its fixed
point is the rational $R_{3,\boldsymbol\varepsilon}$ printed as the
Proposition~5 candidate after the explicit renaming $n=L$. & Printed
Proposition~4's unrestricted identification with $u^L(x)$ is false; it may be
made only for a matching itinerary.  The unnumbered p.~264 paragraph allows
exact period properly dividing $L$.  The final strict position inequality is
reconstructed from the formal run data, not inserted into printed
Proposition~5.\\

$2$-adic parity conjugacy & Bernstein--Lagarias, Eqs.~(1.3)--(1.6),
(4.1)--(4.2), and App.~B \cite{BernsteinLagarias} & The chronological parity
encoder $Q_q$, inverse conjugacy $\Phi_q$, and every odd-$q$ classical
specialization are exact.  For odd prime $q$, the $\Phi_q$ value is
$\iota_2(R_{q,\boldsymbol\varepsilon})$ and the reversed $\chi_q$ value is
$\iota_q(R_{q,\boldsymbol\varepsilon})$ as in
\cref{thm:BL-numen-periodic-overlap}. & $Q_q$ and
$\Phi_q$ retain their inverse roles in $\Z_2$; for odd prime $q$, neither is
the cross-adic map $\chi_q:\Z_2\to\Z_q$.  Composite odd $q$ in the classical conjugacy does not
enlarge the edition's prime-$q$ numen theorem.\\

De Bruijn graph presentation & Laarhoven--de Weger, Thms.~3.1 and~5.1 and
Sec.~5.2 \cite{LaarhovenDeWeger} & Their graph isomorphism is the
chronological parity map denoted $Q_3$ here; their odd-$an+b$ analogue is an
exact classical bridge. & Their paper's symbol $\Phi$ has the orientation of
this edition's $Q_3$, not Bernstein--Lagarias's $\Phi_3$; its $p$-ary
discussion uses appropriately chosen coefficients and is not an arbitrary-
Hydra theorem.\\
\bottomrule
\end{longtable}
\endgroup

\subsection{Source-issue ledger}

\begingroup\small
\begin{longtable}{@{}p{0.10\textwidth}p{0.22\textwidth}p{0.61\textwidth}@{}}
\toprule
Code & Source & Canonicalization issue\\
\midrule
\endfirsthead
\toprule
Code & Source & Canonicalization issue\\
\midrule
\endhead
\textsc{SR-01} & Corpus index & A shared logical-work identifier records bibliographic lineage, not semantic sameness.  The reproducible routing values $0.0248$, $0.0551$, and $0.2870$ do not establish that distinction and do not authorize a merge; the separately archived versioned sources and direct collation control witness status.\\

\textsc{SR-02} & 1909.09733 & A common title is not an identity proof.  Versions v1 and v2 remain separately archived in the apparatus even though v3 is the base; the value $0.6041$ is only a reproducible collation-routing diagnostic.\\

\textsc{SR-03} & 2002.12762 & v2 is a metadata-only withdrawal declaring the paper obsolete and under the wrong subject classification.  It supplies no corrected paper; v1 is historical only.\\

\textsc{SR-04} & 2111.07882 & v2 withdraws the work for an error in Section 3 about a quotient ring.  The archived comment contains the contradictory phrase ``was not incorrectly described''; the withdrawal and existence of an error are clear, but the intended corrected assertion must not be guessed.\\

\textsc{SR-05} & 2601.17030 & The paper expressly replaces potentially inconsistent earlier Hydra definitions.  v3 also reports corrections to minor errors in Proposition 7.6.  Earlier Hydra terminology is historical unless re-derived in the v3 formalism.\\

\textsc{SR-06} & 2307.00436 & The paper expressly takes precedence over the dissertation where their frames/F-series terminology conflicts.  It is more general in some respects and less comprehensive in others; only the conflicting definitions are superseded.\\

\textsc{SR-07} & Dissertation & The August 2022, corrected October 2022, and December 2024 arXiv witnesses are textually distinct.  ``Later upload'' and ``corrected'' are recorded attributes, not a license to discard the other witnesses.\\

\textsc{SR-08} & Web numbering & Web II is part of the web sequence but belongs chiefly to arXiv \emph{paper Part I}; Web III overlaps arXiv \emph{paper Part II}.  Bare labels such as ``Part II'' are therefore prohibited in source-critical references.\\

\textsc{SR-09} & Web I & The current snapshot calls \emph{numen} Greek whereas paper Part I correctly identifies it as Latin; its displayed $q$-adic decay formula omits the subscript $q$ on the norm; and its uniqueness proof says equality on $\Z_q$ where the input domain must be $\Z_2$.  Paper Part I v3 controls these points.  Hash-bound locator: \path{author_site/wordpress_2026-08-23/part_i_numen_hydra_map_2024-07-11.txt}, SHA-256 \nolinkurl{6dfcf3ee7c4056978deae669094e08f41acbfbd7f1410b2cafb3d58293f682ae}, lines 143, 271--275, and 327--329.\\

\textsc{SR-10} & Web II & Claim 5's statement omits the conclusion $\chi_q(z/2)=0$, although its proof supplies it.  The page calls its terminal periodic-point formulation Version 3, while paper Part I v3 canonically labels its corresponding published formulation Version 4.  Its divergent result says ``belongs to a divergent trajectory''; paper Part I v3 states directly that the value is a divergent point.  Hash-bound locator: \path{author_site/wordpress_2026-08-23/part_ii_correspondence_principle_2022-10-11.txt}, SHA-256 \nolinkurl{b5d37b71f4778109fbd853f621b0562c6aeca1487141968d2cde3887c8ba46c6}, lines 283--287, 353--355, and 359--373.\\

\textsc{SR-11} & Web I--IV & Publication and modification timestamps are provenance only.  No page supersedes a formal source solely because its modification date is later, and the public snapshots do not furnish a complete revision history.\\

\textsc{SR-12} & Web III--IV & These are active, evolving expositions.  Statements unique to them require theorem-level checking before incorporation; overlap with paper Part II or the dissertation is not presumed identity.\\

\textsc{SR-13} & Layout PDFs & The Part I and infinite-series layouts are distinct files and must be collated for wording and theorem numbering.  They are not replacement TeX sources merely because they resemble journal typography.\\

\textsc{SR-14} & 2111.07883v2; Web II & In the calculation of $f_q(z^2,x)$ immediately before the printed ordinary generating-function equation, the first term is written as $z/(1-z)$ instead of $z^2/(1-z^2)$.  The resulting Dirichlet functional equation is false for $x\ne0$, and its $q=3$ residue formulas contain spurious $x$-terms.  Part I.5 and the public Web II nevertheless identify the continuation programme.  Chapter 7 recovers it by deriving the corrected equation, proving normal convergence and global strip-by-strip continuation, and recalculating the principal parts.  Denominator zeros remain candidate poles unless noncancellation is proved.\\

\textsc{SR-15} & B\"ohm--Sontacchi (1978) & Printed Proposition~4 is false
under its stated unrestricted quantifiers.  With $x=0$, $m=1$, and
$n_0=n_1=1$, it gives the right side $1/4$, while $u^3(0)=0$; the right side
is instead the formal composition $f\circ g\circ f(0)=1/4$.  The edition
therefore derives the formal affine identity from Proposition~1, invokes
$u^L$ only after proving actual-itinerary correctness, and recovers the
Proposition~5 fixed-return candidate independently.  The possibility of
exact period properly dividing $L$ is located in the unnumbered paragraph
after Proposition~5 on p.~264, not attributed generically to
Propositions~4--6.\\

\bottomrule
\end{longtable}
\endgroup

\begingroup
\small
\sloppy
\def\UrlBreaks{\do\/\do-\do\_}
\Urlmuskip=0mu plus 1mu\relax
\begin{thebibliography}{99}

\bibitem{SiegelDreamcatchers}
Maxwell C. Siegel,
\emph{Conservation of Singularities in Functional Equations Associated to
Collatz-Type Dynamical Systems; or, Dreamcatchers for Hydra Maps},
arXiv:1909.09733.  Version history: v1, 20 September 2019; v2,
4 October 2019; v3, 19 October 2019.  Canonical witness: v3.
\url{https://arxiv.org/abs/1909.09733v3}.

\bibitem{SiegelParityFourier}
Maxwell C. Siegel,
\emph{Fourier Analysis of the Parity-Vector Parameterization of the
Generalized Collatz $px+1$ maps}, arXiv:2002.12762v1, 24 February 2020;
withdrawal revision v2, 14 December 2023.
Substantive v1: \url{https://arxiv.org/abs/2002.12762v1}.
Withdrawal record v2: \url{https://arxiv.org/abs/2002.12762v2}.

\bibitem{SiegelSyracuseV1}
Maxwell C. Siegel,
\emph{Syracuse Random Variables and the Periodic Points of Collatz-type
maps}, arXiv:2007.15936v1, 19 July 2020.  Materially distinct predecessor
under the later Part I lineage.
\url{https://arxiv.org/abs/2007.15936v1}.

\bibitem{SiegelPartI}
Maxwell C. Siegel,
\emph{The Collatz Conjecture \& Non-Archimedean Spectral Theory: Part I---
Arithmetic Dynamical Systems and Non-Archimedean Value Distribution Theory},
arXiv:2007.15936v3, 24 April 2023; v2 submitted 28 March 2023.
Publication DOI: \url{https://doi.org/10.1134/S2070046624020055}.
\url{https://arxiv.org/abs/2007.15936v3}.

\bibitem{SiegelNumberRings}
M. C. Siegel,
\emph{Functional Equations Associated to Collatz-Type Maps on Integer Rings
of Algebraic Number Fields}, arXiv:2111.07882v1, 10 November 2021;
withdrawal revision v2, 16 April 2022.
Substantive v1: \url{https://arxiv.org/abs/2111.07882v1}.
Withdrawal record v2: \url{https://arxiv.org/abs/2111.07882v2}.

\bibitem{SiegelPadicCharacterizationV1}
Maxwell C. Siegel,
\emph{A $p$-Adic Characterization of the Periodic Points of a Class of
Collatz-Type Maps on the Integers}, arXiv:2111.07883v1,
10 November 2021.  Materially distinct predecessor under the later Part I.5
lineage.  \url{https://arxiv.org/abs/2111.07883v1}.

\bibitem{SiegelPartIHalf}
Maxwell C. Siegel,
\emph{The Collatz Conjecture \& Non-Archimedean Spectral Theory: Part I.5--
How To Write The (Weak) Collatz Conjecture As A Contour Integral},
arXiv:2111.07883v2, 15 December 2023.
\url{https://arxiv.org/abs/2111.07883v2}.

\bibitem{SiegelWienerV1}
Maxwell C. Siegel,
\emph{Wiener Tauberian Theorems and the Value Distributions of $(p,q)$-adic
Functions}, arXiv:2208.11082v1, 18 August 2022.  Materially distinct
predecessor under the later Part II lineage.
\url{https://arxiv.org/abs/2208.11082v1}.

\bibitem{SiegelPartII}
Maxwell C. Siegel,
\emph{The Collatz Conjecture \& Non-Archimedean Spectral Theory: Part II---
$(p,q)$-Adic Fourier Analysis and Wiener's Tauberian Theorem},
\emph{$p$-Adic Numbers, Ultrametric Analysis and Applications}
\textbf{17}(2) (2025), 187--232,
\url{https://doi.org/10.1134/S2070046625020062};
arXiv:2208.11082v2, 20 June 2025.
\url{https://arxiv.org/abs/2208.11082v2}.

\bibitem{SiegelInfiniteSeries}
Maxwell C. Siegel,
\emph{Infinite Series Whose Topology of Convergence Varies From Point to
Point}, \emph{$p$-Adic Numbers, Ultrametric Analysis and Applications}
\textbf{15}(2) (2023), 133--167,
\url{https://doi.org/10.1134/S2070046623020061};
arXiv:2307.00436v1, 1 July 2023.
\url{https://arxiv.org/abs/2307.00436v1}.

\bibitem{SiegelDissertationArxiv}
Maxwell C. Siegel,
\emph{$(p,q)$-adic Analysis and the Collatz Conjecture}, Ph.D. dissertation,
University of Southern California, 2022; arXiv:2412.02902v1,
3 December 2024.  \url{https://arxiv.org/abs/2412.02902v1}.

\bibitem{SiegelFSeriesProducts}
Maxwell C. Siegel,
\emph{Algebras of $p$-Adic Distributions Induced by Pointwise Products of
F-Series}, arXiv:2507.13358v1, 1 July 2025.
\url{https://arxiv.org/abs/2507.13358v1}.

\bibitem{SiegelHydraManual}
Maxwell C. Siegel,
\emph{The Hydra Map and Numen Formalisms for Collatz-Type Problems},
arXiv:2601.17030.  Version history: v1, 19 January 2026; v2,
29 January 2026; v3, 12 February 2026.  Canonical witness: v3.
\url{https://arxiv.org/abs/2601.17030v3}.

\bibitem{SiegelWebI}
Maxwell C. Siegel,
\emph{The Collatz Conjecture \& Non-Archimedean Spectral Theory--Part I--
The Numen of a Hydra Map}, author-site post (Web installment I in this
edition), published 2022-08-09 21:45:22 UTC; page reports modification
2024-07-11 21:24:00 UTC; snapshot acquired 23 August 2026.
\url{https://siegelmaxwellc.wordpress.com/2022/08/09/the-collatz-conjecture-non-archimedean-spectral-theory-part-i-the-numen-of-a-hydra-map/}.

\bibitem{SiegelWebII}
Maxwell C. Siegel,
\emph{The Collatz Conjecture \& Non-Archimedean Spectral Theory--Part II--
The Correspondence Principle}, author-site post (Web installment II in this
edition), published 2022-08-10 20:55:28 UTC; page reports modification
2022-10-11 19:55:07 UTC; snapshot acquired 23 August 2026.
\url{https://siegelmaxwellc.wordpress.com/2022/08/10/the-collatz-conjecture-non-archimedean-spectral-theory-part-ii-the-correspondence-principle/}.

\bibitem{FedjaMO333801}
fedja (MathOverflow user 1131),
answer 333801 to \emph{Overconcentration of Poles on the Circle of
Convergence of a Power Series with Bounded Coefficients}, accepted
MathOverflow answer, created 12 June 2019 at 03:24:50 UTC and last edited
18 June 2019 at 01:50:51 UTC, CC BY-SA 4.0.
\url{https://mathoverflow.net/a/333801}.

\bibitem{MatthewsGeneralized}
K. R. Matthews,
\emph{Generalized $3x + 1$ Mappings: Markov Chains and Ergodic Theory},
in Jeffrey C. Lagarias (ed.), \emph{The Ultimate Challenge: The $3x+1$
Problem}, American Mathematical Society, Providence, RI, 2010, pp.~79--103.
The consulted author-form witness states that it contains small additions and
deletions relative to the book chapter.

\bibitem{KohlRCWA}
Stefan Kohl,
\emph{Algorithms for a class of infinite permutation groups},
\emph{Journal of Symbolic Computation} \textbf{43}(8) (2008), 545--581.
\url{https://doi.org/10.1016/j.jsc.2007.12.001}.

\bibitem{BohmSontacchi}
Corrado B\"ohm and Giovanna Sontacchi,
\emph{On the existence of cycles of given length in integer sequences like
$x_{n+1}=x_n/2$ if $x_n$ even, and $x_{n+1}=3x_n+1$ otherwise},
\emph{Atti della Accademia Nazionale dei Lincei. Classe di Scienze Fisiche,
Matematiche e Naturali. Rendiconti}, Series~8, \textbf{64}(3) (1978),
260--264.
\url{https://www.bdim.eu/item?id=RLINA_1978_8_64_3_260_0}.

\bibitem{BernsteinLagarias}
Daniel J. Bernstein and Jeffrey C. Lagarias,
\emph{The $3x + 1$ Conjugacy Map},
\emph{Canadian Journal of Mathematics} \textbf{48}(6) (1996), 1154--1169.
\url{https://doi.org/10.4153/CJM-1996-060-x}.

\bibitem{LaarhovenDeWeger}
Thijs Laarhoven and Benne de Weger,
\emph{The Collatz conjecture and De Bruijn graphs},
arXiv:1209.3495v1, 16 September 2012.
\url{https://arxiv.org/abs/1209.3495v1}.
The consulted local PDF also displays 27 November 2024; that display is not
used as the arXiv submission date.

\bibitem{SiegelWebIII}
Maxwell C. Siegel,
\emph{The Collatz Conjecture \& Non-Archimedean Spectral Theory--Part III--
$(p,q)$-adic Fourier Analysis \& Tauberian Spectral Theory}, author-site post
(Web installment III in this edition), published 2022-08-12 22:51:30 UTC;
archived page reports modification 2025-12-26 22:12:20 UTC.
\url{https://siegelmaxwellc.wordpress.com/2022/08/12/the-collatz-conjecture-non-archimedean-spectral-theory-part-iii-a-pq-adic-fourier-analysis/}.

\bibitem{SiegelWebIV}
Maxwell C. Siegel,
\emph{The Collatz Conjecture \& Non-Archimedean Spectral Theory--Part IV--
A Fourier Analysis of the Numen of the Shortened $qx+1$ map}, author-site
post (Web installment IV in this edition), published 2022-08-15 19:14:38
UTC; archived page reports modification 2025-12-26 22:13:00 UTC.
\url{https://siegelmaxwellc.wordpress.com/2022/08/15/the-collatz-conjecture-non-archimedean-spectral-theory-part-iv-a-fourier-analysis-of-the-numen-of-the-shortened-qx-1-map/}.

\bibitem{SiegelDissertationAugust}
Maxwell C. Siegel,
\emph{$(p,q)$-adic Analysis and the Collatz Conjecture}, Ph.D. dissertation,
University of Southern California, August 2022 author-site edition.
\url{https://siegelmaxwellc.wordpress.com/wp-content/uploads/2022/08/phd_dissertation_final.pdf}.

\bibitem{SiegelDissertationOctober}
Maxwell C. Siegel,
\emph{$(p,q)$-adic Analysis and the Collatz Conjecture}, Ph.D. dissertation,
University of Southern California, corrected October 2022 author-site
edition.  \url{https://siegelmaxwellc.wordpress.com/wp-content/uploads/2022/10/phd_dissertation_final.pdf}.

\bibitem{SiegelPartILayout}
Maxwell C. Siegel,
\emph{The Collatz Conjecture \& Non-Archimedean Spectral Theory: Part I},
author-site journal-layout PDF, same-work witness for arXiv:2007.15936.
\url{https://siegelmaxwellc.wordpress.com/wp-content/uploads/2024/05/blog_post_i_numen_and_correspondence_principle.pdf}.

\bibitem{SiegelInfiniteSeriesLayout}
Maxwell C. Siegel,
\emph{Infinite Series Whose Topology of Convergence Varies From Point to
Point}, author-site journal-layout PDF, same-work witness for
arXiv:2307.00436.
\url{https://siegelmaxwellc.wordpress.com/wp-content/uploads/2023/12/infinite_series_whose_topology_of_convergence_varies_from_point_to_point_simplified.pdf}.

\bibitem{SiegelKansasTalk}
Maxwell C. Siegel,
\emph{Arithmetic Properties of F-series; or, How to $3$-adically Integrate
a $5$-adic Function}, Emporia State University talk notes, March 2025;
author-site PDF uploaded May 2025.
\url{https://siegelmaxwellc.wordpress.com/wp-content/uploads/2025/05/2025_kansas_talk-3.pdf}.

\bibitem{SiegelEpisodeOne}
Maxwell C. Siegel,
\emph{Episode 1---$(p,q)$-adic Analysis and the Collatz Conjecture---A Whole
New World: Script}, author-site companion script, June 2024.
\url{https://siegelmaxwellc.wordpress.com/wp-content/uploads/2024/06/episode-1-script.pdf}.

\bibitem{SiegelPadicLFunctions}
Maxwell C. Siegel,
\emph{$p$-adic $L$-functions}, graduate seminar notes, author-site PDF.
\url{https://siegelmaxwellc.wordpress.com/wp-content/uploads/2023/02/algebraic_number_theory_-_presentation_-_p-adic_l-functions.pdf}.

\bibitem{SiegelFourierNotes}
Maxwell C. Siegel,
\emph{Fourier Series}, course notes, author-site PDF, 2025 site copy.
\url{https://siegelmaxwellc.wordpress.com/wp-content/uploads/2025/02/math_445_souto_-_class_notes_-_chapter_11_-_fourier_analysis.pdf}.

\end{thebibliography}
\endgroup


\end{document}
